% !TEX root = ../thesis.tex
% ==============================================================================
%  CHAPTER 3: LITERATURE REVIEW (فصل سوم: پیشینه پژوهش) (chapters/chapter3_literature.tex)
% ==============================================================================
%  This chapter provides a comprehensive review, taxonomy, and critical
%  comparative survey of existing paradigms and analyzes their theoretical limitations.
% ==============================================================================

\chapter{پیشینه‌ی پژوهش}\label{chap:chapter3_literature}

این فصل با هدف تبیین جایگاه پژوهش حاضر در پیشینه‌ی تفسیرپذیری هوش مصنوعی، به مرور انتقادی و ارائه‌ی دسته‌بندی جامعی از پارادایم‌های تفسیرپذیری می‌پردازد. در این راستا، با واکاوی دقیق رویکردهای پیشین و آسیب‌شناسی محدودیت‌های آن‌ها، شکاف تحقیقاتی و بن‌بست‌های نظری موجود در ارزیابی سیستم‌های جعبه‌سیاه برجسته می‌گردد. در ادامه و در نیمه‌ی دوم فصل، پیشینه‌ی ابزارهای نوین پردازش زبان طبیعی – با تمرکز بر ارتقای واحد معنایی به سطح جمله و بهره‌گیری از مدل‌های مبتنی بر انرژی – که شالوده‌ی محاسباتی و تئوریک روش پیشنهادی را شکل می‌دهند، به‌تفصیل بررسی خواهد شد.



% =========================================================
% بدنه اول
% =========================================================



\section{پارادایم‌های تفسیرپذیری در یادگیری ماشین}\label{sec:3-1}

در این بخش، دسته‌بندی روش‌های تفسیرپذیری به‌تفصیل بررسی می‌شود تا جایگاه و موانع هر رویکرد در چشم‌انداز کنونی هوش مصنوعی مشخص گردد. تفسیرپذیری در یادگیری ماشین صرفاً به معنای ایجاد یک توضیح ساده و شهودی از رفتار مدل نیست؛ بلکه همان‌طور که دوشی-ولز و کیم (۲۰۱۷)~\cite{doshivelez2017rigorous} در یکی از معتبرترین چارچوب‌های نظری این حوزه تبیین کرده‌اند، تفسیرپذیری ابزاری ضروری برای تأیید ویژگی‌های پنهان و حیاتی مدل نظیر انصاف، امنیت و قابلیت‌اطمینان است. در این چارچوب، ارزیابی روش‌های تفسیری در یک دسته‌بندی سه‌گانه تعریف می‌شود: ارزیابی کاربردمحور\enfootnote{application-grounded evaluation} (توسط متخصصان دامنه در شرایط واقعی کاربری)، ارزیابی انسان‌محور\enfootnote{human-grounded evaluation} (توسط افراد عادی در وظایف آزمایشگاهی تقلیل‌یافته) و ارزیابی کارکردمحور\enfootnote{functionally-grounded evaluation} (بدون نیاز به آزمایش‌های انسانی و بر پایه‌ی شاخص‌های صوری و فرمول‌بندی‌شده‌ی محاسباتی نظیر عمق درخت یا میزان خلوتی). این ساختار تئوریک به‌خوبی روشن می‌سازد که چرا با افزایش پیچیدگی سیستم‌ها و هزینه‌ی بالای آزمایش‌های انسانی، بسیاری از پژوهش‌های تفسیرپذیری ناگزیر به‌سوی استفاده از پراکسی‌های محاسباتی و معیارهای صوری سوق یافته‌اند.

از سوی دیگر، با گذار از مدل‌های هوش مصنوعی متمایزگر به مدل‌های مولد و شکل‌گیری شاخه‌ی تفسیرپذیری هوش مصنوعی مولد، این حوزه با چالش‌های بنیادین جدیدی روبه‌رو شده است. اشنایدر (۲۰۲۴)~\cite{schneider2024explainable} در بررسی جامع خود نشان می‌دهد که ابزارهای کلاسیک تفسیرپذیری، که ماهیتاً برای خروجی‌های ساختاریافته‌ی مدل‌های متمایزگر طراحی شده بودند، در رمزگشایی از خروجی‌های درهم‌تنیده، تعاملی و ابعاد بالای مدل‌های زبانی بزرگ (\lr{LLM}) ناکارآمدند. این بررسی مروری با ترسیم ابعاد مختلف تفسیر (مانند گستره‌ی محلی در برابر سراسری)، اثبات می‌کند که برای تحلیل سیستم‌های جعبه‌سیاه مدرن، نیازمند عبور از پارادایم‌های قدیمی و توسعه‌ی متدولوژی‌های تطبیق‌یافته و معنایی هستیم. با در نظر داشتن این گذار تاریخی و تئوریک، در ادامه‌ی این بخش، پارادایم‌های کلان تفسیرپذیری (درونی، پساعملی جعبه‌سفید و پساعملی جعبه‌سیاه) مورد واکاوی و نقد قرار می‌گیرند~\cite{guidotti2018survey, danilevsky2020survey}. 


\subsection{تفسیرپذیری درونی}\label{sec:3-1-1}

تفسیرپذیری درونی به دسته‌ای از مدل‌های یادگیری ماشین اشاره دارد که معماری و ساختار آن‌ها از پایه به‌گونه‌ای طراحی شده است که فرایند تصمیم‌گیری و منطق محاسباتی‌شان ذاتاً برای انسان شفاف و قابل‌درک باشد~\cite{guidotti2018survey}. در این پارادایم، برای فهمیدن چرایی تولید یک خروجی، نیازی به استفاده از ابزارهای ثانویه، مهندسی معکوس، یا مدل‌های مفسر خارجی نیست؛ بلکه خود مدل، ساختار استنتاجی خود را به‌صورت مستقیم در اختیار کاربر قرار می‌دهد. تلاش‌ها در این حوزه از مدل‌های آماری کلاسیک آغاز شده و تا طراحی شبکه‌های عصبی مقید و بلوک‌های ساختاری شفاف امتداد یافته است. با این وجود، کاربرد این پارادایم در مواجهه با سیستم‌های هوش مصنوعی مدرن با چالش‌های نظری و عملی عمیقی روبه‌روست.


\subsubsection{مبانی و سیر تطور}\label{sec:3-1-1-1}

درک مفهوم تفسیرپذیری درونی در یادگیری ماشین، ریشه در مدل‌های کلاسیک نظیر رگرسیون خطی، درخت‌های تصمیم و سیستم‌های مبتنی بر قانون دارد~\cite{molnar2022interpretable}. این مدل‌ها به‌دلیل سادگی ریاضیاتی و محدود بودن فضای فرضیه، بالاترین سطح از شفافیت را ارائه می‌دهند. برای نمونه، در یک مدل رگرسیون خطی، وزن‌های (ضرایب) اختصاص‌یافته به هر ویژگی ورودی، به‌طور دقیق و کمی نشان می‌دهند که آن ویژگی چه میزان و در چه جهتی (مثبت یا منفی) بر خروجی نهایی تأثیر می‌گذارد. به‌طور مشابه، در یک درخت تصمیم، مسیر طی‌شده از ریشه تا برگ نهایی، یک زنجیره‌ی استدلالی از قوانین شرطی (اگر-آنگاه) را شکل می‌دهد که قدم‌به‌قدم برای یک ناظر انسانی قابل‌ردیابی و ارزیابی است.

با این حال، کارایی این مدل‌های شفاف تنها محدود به داده‌های ساختاریافته‌ی جدولی و مسائلی با ابعاد پایین است. با افزایش پیچیدگی داده‌ها – به‌ویژه در مواجهه با داده‌های غیرساختاریافته و با ابعاد بسیار بالا نظیر پردازش زبان طبیعی و متون خام – مدل‌های کلاسیک توانایی یادگیری روابط غیرخطی، سلسله‌مراتبی و به‌شدت درهم‌تنیده‌ی موجود در داده‌ها را از دست می‌دهند. برای دستیابی به دقت بالاتر در این وظایف پیچیده، محققان به سمت استفاده از شبکه‌های عصبی عمیق با میلیون‌ها و میلیاردها پارامتر حرکت کردند که این گذار، به شکل‌گیری یکی از چالش‌های بنیادین در حوزه‌ی یادگیری ماشین با عنوان مصالحه‌ی عملکرد-شفافیت منجر شد~\cite{doshivelez2017rigorous}.

در ادبیات این حوزه، رودین (۲۰۱۹)~\cite{rudin2019stop} استدلال می‌کند که در کاربردهای پرمخاطره (نظیر تشخیص‌های پزشکی یا قضایی)، باید به‌جای تلاش برای توضیح دادن مدل‌های جعبه‌سیاه غیرشفاف پس از آموزش، از همان ابتدا مدل‌هایی طراحی کرد که ذاتاً شفاف باشند. در شبکه‌های عصبی عمیقی که برای وظایف متمایزگر آموزش می‌بینند، تحقق این خواسته نیازمند اعمال سوگیری‌های استقرایی و محدودیت‌های ساختاری شدید بر معماری شبکه است؛ نظیر وادار کردن شبکه به رفتارهای خطی، اعمال خلوتی بر وزن‌ها، یا ایجاد گلوگاه‌های اطلاعاتی باریک.

اما با گذار از مدل‌های متمایزگر به مدل‌های مولد مدرن، اعمال این محدودیت‌های صریح هزینه‌ی بسیار سنگین‌تری در پی دارد. داده‌های تولیدی نظیر زبان طبیعی و تصاویر باکیفیت، دارای ساختاری بی‌نهایت منعطف و چندبُعدی هستند که مدل‌سازی آن‌ها نیازمند ظرفیت عظیم پارامتری برای ایجاد بازنمایی‌های توزیع‌شده است~\cite{elhage2022toymodels}. محدود کردن جریان اطلاعات یا تقلیل ساختار مدل به‌منظور دستیابی به شفافیت در فرایند تولید، نه تنها باعث افت فاحش عملکرد و کاهش روانی خروجی می‌شود~\cite{schneider2024explainable}، بلکه همان‌طور که فریدمن و همکاران (۲۰۲۴)~\cite{friedman2024interpretability} نشان داده‌اند، پدیده‌ای با عنوان توهم تفسیرپذیری ایجاد می‌کند؛ وضعیتی که در آن یک بازنمایی ساده‌شده تنها در داده‌های آشنا و درون‌توزیع وفادار به نظر می‌رسد، اما از بازتاب سازوکارهای واقعی مدل در استنتاج‌های پیچیده و تعمیم به شرایط جدید ناتوان است.

این مصالحه‌ی بنیادین میان ماهیت توزیع‌شده‌ی شبکه‌های مولد و نیاز به گره‌های محاسباتی تفسیرپذیر، مسیر پژوهش‌های تفسیرپذیری درونی را به دو شاخه‌ی متمایز سوق داد: نخست، تلاش برای دستیابی به بازنمایی‌های تفکیک‌شده با هدف مستقل‌سازی ابعاد در فضای پنهان مدل‌های پیوسته، و دوم، رویکردهای متأخر مبنی بر تعریف مفهوم که سعی دارند ساختار استنتاجی سیستم‌های پیچیده‌تر نظیر \lr{LLM}ها را مستقیماً با مفاهیم انسانی هم‌تراز کنند. در ادامه، سیر و چالش‌های نظری این دو جریان بررسی می‌شود.


\subsubsection{مدل‌های مبتنی بر تفکیک فضای پنهان}\label{sec:3-1-1-2}

یکی از نخستین رویکردها برای ایجاد شفافیت در معماری مدل‌های مولد، یادگیری بازنمایی‌های تفکیک‌شده در فضای پنهان است. ایده‌ی محوری در این پارادایم آن است که مدل به‌گونه‌ای آموزش ببیند که هر بُعد از فضای پنهان آن، به‌طور مستقل تنها یک ویژگی معنایی مشخص (مانند زاویه‌ی دید در یک تصویر یا سبک قلم) را کنترل کند. در این راستا، چن و همکاران (۲۰۱۶)~\cite{chen2016infogan} چارچوب \lr{InfoGAN} را به‌عنوان بسطی بر شبکه‌های مولد تخاصمی\enfootnote{generative adversarial networks} (\lr{GAN}) معرفی کردند. در یک \lr{GAN} معمولی، ورودی شبکه‌ی مولد ($G$) تنها یک بردار نویز تصادفی ناهمبسته ($z$) است که ابعاد آن به‌شکلی کاملاً درهم‌تنیده با هم ترکیب می‌شوند و هیچ مؤلفه‌ی آن به مفهوم ملموسی اختصاص ندارد. برای حل این مشکل، \lr{InfoGAN} فضای پنهان ورودی را به دو بخش افراز می‌کند: نویز تصادفی غیرساختاریافته‌ی پایه‌ای ($z$) و کدهای پنهان معنادار ($c$) که قرار است عوامل برجسته‌ی معنایی را بازنمایی کنند. سپس با افزودن یک عبارت منظم‌ساز مبتنی بر نظریه‌ی اطلاعات، بازی رقابتی حداقل-حداکثر استاندارد به‌شکل زیر بازفرمول‌بندی می‌شود:
\begin{equation}\label{eq:chap3_gan}
    \min_G \max_D V(D, G) - \lambda I(c; G(z, c))
\end{equation}
در این معادله، $V(D, G)$ همان تابع هدف کلاسیک \lr{GAN} است که در آن تمایزگر ($D$) می‌کوشد داده‌های واقعی را از خروجی‌های مصنوعی تفکیک کند و مولد ($G$) سعی در فریب دادن $D$ دارد. نوآوری بنیادین در عبارت دوم، یعنی $-\lambda I(c; G(z, c))$، نهفته است که در آن $I(\cdot ; \cdot)$ نشان‌دهنده‌ی اطلاعات متقابل و $\lambda > 0$ ضریب منظم‌سازی است. 

از منظر شهودی، اطلاعات متقابل $I(c; G(z, c)) = H(c) - H(c | G(z, c))$ میزان کاهش عدم‌قطعیت درباره‌ی کد پنهان $c$ را پس از مشاهده‌ی نمونه‌ی تولیدشده‌ی $G(z, c)$ اندازه‌گیری می‌کند. از آنجا که مولد در معادله‌ی~\ref{eq:chap3_gan} تابع هدف را کمینه‌سازی می‌کند، وجود علامت منفی در پشت این عبارت به این معناست که $G$ ملزم به \textit{بیشینه‌سازی} اطلاعات متقابل میان کدهای پنهان و داده‌های خروجی است. به بیان ساده، اگر مولد کد $c$ را نادیده بگیرد، نمونه‌های تولیدی حاوی هیچ اطلاعاتی از $c$ نخواهند بود و در نتیجه $I(c; G(z, c))$ به صفر میل می‌کند. اما با بیشینه‌سازی این مقدار، شبکه مجازات سنگینی در صورت نادیده گرفتن $c$ متحمل می‌شود؛ در نتیجه ناچار می‌گردد ارتباطی صریح و پایدار میان ابعاد $c$ و ویژگی‌های بارز خروجی برقرار سازد، به‌طوری‌که یک ناظر بتواند صرفاً با نگاه به خروجی تولیدی، کد پنهان اولیه‌ی $c$ را بازسازی کند. از آنجا که محاسبه‌ی دقیق این اطلاعات متقابل نیازمند توزیع پسین ناشناخته‌ی $P(c|x)$ است و در عمل حل‌ناپذیر است، چن و همکاران با تعریف یک شبکه‌ی فرعی ($Q$)، یک کران پایین متغیر\enfootnote{variational lower bound} برای آن بهینه‌سازی می‌کنند تا ابعاد $c$ به ویژگی‌های معنایی مجزا و پیوسته (مانند چرخش، ضخامت یا زاویه‌ی تابش) تبدیل گردند.

به‌موازات این پیشرفت‌ها، تلاش‌های مشابهی در حوزه‌ی خودرمزگذارهای متغیر\enfootnote{variational autoencoders} (\lr{VAE}) نیز صورت گرفت. هیگینز و همکاران (۲۰۱۷)~\cite{higgins2017betavae} برای دستیابی به فضاهای پنهان تفسیرپذیر، چارچوب \lr{$\beta$-VAE} را ارائه دادند. در فرمول‌بندی استانداردِ \lr{VAE}، تابع هدف بر پایه‌ی بهینه‌سازی کران پایین شواهد\enfootnote{evidence lower bound (\lr{ELBO})} تعریف می‌شود که شامل درست‌نمایی بازسازی به‌همراه یک جریمه‌ی منظم‌ساز بر اساس واگرایی کولبک-لایبلر است تا توزیع پنهان را به یک توزیع نرمال استاندارد نزدیک کند. نوآوری \lr{$\beta$-VAE} در این بود که با اعمال یک ضریب وزنی تشدیدیافته ($\beta > 1$) بر روی عبارت منظم‌ساز، تابع هدف به‌صورت زیر اصلاح گردید:
\begin{equation}\label{eq:chap3_beta_vae}
    \mathcal{L}(\theta, \phi; x, z) = \mathbb{E}_{q_\phi(z|x)}[\log p_\theta(x|z)] - \beta D_{KL}(q_\phi(z|x) \parallel p(z))
\end{equation}
در این معادله، $\mathcal{L}(\theta, \phi; x, z)$ کران پایین تعمیم‌یافته‌ای است که مدل در فرایند آموزش به دنبال بیشینه‌سازی آن (یا کمینه‌سازی معکوس آن به‌عنوان تابع زیان) است. این رابطه تعادلی پویا میان دو نیروی متضاد را مدل‌سازی می‌کند:
\begin{enumerate}
    \item عبارت اول، $\mathbb{E}_{q_\phi(z|x)}[\log p_\theta(x|z)]$، امید ریاضی درست‌نمایی بازسازی داده است که عملکرد رمزگشا ($p_\theta$) را در بازتولید وفادارانه‌ی داده‌ی ورودی $x$ از روی متغیر پنهان $z$ (استخراج‌شده توسط رمزگذار $q_\phi$) می‌سنجد. بیشینه‌سازی این عبارت به‌تنهایی باعث می‌شود شبکه تمایل یابد کدهایی با ابعاد پیچیده و درهم‌تنیده بیاموزد تا کوچک‌ترین جزئیات ورودی بازسازی شوند.
    \item عبارت دوم، $D_{KL}(q_\phi(z|x) \parallel p(z))$، فاصله‌ی آماری میان توزیع پنهان استنتاج‌شده ($q_\phi$) و توزیع پیشین ($p(z) \sim \mathcal{N}(0, I)$) را اندازه می‌گیرد. نکته‌ی کلیدی در اینجاست که توزیع پیشین فرضی، یک توزیع نرمال با ماتریس کوواریانس قطری و همانی است؛ یعنی متغیرهای آن کاملاً مستقل از یکدیگرند.
\end{enumerate}

منطق شهودی حضور ضریب $\beta > 1$ در ایجاد یک گلوگاه اطلاعاتی\enfootnote{information bottleneck} فشرده نهفته است. در یک \lr{VAE} عادی ($\beta = 1$)، وزن بازسازی بالاتر بوده و شبکه می‌تواند همبستگی‌های پنهان را حفظ کند. اما هنگامی که ضریب جریمه‌ی $\beta$ به مقادیر بزرگ‌تر از یک افزایش می‌یابد، پهنای باند کانال انتقال اطلاعات به‌شدت محدود می‌شود. در چنین شرایطی، شبکه دیگر گنجایش کافی برای عبور دادن ویژگی‌های همبسته و نویزهای فرعی را ندارد و از نظر ریاضی ناچار می‌شود کارآمدترین، مستقل‌ترین و عام‌ترین مؤلفه‌های داده (مانند اندازه، موقعیت مکانی و جهت) را کشف کند و هر مؤلفه را به یک محور مختصات مستقل در فضای پنهان تخصیص دهد. البته بهای این تفکیک و استقلال ابعاد، افت اجتناب‌ناپذیر کیفیت بازسازی و ایجاد تاری در خروجی‌ها به‌دلیل از دست رفتن جزئیات فرعی است.

اگرچه این روش‌ها در ایزوله کردن مفاهیم پایه در داده‌های بصریِ ساده موفق عمل کردند، اما امیدواری به ساخت مدل‌های مولد کاملاً شفاف با انتشار مطالعه‌ی نظری لوکاتلو و همکاران (۲۰۱۹)~\cite{locatello19challenging} با چالش‌های جدی مواجه شد. این پژوهش با یک اثبات جامع ریاضی نشان داد که یادگیری بدون‌نظارت بازنمایی‌های تفکیک‌شده، بدون اعمال سوگیری‌های استقرایی صریح بر روی مدل یا داده‌ها، از اساس غیرممکن است. پیامد این قضیه برای \lr{LLM}ها بسیار حائز اهمیت است؛ از آنجا که این مدل‌ها به‌صورت کاملاً خودنظارتی بر روی حجم عظیمی از متون بدون برچسب در مقیاس وب آموزش می‌بینند، انتظار اینکه بتوانند به‌طور خودکار و ارگانیک فضاهای پنهان شفاف و تفکیک‌شده تولید کنند، فاقد پشتوانه ریاضی است. به‌ویژه آنکه خروجی این مدل‌ها دنباله‌ای گسسته از نشانه‌هاست و نگاشت یک فضای پنهان پیوسته به مفاهیم انتزاعی زبانی بدون نظارت صریح، به لحاظ نظری نامعین باقی می‌ماند؛ زیرا خمینه‌های زبانی به‌شدت ناپیوسته هستند و یک انتقال هموار در فضای هندسی پنهان، لزوماً معادل یک انتقال منطقی و هموار در ساختار معنایی زبان نیست.

در پاسخ به این محدودیت نظری، پژوهش‌ها به سمت تزریق دانش پیشین تغییر مسیر دادند. برای مثال، موران و آراگام (۲۰۲۶)~\cite{moran2026towardinterpretable} با تحلیل آماری حوزه‌ی یادگیری بازنمایی علّی\enfootnote{causal representation learning}، شرایط ریاضی و قضایای همانندپذیری\enfootnote{identifiability} را برای ادغام متغیرهای پنهان با گراف‌های علّی ساختاری تدوین کردند تا مدل بتواند متغیرهای حقیقی تولیدکننده‌ی داده را به‌طور قطعی بازیابی نماید. با این وجود، این چارچوب‌ها وابستگی مطلقی به داشتن دانش پیشین ساختاریافته پیرامون نمودار و توزیع علّی داده‌ها دارند؛ فرضی که برای دامنه‌ی باز و بی‌کران زبان طبیعی اساساً ناممکن است.

علاوه بر الزامات آموزشی، اعمال سوگیری‌های استقرایی با یک چالش ماهوی دیگر در حوزه‌ی پردازش متن روبه‌روست. در بینایی ماشین، محققانی نظیر ژو و همکاران (۲۰۲۱)~\cite{zhu2021whereandwhat} با استفاده از ماسک‌های مکانی\enfootnote{spatial masks} توانستند اثر هر بُعد پنهان را به یک ناحیه‌ی خاص از تصویر محدود کنند و شفافیت را افزایش دهند. اما موفقیت این تکنیک‌ها به‌شدت وابسته به پیوستگی هندسی و دوبعدی تصاویر است. در مقابل، خروجی مدل‌های زبانی ماهیتی گسسته و متوالی دارد که توکن‌به‌توکن تولید می‌شود. با این حال، در زبان طبیعی، یک مفهوم معنادار (مانند یک استدلال منطقی یا پیش‌فرض) در سطح قطعات خُرد نحوی شکل نمی‌گیرد، بلکه دارای ماهیتی گزاره‌ای است. بنابراین، پدیده‌ای نظیر تغییر و انتقال مفاهیم در متون، به‌صورت توزیع‌شده در سراسر جملات رخ می‌دهد و تحلیل آن در سطح توکن یا با اعمال محدودیت‌های مکانی ساده امکان‌پذیر نیست.

در مجموع، بررسی این مسیر پژوهشی نشان می‌دهد که تلاش برای دستیابی به تفسیرپذیری از طریق یادگیری بدون‌نظارت فضاهای پنهان، به‌ویژه برای زبان، با بن‌بست روبه‌روست. این روش‌ها نه تنها نیازمند تغییرات بنیادین در تابع زیان طی فاز پیش‌آموزش هستند، بلکه نمی‌توانند تضمین کنند که ابعاد تفکیک‌شده، با ساختار گزاره‌ای و درک انسانی از متن هم‌تراز باشند. این ناتوانی، محققان را ناگزیر ساخت تا از رویکردهای بدون‌نظارت و ضمنی عبور کرده و به سمت روش‌های نظارت‌شده و صریح حرکت کنند؛ جایی که مفاهیم از پیش توسط انسان تعریف شده و مستقیماً به مدل دیکته می‌شوند.


\subsubsection{مدل‌های مبتنی بر تعریف مفهوم}\label{sec:3-1-1-3}

برای غلبه بر چالش فقدان معنا در بازنمایی‌های تفکیک‌شده، جریان قدرتمند دیگری شکل گرفت که تلاش می‌کرد شفافیت را از طریق اجبار شبکه به پیش‌بینی مفاهیم انسانی و از پیش تعریف‌شده محقق سازد. برجسته‌ترین نمونه در این پارادایم، مدل‌های گلوگاه مفهومی (\lr{CBM}) است که توسط کو و همکاران (۲۰۲۰)~\cite{koh2020cbm} برای وظایف متمایزگر معرفی شد.

ایده‌ی محوری در این معماری، تغییر مسیر سنتی پیش‌بینی و شکستن آن به دو مرحله‌ی مجزا و پی‌درپی است. در شبکه‌های طبقه‌بند متداول، ورودی خام $x \in \mathcal{X}$ مستقیماً به برچسب هدف نهایی $y \in \mathcal{Y}$ نگاشت می‌شود ($x \rightarrow y$). اما در \lr{CBM}، یک لایه‌ی میانی متشکل از $k$ مفهوم معنادار انسانی نظیر $c = (c_1, c_2, \dots, c_k) \in \mathcal{C}$ به‌عنوان یک گلوگاه محاسباتی تعبیه می‌شود. این فرایند دوگانه را می‌توان با روابط ریاضی زیر نمایش داد:
\begin{equation}\label{eq:chap3_cbm}
    \hat{c} = g(x), \quad \hat{y} = f(\hat{c})
\end{equation}
در این فرمول‌بندی، تابع $g: \mathcal{X} \rightarrow \mathcal{C}$ (معمولاً یک شبکه‌ی عصبی عمیق) وظیفه‌ی استخراج مفاهیم میانی و پیش‌بینی حضور آن‌ها را بر عهده دارد (به‌عنوان مثال، $\hat{c}_j \in [0, 1]$ احتمال حضور یک ویژگی ظاهری را نشان می‌دهد). سپس تابع $f: \mathcal{C} \rightarrow \mathcal{Y}$ که غالباً یک مدل خطی شفاف یا شبکه‌ای کم‌عمق است، تصمیم نهایی $\hat{y}$ را \textit{منحصراً} بر اساس این مفاهیم پیش‌بینی‌شده استنتاج می‌کند، بدون آنکه دسترسی مستقیمی به ویژگی‌های خام ورودی $x$ داشته باشد. آموزش این شبکه می‌تواند به‌صورت توأم با تابع زیان زیر انجام گیرد:
\begin{equation}\label{eq:chap3_cbm_loss}
    \mathcal{L}_{\text{\lr{CBM}}} = \mathcal{L}_y(f(\hat{c}), y) + \lambda \sum_{j=1}^k \mathcal{L}_c(\hat{c}_j, c_j)
\end{equation}
که در آن $\mathcal{L}_c$ خطای پیش‌بینی تک‌تک مفاهیم، $\mathcal{L}_y$ خطای طبقه‌بندی وظیفه‌ی اصلی و $\lambda$ ضریب تعادل میان یادگیری مفاهیم و برچسب نهایی است. 

به زبان شهودی، شفافیت ذاتی این مدل ناشی از آن است که استدلال نهایی کاملاً روی میز و در قالب مفاهیم زبانی و انسانی صورت می‌گیرد؛ برای نمونه در تشخیص نوع یک پرنده، مدل ابتدا ویژگی‌هایی نظیر «منقار زرد»، «بال سیاه» و «طول دم بلند» را پیش‌بینی می‌کند و سپس دسته‌بندی نهایی را بر مبنای وزن و ترکیب این ویژگی‌ها انجام می‌دهد. مهم‌ترین مزیت عملیاتی این ساختار، قابلیت‌مداخله در زمان آزمون\enfootnote{test-time intervention} است: اگر مدل در پیش‌بینی یکی از مفاهیم دچار لغزش شود (مثلاً ویژگی «منقار زرد» را به اشتباه صفر تخمین بزند)، یک متخصص می‌تواند بدون نیاز به آموزش مجدد شبکه، مقدار این گره را در بردار مفاهیم اصلاح کند ($\hat{c}_{\text{\lr{int}}}$)؛ با این تصحیح، لایه‌ی نهایی $f$ تصمیم خود را بلافاصله بر مبنای مقدار تصحیح‌شده به‌روزرسانی می‌کند.

با این وجود، این چارچوب با چالش‌های بنیادینی روبه‌روست که مسیر تکامل آن را شکل داد:

\paragraph{۱. مصالحه‌ی عملکرد-شفافیت:} وادار کردن شبکه به فشرده‌سازی تمام اطلاعات یک ورودی پیچیده به چند عدد اسکالر ساده در لایه‌ی گلوگاه، باعث افت شدید دقت مدل نسبت به شبکه‌های دسته‌بند معمولی می‌شد. برای رفع این مشکل، اسپینوزا زارلنگا و همکاران (۲۰۲۲)~\cite{zarlenga2022cem} معماری مدل‌های تعبیه‌ی مفهوم\enfootnote{concept embedding models} (\lr{CEM}) را پیشنهاد دادند. آن‌ها استدلال کردند که یک عدد به‌تنهایی نمی‌تواند تمام ظرایف را در خود نگه دارد؛ لذا برای هر مفهوم، به‌جای یک مقدار اسکالر، یک بردار با ابعاد بالا (تعبیه) تولید کردند. این بردارها به مدل اجازه می‌دادند تا علاوه بر خود مفهوم، اطلاعات جانبی را نیز رمزگذاری کنند. اگرچه این نوآوری دقت مدل را به سطح شبکه‌های استاندارد بازگرداند، اما به قیمت نشت اطلاعات\enfootnote{information leakage} تمام شد. در واقع، مدل یاد گرفت اطلاعات مفیدی را که فاقد معنای انسانی بودند در فضای وسیع این بردارها پنهان کرده و از گلوگاه عبور دهد. در نتیجه، تصمیم نهایی دیگر منحصراً وابسته به مفاهیم شفاف نبود و ذات تفسیرپذیری درونی مخدوش شد.

\paragraph{۲. شکنندگی در برابر داده‌های پرنویز:} کیم و همکاران (۲۰۲۳)~\cite{kim2023probabilistic} نشان دادند که پیش‌بینی یک عدد قطعی برای مفاهیم در دنیای واقعی کارآمد نیست. برای مثال اگر بخشی از تصویر تار باشد، \lr{CBM} کلاسیک همچنان یک عدد قطعی تولید می‌کند و کاربر متوجه خطای احتمالی نمی‌شود. آن‌ها با معرفی \lr{ProbCBM}، مدل را وادار به یادگیری توزیع‌های احتمالی برای مفاهیم کردند. این کار به مدل اجازه می‌دهد تا میزان عدم‌قطعیت خود را در تشخیص یک ویژگی اعلام کند (مثلاً بگوید با واریانس بالا حدس می‌زند منقار زرد است) تا کاربر بداند چه زمانی نباید به گلوگاه اعتماد کند.

\paragraph{۳. فراموشی مداخلات انسانی:} در مدل‌های استاندارد، اگر کاربر با استفاده از قابلیت‌مداخله اشتباه مدل را تصحیح کند، این اصلاح فقط برای همان یک نمونه اعمال می‌شود و در صورت رویارویی با داده‌ی مشابه در آینده، مدل دوباره همان اشتباه را تکرار می‌کند. اشتاینمن و همکاران (۲۰۲۴)~\cite{steinmann2024learning} با توسعه‌ی چارچوب \lr{CB2M} و افزودن یک ماژول حافظه‌ی دوگانه، مداخلات انسانی را ذخیره کرده تا در رویارویی مجدد با نمونه‌های مشابه، اصلاحات به‌صورت خودکار اعمال شوند.

\paragraph{۴. تأمین و تعریف مفاهیم:} یکی از جدی‌ترین محدودیت‌های پارادایم گلوگاه، وابستگی مطلق به فرایند استخراج و برچسب‌گذاری مفاهیم است. در رویکردهای کلاسیک، این کار نیازمند مجموعه‌داده‌هایی بود که مفاهیم در آن‌ها توسط انسان برچسب‌گذاری شده باشند؛ فرایندی که حتی برای داده‌های بصری نیز پرهزینه و مقیاس‌ناپذیر است. برای خودکارسازی این روند، روش‌هایی نظیر \lr{Label-free CBM}~\cite{oikarinen2023labelfree} در بینایی ماشین (که اساساً یک رویکرد پساعملی است، اما در اینجا به‌دلیل ابتکار در استخراج خودکار مفهوم ذکر می‌شود) و فاز نخست چارچوب \lr{CB-LLM}~\cite{sun2025concept} در پردازش زبان معرفی شدند که از مدل‌های قدرتمندی نظیر \lr{GPT}~\cite{brown2020language} برای استخراج مجموعه‌ی مفاهیم و امتیازدهی به متون استفاده می‌کنند.

با این حال، این راه‌حل نیز با چالش‌های عمیقی مواجه است. نخست آنکه ابهام در تعریف ماهوی مفهوم مانع از آن می‌شود که بتوان مرز دقیقی برای مفاهیم استخراج‌شده قائل شد. دوم آنکه مفاهیم استخراج‌شده کاملاً به دانش، سوگیری‌ها و نقاط کور مدل خارجی وابسته‌اند و ریسک بروز توهم در لایه‌ی گلوگاه را افزایش می‌دهند. این بحران تعریف، در مواجهه با زبان طبیعی به‌مراتب حادتر می‌شود، چرا که دامنه‌ی مفاهیم منطقی و انتزاعی در زبان برخلاف ویژگی‌های فیزیکی تصاویر، بی‌نهایت است.

\paragraph{۵. تطبیق با معماری تولید متن:} چالش دیگر، تفاوت ساختاری میان مدل‌های متمایزگر و مولد است. در یک دسته‌بند، می‌توان گلوگاه را به‌سادگی پیش از لایه‌ی تصمیم‌گیرنده‌ی نهایی قرار داد؛ اما در مدل‌های زبانی ترنسفورمری، استنتاج یک فرایند متوالی و خودهمبسته در سطح نشانه‌هاست و نقطه‌ی واحدی برای اعمال گلوگاه وجود ندارد. سان و همکاران (۲۰۲۵)~\cite{sun2025concept} در چارچوب \lr{CB-LLM} تلاش کردند با تعبیه‌ی یک لایه‌ی گلوگاه مفهومی اختصاصی (\lr{CBL}) با نورون‌های تفسیرپذیر و آموزش آن‌ها به شیوه‌ی متضاد، نورون‌ها را به‌عنوان نماینده‌ی مفاهیم مشخص درآورند تا کاربر بتواند با دستکاری فعال‌سازی این نورون‌ها در حین استنتاج خودهمبسته، مسیر تولید مدل را هدایت کند. با این حال، از آنجا که محدود کردن کامل فرایند تولید توکن‌به‌توکن به چند مفهوم، روانی متن را مخدوش می‌ساخت، آن‌ها ناچار شدند یک مسیر موازی بدون‌نظارت را در کنار گلوگاه حفظ نمایند. این امر نشان می‌دهد که اجبار نورون‌های داخلی به هم‌ترازی یک‌به‌یک با مفاهیم انسانی، با ویژگی بنیادین چندمعنایی\enfootnote{polysemanticity} در شبکه‌های عصبی بزرگ در تضاد است؛ پدیده‌ای که در آن هر نورون منفرد، بسته به بافتار، در بازنمایی توزیع‌شده‌ی چندین مفهوم مختلف مشارکت دارد~\cite{elhage2022toymodels}.

در کنار این تلاش‌ها، رویکرد متفاوتی برای غلبه بر بحران تعریف مفاهیم و مشکل تولید توکن‌به‌توکن در مدل‌های مفهومی بزرگ\enfootnote{large concept models} (\lr{LCM}) ارائه شده است~\cite{lcmteam2024lcm}. برخلاف ترنسفورمرهای سنتی، \lr{LCM} زبان را به‌جای سطح توکن، در یک فضای بازنمایی پیوسته‌ی سطح جمله مدل‌سازی می‌کند. در این معماری، مدل ابتدا بر روی مفاهیم در فضای پنهان استنتاج کرده و تعبیه‌ی معنایی جمله‌ی بعدی را پیش‌بینی می‌کند. سپس برای تبدیل این بازنمایی‌های انتزاعی به متن قابل‌فهم، از یک رمزگشای غیرخودهمبسته (یا مدل‌های انتشاری\enfootnote{diffusion models}) استفاده می‌شود که وظیفه‌ی نگاشت بردار مفهومی به دنباله‌ای از کلمات را بر عهده دارد. این تفکیک میان برنامه‌ریزی مفهومی و تولید نحوی، نه تنها نیاز به برچسب‌گذاری خارجی مفاهیم را از بین می‌برد، بلکه نشان می‌دهد که تحلیل استدلال‌ها در سطح گزاره‌ای پایدارتر از تحلیل‌های سطح توکن است.

با این وجود، علی‌رغم نوآوری در بازتعریف واحد محاسباتی، \lr{LCM}ها با چالش‌های جدی در زمینه‌ی عملکرد و تفسیرپذیری واقعی روبه‌رو هستند. نخست آنکه در مقایسه با مدل‌های زبانی استاندارد، این مدل‌ها در وظایف استدلالی پیچیده همچنان دقت پایین‌تری نشان می‌دهند. دوم آنکه اگرچه تغییر از توکن به جمله، فرایند استنتاج را با ساختار گزاره‌ای ذهن انسان هم‌سو می‌کند، اما بردارهای پنهان در این فضای معنایی همچنان ماهیتی متراکم و غیرشفاف دارند؛ به‌عبارتی، \lr{LCM} تنها واحد پردازش را تفسیرپذیرتر کرده است، اما محتوای درونی هر مفهوم همچنان برای ناظر انسانی به‌سادگی قابل‌تجزیه نیست. در نهایت، وابستگی مطلق به یک رمزگشای ثانویه برای مشاهده‌ی خروجی، ریسک بروز خطا در مرحله‌ی نگاشت معنا به متن را افزایش داده و مانع از نظارت مستقیم بر فرایند استدلال درونی مدل می‌شود.


\subsubsection{چالش‌های بنیادین}\label{sec:3-1-1-4}

مجموع بررسی‌های فوق نشان می‌دهد که پارادایم تفسیرپذیری درونی، علی‌رغم نوآوری‌های چشمگیر، در مواجهه با اکوسیستم فعلی هوش مصنوعی به یک بن‌بست عملی رسیده است. این بن‌بست ریشه در سه محدودیت بنیادین دارد.

نخستین مانع، مصالحه‌ی گریزناپذیر میان عملکرد و شفافیت است. طراحی معماری‌هایی که فرایند تصمیم‌گیری آن‌ها ذاتاً تفسیرپذیر باشد، مستلزم اعمال محدودیت‌های ساختاری و اطلاعاتی شدید بر شبکه است. این محدودیت‌ها در مدل‌های به‌خصوص زبانی باعث افت محسوس دقت، کاهش انعطاف‌پذیری و تنزل کیفیت متن تولیدی می‌شوند، به‌گونه‌ای که این مدل‌های شفاف توان رقابت با همتایان جعبه‌سیاه خود را در وظایف پیچیده از دست می‌دهند.

دومین مانع، تقلیل‌گرایی در بازنمایی زبان است. رویکردهای درونی کلاسیک غالباً تلاش می‌کنند مفاهیم پیچیده را در گره‌های محاسباتی منفرد و ایزوله (مانند یک نورون خاص یا یک بُعد پنهان مجزا) فشرده سازند. با این حال در مدل‌های زبانی بزرگ، مفاهیم دارای بازنمایی‌های توزیع‌شده در فضا‌های برداری فعال‌سازی هستند و معانی در سطوح بالاتر گزاره‌ای متبلور می‌شوند~\cite{zou2023representationengineering}. همان‌طور که فریدمن و همکاران (۲۰۲۴)~\cite{friedman2024interpretability} به‌صورت تجربی اثبات کرده‌اند، تلاش برای تقلیل و فشرده‌سازی این بازنمایی‌های غنی به واحدهای ساده‌شده‌ی تفسیرپذیر، زمینه‌ساز بروز توهم تفسیرپذیری می‌گردد؛ چرا که این مدل‌های ساده‌شده اگرچه در داده‌های توزیع آموزش وفادار به نظر می‌رسند، اما در تعمیم به شرایط خارج از توزیع و انعکاس استدلال‌های چندوجهی مدل اصلی به‌شدت شکست می‌خورند.

سومین و مهم‌ترین مانع، تضاد هدف این پارادایم با واقعیت کاربرد مورد نیاز اکثر کاربران است. رسالت تفسیرپذیری درونی، طراحی و آموزش مدل‌های جدید از پایه است و اساساً برای تفسیر مدل‌های پیش‌آموزش‌دیده طراحی نشده است. امروزه پیشرفته‌ترین مدل‌های زبانی که به‌طور گسترده در کاربردهای حساس استفاده می‌شوند، سیستم‌هایی عظیم و منجمد هستند که پشت واسط‌های برنامه‌نویسی محافظت می‌شوند. از آنجا که امکان بازطراحی معماری یا آموزش مجدد این مدل‌های جعبه‌سیاه وجود ندارد، رویکردهای درونی برای ممیزی سیستم‌های مسلط امروزی کاملاً بی‌کاربرد هستند.

این محدودیت‌های نظری و ساختاری اثبات می‌کنند که برای ارزیابی مدل‌های زبانی تجاری و پیشرفته، باید مدل را به‌عنوان یک موجودیت ثابت در نظر گرفت و رفتار آن را پس از اتمام آموزش تفسیر کرد. این ضرورت عملی، پایه‌گذار پارادایم تفسیرپذیری پساعملی گردید که چارچوب اصلی آن در بخش بعدی بررسی خواهد شد.


\subsection{تفسیرپذیری پساعملی: رویکردهای جعبه‌سفید}\label{sec:3-1-2}

با توجه به محدودیت‌های بنیادین و افت عملکرد ناشی از اعمال محدودیت‌های ساختاری در رویکردهای تفسیرپذیری درونی، تمرکز جامعه‌ی تحقیقاتی به سمت پارادایم تفسیرپذیری پساعملی معطوف گردید. در این پارادایم، مدل پیچیده (نظیر یک شبکه‌ی عصبی عمیق یا \lr{LLM}) به‌طور کامل آموزش می‌بیند و هیچ‌گونه محدودیتی بر معماری یا ظرفیت یادگیری آن اعمال نمی‌شود. پس از اتمام آموزش و ثابت شدن وزن‌ها، از ابزارها و الگوریتم‌های ثانویه برای تحلیل رفتار، استخراج منطق تصمیم‌گیری و درک مکانیسم‌های داخلی مدل استفاده می‌گردد.

رویکردهای پساعملی بر اساس میزان دسترسی به اطلاعات درون مدل، به دو دسته‌ی کلی جعبه‌سفید و جعبه‌سیاه تقسیم می‌شوند~\cite{guidotti2018survey, danilevsky2020survey}. روش‌های جعبه‌سفید بر این پیش‌فرض استوارند که تحلیل‌گر به تمامی اجزای داخلی مدل هدف – شامل معماری، وزن‌ها، توابع فعال‌سازی و قابلیت‌محاسبه‌ی گرادیان‌ها – دسترسی کامل دارد. در این بخش، جریان‌های تحقیقاتی در پارادایم جعبه‌سفید مورد بررسی و نقد قرار می‌گیرند تا مشخص شود چگونه دسترسی به سیگنال‌های داخلی، درک ما را از رفتار مدل ارتقا داده است.


\subsubsection{انتساب اعتبار مبتنی بر سیگنال‌های پایه}\label{sec:3-1-2-1}

یکی از نخستین و شهودی‌ترین پرسش‌ها در ارزیابی یک مدل یادگیری ماشین این است که: «کدام بخش‌های ورودی، بیشترین سهم را در تولید این خروجی خاص داشته‌اند؟». پاسخ به این پرسش در قالب مسئله‌ی انتساب اعتبار\enfootnote{credit attribution} یا اهمیت ویژگی\enfootnote{feature importance} فرمول‌بندی می‌شود~\cite{danilevsky2020survey}. در شبکه‌های عصبی، مستقیم‌ترین راه برای سنجش این حساسیت، استفاده از سیگنال‌های پایه‌ای شبکه‌ی عصبی، یعنی مشتقات (گرادیان‌ها) و وزن‌های مسیریابی (مانند مکانیزم توجه) است.

\paragraph{انتساب مبتنی بر گرادیان.}

استفاده از گرادیان به‌عنوان معیار اهمیت، ریشه در کارهای بنیادین حوزه‌ی بینایی ماشین دارد. سیمونیان و همکاران (۲۰۱۴)~\cite{simonyan2014deepinside} با معرفی نقشه‌های برجستگی\enfootnote{saliency maps}، نشان دادند که با محاسبه‌ی مشتق جزئیِ امتیاز کلاس خروجی ($y$) نسبت به پیکسل‌های تصویر ورودی ($x_i$)، یعنی $\frac{\partial y}{\partial x_i}$، می‌توان مشخص کرد که تغییرات اندک در کدام پیکسل‌ها، بیشترین تأثیر را بر تصمیم مدل می‌گذارد. به بیان شهودی، این گرادیان محلی نشان می‌دهد که اگر روشنایی یا مقدار یک پیکسل را ذره‌ای تغییر دهیم، امتیاز پیش‌بینی مدل با چه شیبی دگرگون خواهد شد.

در همین راستا و برای غلبه بر نویزِ نقشه‌های برجستگی خام در تصاویر، سلواراجو و همکاران (۲۰۱۷)~\cite{selvaraju2017gradcam} چارچوب \lr{Grad-CAM} را معرفی کردند. این روش به‌جای تمرکز بر تک‌تک پیکسل‌های خام، گرادیان امتیاز کلاس هدف ($y^c$) را نسبت به نقشه‌های ویژگی آخرین لایه‌ی پیچشی ($A^k$) محاسبه می‌کند. وزن اهمیت هر نقشه‌ی ویژگی ($\alpha_k^c$) از طریق تجمیع میانگین سراسری\enfootnote{global average pooling} به‌دست می‌آید:
\begin{equation}\label{eq:chap3_grad_cam}
    \alpha_k^c = \frac{1}{Z}\sum_i\sum_j \frac{\partial y^c}{\partial A_{i,j}^k}
\end{equation}
که در آن $Z$ تعداد نورون‌های موجود در نقشه‌ی ویژگی (حاصل‌ضرب ابعاد طول و عرض) است. از نظر شهودی، عبارت $\frac{\partial y^c}{\partial A_{i,j}^k}$ بیانگر حساسیت محلی امتیاز کلاس به فعال‌سازی مکان $(i, j)$ در نقشه‌ی $k$ام است؛ با میانگین‌گیری از این مقادیر در کل پهنه‌ی تصویر، وزن $\alpha_k^c$ اهمیت سراسری آن الگوی بصری را برای کلاس هدف $c$ خلاصه می‌کند. سپس نقشه‌ی نهایی \lr{Grad-CAM} از ترکیب خطی این نقشه‌ها و عبور از یک تابع یکسوکننده‌ی خطی (\lr{ReLU}) حاصل می‌شود:
\begin{equation}\label{eq:chap3_grad_cam_map}
    L_{\text{\lr{Grad-CAM}}}^c = \text{\lr{ReLU}}\left( \sum_k \alpha_k^c A^k \right)
\end{equation}
کاربرد شهودی تابع \lr{ReLU} در این رابطه، پالایش الگوهاست؛ به این معنا که تنها الگوهایی حفظ می‌شوند که اثر مثبت در افزایش احتمال کلاس هدف دارند و ویژگی‌هایی که متعلق به کلاس‌های دیگر بوده یا تصمیم مدل را تضعیف می‌کنند، حذف می‌گردند. در نتیجه، نواحی تعیین‌کننده با وضوحی کمتر از سطح پیکسل اما با تفکیک معنایی غنی، به‌صورت نقشه‌ی حرارتی برجسته می‌شوند.

این منطق ریاضی به‌سرعت به پردازش زبان طبیعی نیز تعمیم یافت، اما پیاده‌سازی آن با یک چالش ماهوی روبه‌رو بود؛ برخلاف پیکسل‌های تصویر، کلمات ماهیتی گسسته دارند و مشتق‌گیری مستقیم از آن‌ها امکان‌پذیر نیست. لی و همکاران (۲۰۱۶)~\cite{li2016visualizing} برای حل این مشکل نشان دادند که در وظایف زبانی، مشتق باید نسبت به تعبیه‌های پیوسته‌ی کلمات محاسبه شود. در این روش که برای وظایف متمایزگر (نظیر تحلیل احساسات) توسعه یافت، یک رویکرد استاندارد برای محاسبه‌ی اهمیت کلمه‌ی $i$ام، استفاده از تکنیک گرادیان ضرب‌در ورودی\enfootnote{gradient $\times$ input} است:
\begin{equation}\label{eq:chap3_grad_x_input}
    I(x_i) = {\left( \frac{\partial y}{\partial e_i} \right)}^T \cdot e_i
\end{equation}
که در آن $e_i$ بردار تعبیه‌ی کلمه و $y$ تابع لوجیت کلاس نهایی است. از منظر شهودی، گرادیان $\frac{\partial y}{\partial e_i}$ صرفاً حساسیت نظری مدل به ابعاد تعبیه را ارزیابی می‌کند؛ اما ممکن است مدلی به یک جهت حساس باشد در حالی که آن ویژگی اصلاً در ورودی حضور نداشته باشد. ضرب داخلی گرادیان در خود بردار تعبیه ($e_i$) که تقریب مرتبه‌اول بسط تیلور تابع تصمیم در مبدأ است ($y(e_i) - y(0) \approx \nabla_{e_i} y \cdot e_i$)، میزان حساسیت را با شدت حضور واقعی کلمه وزن‌دهی می‌کند و یک امتیاز واحد (اسکالر) به‌دست می‌دهد که سهم فعال آن کلمه را در تصمیم نهایی بازتاب می‌دهد.

با ظهور \lr{LLM}ها، این رویکرد به وظایف مولد و تولید متن نیز بسط یافت. در تولید متن، از آنجا که خروجی یک کلاسِ منفرد نیست بلکه دنباله‌ای از کلمات است، محاسبه‌ی میزان اثرگذاری پیش‌آیندها باید به‌صورت گام‌به‌گام انجام شود. سارتی و همکاران (۲۰۲۳)~\cite{sarti2023inseq} با درک این پیچیدگی و برای پر کردن شکاف نظری میان مدل‌های دسته‌بندی و تولید توالی، چارچوب و کتابخانه‌ی جامع \lr{Inseq} را توسعه دادند. آن‌ها فرمول‌بندی کردند که چگونه می‌توان میزان تأثیر توکن‌های ورودی ($X$) را به‌عنوان پیش‌آیند در تولید توکن هدف در گام $t$ ($y_t$) برای مدل‌های خودهمبسته سنجید:
\begin{equation}\label{eq:chap3_inseq}
    I(x_i \rightarrow y_t) = {\left( \frac{\partial \log P(y_t \mid \mathbf{x}, y_{<t})}{\partial e_i} \right)}^\top \cdot e_i
\end{equation}
منطق شهودی این فرمول در آن است که هدف انتساب در هر گام از رمزگشایی، میزان افزایش لگاریتم احتمال تولید توکن $y_t$ با شرطِ حضور پیشوند ($y_{<t}$) و پرامپت ($X$) است. مشتق جزئی مشخص می‌کند که تغییر در تعبیه‌ی کلمه‌ی ورودی $x_i$ چه اثری بر گرایش مدل به انتخاب کلمه‌ی $y_t$ دارد و ضرب آن در $e_i$، سهم مستقیم آن کلمه‌ی ورودی را در این رویداد زبانی کمّی می‌سازد. این چارچوب نه‌تنها تکنیک پایه‌ی روش گرادیان ضرب‌در ورودی را برای متون پیاده‌سازی کرد، بلکه روش‌های پیشرفته‌تری نظیر \lr{IG} و \lr{DeepLIFT} را نیز به‌صورت بومی برای فرایند رمزگشایی گام‌به‌گام در معماری‌های مختلف بازنویسی و استانداردسازی نمود.

با این حال، استفاده از گرادیان خام با چالشی اساسی به‌نام اشباع گرادیان\enfootnote{gradient saturation} روبه‌روست. در شبکه‌های عمیق، اگر یک ویژگی ورودی بسیار مهم و قوی باشد، ممکن است خروجی توابع فعال‌سازی (نظیر سیگموئید، \lr{ReLU} یا لایه‌های بیشینه) به ناحیه‌ی افقی و مسطح خود برسد؛ در این حالت، مشتق محلی صفر می‌شود و مدل به اشتباه اهمیت آن ویژگی کلیدی را صفر ارزیابی می‌کند، در حالی که دقیقاً همان ویژگی عاملِ اصلی رسیدن نورون به حالت اشباع بوده است~\cite{shrikumar17learning}.

برای غلبه بر مشکل اشباع گرادیان در پردازش تصاویر، سانداراراجان و همکاران (۲۰۱۷)~\cite{sundararajan2017axiomatic} چارچوب روش گرادیان‌های یکپارچه\enfootnote{integrated gradients} (\lr{IG}) را معرفی کردند. آن‌ها استدلال کردند که یک روش انتساب معتبر باید اصل کامل‌بودن\enfootnote{completeness} را ارضا کند؛ یعنی مجموع امتیازات انتساب‌یافته به تمام ویژگی‌ها، باید دقیقاً برابر با اختلاف خروجی مدل برای ورودی واقعی ($x$) و یک ورودی پایه یا خنثی ($x'$) باشد: $\sum_i \text{\lr{IG}}_i(x) = F(x) - F(x')$. متدولوژی \lr{IG} این اصل را با محاسبه‌ی انتگرال مسیر گرادیان‌ها در طول یک خط راست از ورودی پایه (مثلاً یک تصویر کاملاً سیاه) تا ورودی واقعی محقق می‌سازد. از نظر ریاضی، اهمیت ویژگی $i$ام ($\text{\lr{IG}}_i$) به‌صورت زیر محاسبه می‌شود:
\begin{equation}\label{eq:chap3_integrated_gradients}
    \text{\lr{IG}}_i(x)=(x_i-x'_i)\int_{\alpha=0}^{1}\frac{\partial F(x'+\alpha(x-x'))}{\partial x_i}d\alpha
\end{equation}
که در آن $F$ تابع پیش‌بینی مدل و $\alpha \in [0, 1]$ ضریب درون‌یابی پیوسته در امتداد مسیر است. مفهوم شهودی انتگرال‌گیری در این فرمول آن است که ویژگی $i$ام را از حالت خنثی ($x'_i$) به‌صورت تدریجی تا مقدار واقعی‌اش ($x_i$) روشن می‌کنیم و در تمام نقاط میانی این گذار، حساسیت مدل ($\frac{\partial F}{\partial x_i}$) را ثبت و با هم جمع می‌بندیم. بدین‌ترتیب، حتی اگر در انتهای مسیر تابع اشباع شده و شیب محلی به صفر رسیده باشد، مشارکت آن ویژگی در طول مسیر ثبت شده و سهم کامل آن، مقیاس‌شده با کل جابه‌جایی ورودی ($x_i - x'_i$)، به‌دست می‌آید.

به‌موازات این روش، شریکومار و همکاران (۲۰۱۷)~\cite{shrikumar17learning} الگوریتم «انتشار تفاوت فعال‌سازی‌ها»\enfootnote{propagating activation differences} یا \lr{DeepLIFT} را برای مدل‌های بینایی پیشنهاد دادند. این روش به‌جای محاسبه‌ی انتگرال‌های پرهزینه، مستقیماً تفاوت میان فعال‌سازی هر نورون برای ورودی واقعی و ورودی پایه ($\Delta y = f(x) - f(x')$) را محاسبه می‌کند. به‌طور دقیق‌تر، با استفاده از قانون زنجیره‌ایِ گسسته در طی یک بار عبور پس‌رو، سهم‌های مجزایی ($C_{\Delta x_i \Delta y}$) را به هر ویژگی ورودی انتساب می‌دهد تا شرط کامل‌بودن ($\sum_i C_{\Delta x_i \Delta y} = \Delta y$) ارضا شود. این روش ضمن حل مشکل اشباع، توانایی تشخیص اثرات مثبت و منفی غیرخطی را نیز دارد.

با این وجود، استفاده از این روش‌های مبتنی بر ورودی مبنا (\lr{IG} و \lr{DeepLIFT}) در مواجهه با متون و زبان طبیعی با یک ابهام معنایی و نظری گره خورده است: انتخاب مبنای خنثی. در حالی که در تصویر می‌توان یک پیکسل سیاه را به‌عنوان پایه در نظر گرفت، تعریف یک بردار تعبیه‌ی صفر یا خنثی در فضای گسسته‌ی متن، فاقد معنای روشن زبانی است. ادغام یک مسیر پیوسته بین تعبیه‌های گسسته نیز از نظر معنایی دچار ابهام است و انتخاب مبنای نادرست می‌تواند به انتساب‌های کاملاً گمراه‌کننده منجر شود.

علاوه بر ابهام مبنا، تعمیم رویکردهای مبتنی بر مبنا (به‌ویژه منطق انتگرال‌گیری \lr{IG}) از مدل‌های متمایزگر به مدل‌های مولد خودهمبسته، با یک چالش نظری بنیادین دیگر روبه‌روست. در وظایف دسته‌بندی، مدل برای تفکیک کلاس‌ها آموزش می‌بیند و گرادیانِ امتیاز یک کلاس، عوامل متمایزکننده‌ی آن را به‌روشنی نشان می‌دهد. اما مدل‌های زبانی خودهمبسته منحصراً با داده‌های مثبت آموزش می‌بینند تا احتمال تولید توالی‌های واقعی را بیشینه کنند؛ در نتیجه، فاقد کلاس منفی یا مرز تصمیم مشخص برای لنگر انداختن محاسبات انتساب هستند. 

برای غلبه بر این دو معضل، فرانک و همکاران (۲۰۲۵)~\cite{frank2025attribution} چارچوب \lr{GAMA} را توسعه دادند. نوآوری اصلی این روش در آن است که برای استخراج ویژگی‌های معنادار در غیاب داده‌های منفی، به مقایسه‌ی رفتار مدل آموزش‌دیده با یک مدل مرجع تصادفی متوسل می‌شود. در این چارچوب، ابتدا الگوریتم \lr{IG} در وظیفه‌ی پیش‌بینی توکن بعدی، برای مدل اصلی محاسبه می‌شود. سپس، دقیقاً همین محاسبات انتگرال‌گیری برای یک مدل با معماری کاملاً مشابه اما با وزن‌های تصادفی و آموزش‌ندیده تکرار می‌گردد. در نهایت، \lr{GAMA} با اندازه‌گیری تفاوت میان امتیازات انتساب‌یافته توسط مدل اصلی و مدل تصادفی، نویزهای ساختاری شبکه را خنثی کرده و منحصراً الگوها و وابستگی‌هایی را برجسته می‌سازد که مدل طی فرایند آموزش از توزیع داده‌های مثبت یاد گرفته است.

اگرچه \lr{GAMA} مشکل فقدان کلاس منفی را در مدل‌های مولد حل می‌کند، اما پیاده‌سازی آن در مدل‌های خودهمبسته نیازمند محاسبات پیوسته و سنگین گرادیان در هر گام از تولید است که بار محاسباتی شدیدی را تحمیل می‌کند. افزون بر این، \lr{GAMA} همچنان متکی بر ادغام مسیر در فضای گسسته‌ی تعبیه‌هاست و نتوانسته چالش نظری عدم‌تطابق فضای پیوسته و گسسته در زبان را به‌طور کامل رفع کند.

چالش دیگر روش‌های مبتنی بر مشتق، پرنویز بودن نقشه‌های اهمیت است. توابع فعال‌سازی در شبکه‌های عمیق دارای نوسانات محلی شدیدی هستند که باعث می‌شود گرادیان خام برای انسان غیرتفسیرپذیر باشد. اسمیلکوف و همکاران (۲۰۱۷)~\cite{smilkov2017smoothgrad} با ارائه تکنیک \lr{SmoothGrad} نشان دادند که می‌توان سیگنال گرادیان را با تجمیع هموار کرد. در این متدولوژی، با افزودن مقادیر جزئی از نویز تصادفی گاوسی ($\mathcal{N}(0, \sigma^2)$) به ورودی اصلی، ده‌ها نسخه‌ی نویزدار تولید می‌شود. سپس گرادیان نقشه‌ی برجستگی ($M_c$) برای هر نسخه به‌صورت جداگانه محاسبه شده و میانگین آن‌ها به‌عنوان نقشه‌ی نهایی استخراج می‌گردد:
\begin{equation}\label{eq:chap3_smoothgrad}
    \hat{M}_c(x)=\frac{1}{n}\sum_{i=1}^n M_c(x+\mathcal{N}(0,\sigma^2))
\end{equation}
شهود فیزیکی این روش مشابه فیلترهای پایین‌گذر در پردازش سیگنال است؛ تپه‌های تیز و نوسانات با فرکانس بالای چشم‌انداز زیان در اطراف نقطه $x$، بر اثر جابه‌جایی‌های تصادفی در جهات مختلف همدیگر را خنثی می‌کنند (تداخل ویرانگر)، در حالی که جهت شیب کلی و پایدار که نشان‌دهنده‌ی اهمیت ساختاری ویژگی‌هاست تقویت می‌شود. با این حال، اعمال این روش مستلزم ده‌ها بار عبور پیش‌رو و پس‌رو در شبکه است که هزینه‌ی محاسباتی را به‌شدت بالا می‌برد؛ به‌علاوه اینکه در پردازش متن، افزودن مستقیم نویز پیوسته به فضای گسسته‌ی تعبیه‌ها مجدداً با همان چالش‌های نظری عدم‌تطابق فضای پیوسته و گسسته روبه‌رو می‌شود.

\paragraph{پراکسی‌های مکانیزم توجه.}

با توجه به سربار محاسباتی سنگین روش‌های مبتنی بر پس‌انتشار در شبکه‌های عمیق، پژوهشگران حوزه‌ی پردازش زبان طبیعی به بررسی مستقیم مکانیزم مسیریابی درون معماری ترنسفورمرها، یعنی وزن‌های توجه، روی آوردند. در این رویکرد فرض بر این بود که ماتریس‌های خودتوجهی ($\alpha_{i,j}$)، مستقیماً نشان می‌دهند که مدل در حین تولید یک توکن خاص ($i$)، به کدام توکن‌های پیشین ($j$) نگاه کرده است؛ در نتیجه می‌توان از توزیع توجه به‌عنوان یک پراکسیِ ارزان برای انتساب اعتبار استفاده کرد.

با این حال، این پیش‌فرض شهودی به‌زودی به یکی از بحث‌برانگیزترین مناظرات علمی در ادبیات تفسیرپذیری تبدیل شد. جین و والاس (۲۰۱۹)~\cite{jain2019attention} در مقاله‌ی مشهور «توجه، توضیح نیست»، وفاداریِ وزن‌های توجه به خروجی نهایی مدل را به چالش کشیدند. آزمایش‌های آن‌ها نشان داد که همبستگی میان وزن‌های توجه و معیارهای تثبیت‌شده‌ی مبتنی بر گرادیان بسیار ضعیف است. از آن مهم‌تر، آن‌ها با استفاده از یک روش جستجوی متخاصم، توزیع‌های توجه کاملاً متفاوتی یافتند که با تزریق در شبکه، دقیقاً همان خروجی پیشین را تولید می‌کردند؛ امری که نشان می‌داد توجه بالا به یک کلمه، لزوماً به معنای نقش علّی آن در تصمیم‌گیری نیست.

در پاسخ به این نقد، ویگرف و پینتر (۲۰۱۹)~\cite{wiegreffe2019attention} با مجموعه‌ای از آزمایش‌های کنترلی استدلال کردند که یافتن توزیع‌های متخاصم به‌تنهایی برای رد کارایی توجه کافی نیست. آن‌ها نشان دادند که مدل‌های جایگزینی که از پایه با این توجه‌های متخاصم آموزش دیده بودند، افت عملکرد فاحشی را تجربه می‌کردند. این مناظره‌ی علمی گرچه نشان داد تحت شرایطی می‌توان سیگنال‌های معناداری استخراج کرد، اما اثبات کرد که استخراج استدلال علیِ قطعی از وزن‌های خام ترنسفورمرها، رویکردی ذاتاً مبهم است.

دلیل اصلی این ابهام ساختاری، پدیده‌ای به‌نام اختلاط اطلاعات است. آبنار و زویدما (۲۰۲۰)~\cite{abnar2020quantifying} استدلال کردند که در معماری ترنسفورمر، اطلاعات توکن‌ها در هر لایه با یکدیگر ترکیب می‌شوند؛ بنابراین، بازنماییِ یک توکن در لایه‌های عمیق‌تر، دیگر منحصراً متعلق به همان کلمه‌ی اولیه نیست و اتکا به ماتریس توجه یک لایه‌ی خاص به‌شدت گمراه‌کننده است. آن‌ها برای کمّی‌سازی جریان واقعی اطلاعات از ورودی تا خروجی، روش بسط توجه را پیشنهاد دادند. در این روش، برای در نظر گرفتن اتصالات باقیمانده، ماتریس همانی ($I$) به ماتریس توجه خام هر لایه ($A^{(l)}$) افزوده می‌شود. پس از نرمال‌سازی سطریِ این ترکیب برای حفظ خاصیت توزیع احتمال (ماتریس $\hat{A}^{(l)}$)، این ماتریس‌ها به‌صورت متوالی در یکدیگر ضرب می‌شوند تا ماتریس بسط‌یافته‌ی نهایی ($\tilde{A}$) به‌دست آید:
\begin{equation}\label{eq:chap3_attention_rollout}
    \tilde{A} = \hat{A}^{(L)} \cdot \hat{A}^{(L-1)} \cdots \hat{A}^{(1)}
\end{equation}
شهود ریاضی پشت این فرمول‌بندی ضرب زنجیره‌ای، مدل‌سازیِ گذر اطلاعات در یک گراف چندلایه‌ی جهت‌دار است. جمع ماتریس توجه با ماتریس همانی ($A^{(l)} + I$) به این معناست که بخشی از اطلاعات هر توکن از طریق مسیرهای میان‌بر (اتصالات باقیمانده) بدون هیچ تغییری مستقیماً به لایه‌ی بعد منتقل می‌شود؛ سپس ضرب متوالی ماتریس‌ها معادل محاسبه‌ی مجموع وزن تمام مسیرهای هدایت اطلاعات از توکن‌های لایه‌ی ورودی تا لایه‌ی $L$ام است. آن‌ها همچنین روش جریان توجه را معرفی کردند که شبکه را به‌عنوان یک گراف جهت‌دار مدل‌سازی کرده و از الگوریتم بیشینه-جریان\enfootnote{max-flow} برای یافتن ظرفیت مسیرهای انتقال اطلاعات استفاده می‌کند. برای عملیاتی کردن این مفاهیم پیچیده در مدل‌های نوین تولید متن، چارچوب \lr{Inseq}~\cite{sarti2023inseq} استخراج وزن‌های توجه و محاسبه‌ی بسط توجه را برای هر دو معماریِ رمزگذار-رمزگشا (نظیر \lr{T5}~\cite{raffel2020exploring}) و فقط-رمزگشا (نظیر سری \lr{GPT}~\cite{brown2020language}) از نظر ریاضی استانداردسازی کرد. این ابزار به محققان اجازه داد تا بدون درگیری با پیچیدگی‌های پیاده‌سازی مدل‌های خودهمبسته، سهم توجه اختصاص‌یافته به هر توکن را در هر گام از تولید توالی به‌صورت مجزا و شفاف ردیابی کنند.

با این وجود، روش‌های مبتنی بر بسط یا جریان توجه همچنان از یک نقص نظری رنج می‌بردند؛ آن‌ها صرفاً تأثیر مکانیزم توجه را مدل می‌کردند و تأثیر بردارهای ارزش و شبکه‌های پیش‌خور (که مخازن اصلی دانش پارامتریک مدل هستند) را نادیده می‌گرفتند. برای رفع این نقیصه‌ی بنیادین، چفر و همکاران (۲۰۲۱)~\cite{chefer2021transformer} الگوریتم «انتشار لایه‌به‌لایه‌ی ارتباط»\enfootnote{layer-wise relevance propagation} یا \lr{Transformer-LRP} را بر پایه‌ی اصل تجزیه‌ی تیلور عمیق\enfootnote{deep Taylor decomposition} ارائه کردند. این چارچوب با معرفی قواعد مشتق‌گیریِ اختصاصی برای مکانیزم‌های توجه چندسره، تضمین می‌کند که قانون بقای ارتباط حفظ شود؛ یعنی مجموع ارتباط توکن‌های ورودی ($R_i$) باید دقیقاً برابر با خروجی نهایی مدل ($f(x)$) باشد:
\begin{equation}\label{eq:chap3_conservation_of_relevance}
    \sum_i R_i = f(x)
\end{equation}
شهود حاکم بر این اصل بقا، مشابه قوانین بقای شار در فیزیک است: تصمیم نهایی شبکه ($f(x)$) مانند مقداری مشخص از سیال یا انرژی فرض می‌شود که باید در مسیر پس‌رو، بدون هیچ‌گونه اتلاف یا تولید ناگهانی، در میان گره‌های گراف محاسباتی تقسیم گردد ($\sum_j R_j^{(l+1)} = \sum_i R_i^{(l)}$). بدین‌ترتیب، ارتباط توکن‌های ورودی سهم دقیق و بدون نشتِ هر کلمه را از مقدار کل پیش‌بینی بازتاب می‌دهد.

در تکاملی دیگر برای بررسی دقیق‌تر تعاملات میان توکن‌ها، هائو و همکاران (۲۰۲۱)~\cite{hao2021selfattention} الگوریتم «انتساب خودتوجهی»\enfootnote{self-attention attribution} را توسعه دادند. این پژوهش نشان داد که ترکیب خطیِ ساده‌ی گرادیان‌ها و وزن‌های توجه (مانند ارزیابی بر اساس $\alpha_{i,j} \times \frac{\partial y}{\partial \alpha_{i,j}}$) به‌دلیل بروز مشکلاتی نظیر اشباع گرادیان، نمی‌تواند وابستگی‌های غیرخطی عمیق را به‌درستی مدل‌سازی کند. از این رو، آن‌ها منطق الگوریتم \lr{IG} را مستقیماً بر روی ماتریس‌های توجه پیاده‌سازی کردند. در این روش، لبه‌های ارتباطی ماتریس توجه به‌عنوان ویژگی‌های ورودی در نظر گرفته شده و امتیاز انتساب هر پیوند از طریق زیر به‌دست می‌آید:
\begin{equation}\label{eq:chap3_selfattention_attribution}
    \text{\lr{Attr}}(\alpha_{i,j}) = \alpha_{i,j} \times \int_{\beta=0}^{1} \frac{\partial F(\beta A)}{\partial \alpha_{i,j}} d\beta
\end{equation}
شهود این فرمول آن است که ماتریس کامل توجه به‌عنوان ورودی مدل تلقی شده و از یک مبنای خنثی ($\beta=0$، یعنی زمانی که تمام توکن‌ها به‌طور کامل از هم ایزوله هستند و هیچ توجهی ردوبدل نمی‌شود) تا وضعیت کامل و عملیاتی شبکه ($\beta=1$) مقیاس‌دهی می‌شود. انتگرال‌گیری پیوسته بر روی ضریب $\beta$، حساسیت مدل به اتصال میان توکن $i$ و $j$ را در سراسر این بازه‌ی شکل‌گیری ارتباط محاسبه می‌کند. ضرب این حساسیتِ انباشته در وزن خام توجه ($\alpha_{i,j}$)، تضمین می‌کند که اثرات غیرخطی لحاظ شده و لبه‌های کاذب (که وزن بالایی دارند اما تأثیری بر خروجی ندارند) حذف شوند. در نهایت، با استفاده از این امتیازات دقیق ریاضی، یک درخت انتساب استخراج می‌گردد که ساختار سلسله‌مراتبی و نحوی تعاملات اطلاعاتی درون ترنسفورمر را به‌طور شفاف نمایان می‌سازد.

\paragraph{تبیین و انتساب تقابلی در سطح توکن.}

با وجود پیشرفت‌های حاصل در انتساب گرادیان و توجه، پژوهشگران با الهام از علوم شناختی متوجه شدند که توضیحات انسانی در بستر زبان ذاتاً ماهیتی تقابلی دارند. انسان‌ها برای درک یک رویداد، معمولاً نمی‌پرسند چرا یک رویداد خاص رخ داد، بلکه می‌پرسند چرا آن رویداد رخ داد در حالی که انتظار می‌رفت پادنمونه‌ی آن رخ دهد. جاکووی و همکاران (۲۰۲۱)~\cite{jacovi2021contrastive} استدلال کردند که برای پیاده‌سازی این منطق در شبکه‌های زبانی، در یک لایه‌ی دسته‌بند خطی با بردارهای وزن $w_A$ و $w_B$، مرز تصمیم میان این دو کلاس معادل بردار تفاضل آن‌ها ($w_A - w_B$) است. با تصویر کردن بازنمایی پنهان متن بر روی این محور تقابلی، می‌توان مشخص کرد کدام توکن‌ها مستقیماً شبکه را به سمت کلاس هدف سوق داده و هم‌زمان از کلاس پادنمونه\enfootnote{foil class} دور کرده‌اند.

تلاش برای بسط این منطق تقابلی به مدل‌های زبانی مولد، ظرایف ساختاری جدیدی را طلب می‌کرد. در وظیفه‌ی پیش‌بینی توکن بعدی، خروجی مدل حاصل رقابت ده‌ها هزار واژه است. یین و نوبیگ (۲۰۲۲)~\cite{yin2022interpreting} برای غلبه بر این پیچیدگی، تابع هدف روش‌های مرسوم انتساب اعتبار را بازنویسی کردند. در این متدولوژی مشتقات به‌جای آنکه نسبت به لوجیت یک تک‌توکن هدف محاسبه شوند، مستقیماً نسبت به اختلاف لوجیت توکن هدف ($z_t$) و توکن پادنمونه ($z_f$) استخراج می‌گردند:
\begin{equation}\label{eq:chap3_contrastive_attribution}
    c_i = \left( \nabla_{x_i} z_t - \nabla_{x_i} z_f \right) \cdot x_i
\end{equation}
منطق شهودی پشت تفاضل گرادیان‌ها ($\nabla_{x_i} z_t - \nabla_{x_i} z_f$) آن است که در تولید متن، بسیاری از توکن‌های زمینه (مانند کلمات رایج دستوری یا موضوع عمومی متن)، لوجیت تمامی واژگان مرتبط با آن حوزه را به‌طور هم‌زمان بالا می‌برند. در انتساب معمولی، این کلمات نویز زیادی ایجاد می‌کنند؛ اما در فرمول‌بندی تقابلی، اثری که به یک اندازه بر هر دو توکن هدف و پادنمونه تأثیر بگذارد خنثی می‌شود ($\nabla_{x_i} z_t \approx \nabla_{x_i} z_f$). بنابراین، امتیاز $c_i$ تنها کلماتی را برجسته می‌سازد که به‌طور مشخص ترازو را به‌نفع توکن هدف سنگین کرده و توکن رقیب را سرکوب نموده‌اند. برای ارزیابی این رویکرد، نویسندگان از جفت‌های زبانی با ساختارهای حداقلی و تقابلی بهره بردند؛ همچنین در ارزیابی استدلال‌های خلاف‌واقع و تقابلی در مدل‌های زبانی، محک‌هایی نظیر مجموعه‌داده‌ی \lr{CRASS}~\cite{frohberg2022crass} برای تدوین سناریوهای پادنمونه و سنجش تصمیمات ظریف زبانی به‌کار گرفته می‌شوند.

برای کاربردی کردن این استراتژی در مدل‌های مقیاس‌بزرگ، سارتی و همکاران (۲۰۲۳)~\cite{sarti2023inseq} در کتابخانه‌ی \lr{Inseq}، «انتساب ویژگی تقابلی»\enfootnote{contrastive feature attribution} را به‌عنوان یکی از ارکان اصلی چارچوب خود پیاده‌سازی کردند. در این سیستم، کاربر می‌تواند یک توکن هدف ($y$) و یک توکن پادنمونه ($y'$) را مشخص کند و ابزار با محاسبه‌ی تفاضل لوجیت‌ها ($\Delta \text{\lr{logits}} = \text{\lr{logit}}(y) - \text{\lr{logit}}(y')$)، نقشه‌های حرارتی دقیقی تولید می‌کند که نشان می‌دهد کدام بخش‌های پرامپت منحصراً سبب برتری توکن منتخب بر رقیب شده‌اند. با این حال، اعمال این استراتژی در تولید متن آزاد و طولانی همچنان با دشواری‌های جدی مواجه است، زیرا تعریف یک پادنمونه یگانه و صریح برای کل یک پاراگراف یا یک فرایند استدلال چندمرحله‌ای همواره امکان‌پذیر نیست.

علی‌رغم دستاورد بزرگ ابزارهایی نظیر \lr{Inseq} در استانداردسازی و تعمیم ریاضی روش‌های انتساب (از گرادیان و توجه تا استدلال تقابلی) برای مدل‌های مولد، این پیشرفت‌های مهندسی محدودیت‌های بنیادین پارادایم مبتنی بر توکن را عریان‌تر کرد. حتی با محاسبه‌ی دقیق‌ترین گرادیان‌های تقابلی در گام‌های رمزگشایی، خروجی نهایی تمامی این روش‌ها همچنان یک نقشه‌ی حرارتی پرنویز در سطح کلمات خُرد است. فقدان یک ساختار گزاره‌ای و منطقی در این خروجی‌ها و ناتوانی آن‌ها در تجمیع این سیگنال‌های خلوت به مفاهیمی که برای ذهن انسان درک‌پذیر باشند، لزوم گذار از تحلیل سیگنال‌های پایه به رویکردهای مبتنی بر استخراج مفهوم را بیش از پیش توجیه می‌نماید.


\subsubsection{تحلیل و استخراج مفاهیم در فضای پنهان}\label{sec:3-1-2-2}

با وجود دقت روش‌های انتساب بر پایه‌ی سیگنال‌های اولیه، این رویکردها با یک محدودیت معنایی بنیادین روبه‌رو هستند. خروجی آن‌ها در پایین‌ترین سطح وضوح، یعنی در سطح پیکسل‌های تصویر یا توکن‌های منفرد متن ارائه می‌شود. درک اینکه کدام توکن‌ها مهم هستند، لزوماً به درک منطق و استدلال شبکه منجر نمی‌شود، زیرا مفاهیم شناختی انسان (مانند یک ویژگی فیزیکی، یک لحن خاص، یا یک پیش‌فرض منطقی) ماهیتی توزیع‌شده دارند و در یک نشانه یا توکن واحد خلاصه نمی‌شوند.

پیش‌تر در بخش~\ref{sec:3-1-1-3} در تفسیرپذیری درونی دیدیم که چگونه مدل‌های مفهومی تلاش می‌کردند معماری شبکه را از پایه برای استدلال بر روی مفاهیم انسانی بهینه‌سازی کنند. در پارادایم پساعملی اما هدف کاملاً متفاوت است. در اینجا با یک مدل پیش‌آموزش‌دیده و منجمد روبه‌رو هستیم و می‌خواهیم فضای پنهان آن را برای کشف مفاهیم کاوش کنیم. پارادایم تفسیرپذیری مبتنی بر مفهوم در این بخش، بر این فرضیه استوار است که می‌توان بردارهای متراکم داخلی شبکه را بدون تغییر در معماری اصلی، به مفاهیم شناختی انسان پیوند داد و نشان داد خروجی مدل تا چه حد به حضور یک مفهوم مشخص در گراف محاسباتی‌اش وابسته است.

پژوهش‌ها در این حوزه را می‌توان در قالب یک سیر منطقی ساختاردهی کرد. ادبیات این بخش ابتدا از منظر استخراج نظارت‌شده و نگاشت مستقیم مفاهیم بررسی می‌شود. در پاسخ به چالش‌های مقیاس‌پذیری این متدولوژی‌ها، رهیافت‌های مبتنی بر کشف خودکار مفاهیم و تجزیه‌ی فضای پنهان مورد واکاوی قرار می‌گیرند. سپس برای پر کردن شکاف معنایی ناشی از رویکردهای صرفاً ریاضی، مکانیزم‌های ارزیابی و اثبات علیت معرفی می‌شوند. در نهایت برای پاسخ به چالش‌های بنیادین زبان طبیعی، ضرورت کاوش در هندسه‌ی درونی و دینامیک نظری مختص مدل‌های زبانی تبیین می‌گردد.

\paragraph{استخراج نظارت‌شده و نگاشت مستقیم مفاهیم.}

نخستین و شهودی‌ترین گام برای استخراج پساعملی مفاهیم، متکی بر تعریف و برچسب‌گذاری صریح آن‌ها توسط انسان بود. ایده‌ی محوری در این رویکرد آن است که پس از تعریف یک مفهوم، جهت خطی متناظر با آن در فضای پنهان مدل جستجو شود. ریشه‌ی نظریِ این ایده به پژوهش بنیادین آلین و بنجیو (۲۰۱۷)~\cite{alain2017understanding} در معرفی کاوشگرهای خطی بازمی‌گردد. این پژوهشگران برای درک اینکه لایه‌های میانی یک شبکه‌ی عصبی عمیق (به‌ویژه شبکه‌های متمایزگر در بینایی ماشین) چه اطلاعاتی را در خود ذخیره کرده‌اند، پیشنهاد دادند که یک دسته‌بند خطیِ ساده به بازنمایی‌های هر لایه متصل شود.

از نظر ریاضی، با فرض اینکه $h_k(x)$ بردار فعال‌سازی لایه‌ی $k$ام برای ورودی $x$ باشد، یک کاوشگر با فرمول $f_k(h_k(x)) = W_k h_k(x) + b_k$ آموزش می‌بیند تا برچسب هدف یا ویژگی خاصی را پیش‌بینی کند، در حالی که وزن‌های مدل اصلی کاملاً منجمد و ثابت باقی می‌مانند. یافته‌ی کلیدی آن‌ها این بود که تفکیک‌پذیری خطی\enfootnote{linear separability} با حرکت به سمت لایه‌های عمیق‌تر، به‌صورت یکنواخت افزایش می‌یابد؛ امری که ثابت می‌کرد شبکه‌های عصبی به‌تدریج اطلاعات درهم‌تنیده را باز کرده و مفاهیم انتزاعی را در قالب خطوط راست در فضاهای پنهان خود سازماندهی می‌کنند.

با اتکا به این پشتوانه‌ی نظری برای جستجوی بردارهای خطی، کیم و همکاران (۲۰۱۸)~\cite{kim2018interpretability} با تمرکز بر شبکه‌های بینایی ماشین، چارچوب ارزیابی کمّی با بردارهای فعال‌سازی مفهوم\enfootnote{concept activation vectors} (\lr{TCAV}) را به‌عنوان بنیان این تفکر معرفی کردند. مسئله‌ی محوری پژوهش آن‌ها رفع رسا نبودن نقشه‌های برجستگی بود. برای مثال، نقشه‌های سنتی تنها نشان می‌دادند کدام پیکسل‌ها برای تشخیص کلاس یک حیوان مهم بوده‌اند، اما مشخص نمی‌کردند که آیا مدل به مفهوم انتزاعیِ «راه‌راه بودن» نیز در تصمیم‌گیری خود توجه کرده است یا خیر.

متدولوژی \lr{TCAV} برای کشف این وابستگی از یک فرایند نظارت‌شده‌ی مستقیم بهره می‌برد. کاربر مجموعه‌ای از تصاویر نمایانگر مفهوم هدف $C$ و مجموعه‌ای از تصاویر تصادفی را به‌عنوان داده‌های منفی فراهم می‌کند. این داده‌ها از مدل مورد نظر عبور داده می‌شوند تا فعال‌سازی‌های مربوط به یک لایه‌ی پنهان مشخص ($l$) استخراج گردند. سپس، یک دسته‌بند خطی نظیر ماشین بردار پشتیبان برای تفکیک فعال‌سازی‌های مثبت و منفی آموزش می‌بیند. بردار عمود بر ابرصفحه‌ی این دسته‌بند، همان بردار \lr{CAV} (با نماد $v_C^l$) است که جهت هندسی حضور مفهوم را در فضای فعال‌سازی نشان می‌دهد. در نهایت، برای سنجش اینکه پیش‌بینی مدل برای یک کلاس خاص ($k$) تا چه حد به این مفهوم حساس است، از مشتق جهت‌دار استفاده می‌شود. اگر $h_k(f_l(x))$ تابع خروجی مدل (لوجیت کلاس $k$) بر اساس فعال‌سازی‌های لایه $l$ در نقطه $x$ باشد، حساسیت شبکه به مفهوم $C$ به‌صورت زیر محاسبه می‌گردد:
\begin{equation}\label{eq:chap3_tcav}
    S_{C,k,x}^l = \nabla h_k(f_l(x)) \cdot v_C^l
\end{equation}
از منظر شهودی، گرادیان $\nabla h_k(f_l(x))$ جهت بیشترین افزایش امتیاز کلاس $k$ را در فضای پنهان نشان می‌دهد و ضرب داخلی آن با بردار یکه‌ی مفهوم ($v_C^l$)، تصویر این حساسیت بر جهت مفهوم است. این مشتق جهت‌دار به این پرسش پاسخ می‌دهد که «اگر بازنمایی درونی ورودی $x$ را در لایه‌ی $l$ به اندازه‌ی یک گام کوچک به سمت مفهوم $C$ حرکت دهیم، اطمینان مدل به کلاس $k$ افزایش می‌یابد ($S > 0$) یا افت می‌کند ($S < 0$)؟». برای تجمیع این اثر موضعی در کل مجموعه‌داده‌ی کلاس $k$ ($X_k$)، شاخص کلی \lr{TCAV} به‌صورت نسبت نمونه‌هایی تعریف می‌شود که دارای مشتق جهت‌دار مثبت هستند (یعنی $\text{\lr{TCAV}}_{k,C}^l = \frac{|\{x \in X_k : S_{C,k,x}^l > 0\}|}{|X_k|}$). این شاخص درصدی (عددی بین ۰ و ۱)، تصویری کلان و کمّی از میزان اتکای مدل به یک مفهوم انسانی (مثلاً میزان اثرگذاری مفهوم «راه‌راه بودن» بر تشخیص «گورخر») در سراسر توزیع داده‌ها ارائه می‌دهد.

منطق استفاده از مرزهای تصمیم خطی به‌سرعت از مدل‌های متمایزگر به مدل‌های مولد نیز تعمیم یافت؛ با این تفاوت که در مدل‌های مولد، هدف صرفاً سنجش حساسیت نبود، بلکه کنترل خروجی مدنظر قرار داشت. شن و همکاران (۲۰۲۰)~\cite{shen2020interpreting} در چارچوب \lr{InterFaceGAN} از آموزش دسته‌بندهای خطی بر روی بردارهای پنهان برچسب‌دار استفاده کردند تا جهت‌های معنایی کلان نظیر سن یا جنسیت را در فضای پنهان استخراج کنند. آن‌ها نشان دادند که در مدل‌های مولد، حرکت دادن بردار نویز پایه در امتداد این بردارهای نرمال، امکان کنترل مستقیم و تفسیرپذیرِ ویژگی‌های تصویر تولیدی را فراهم می‌آورد.

در قلمروی پردازش زبان طبیعی نیز، استفاده از کاوشگرهای خطی برای کشف مفاهیم عمیق‌تر و انتزاعی‌تر به‌کار گرفته شد. آزاریا و میچل (۲۰۲۳)~\cite{azaria2023theinternal} در پژوهشی جالب‌توجه نشان دادند که می‌توان مفهوم «حقیقت» را مستقیماً از فضای پنهان مدل‌های زبانی استخراج کرد. آن‌ها با تغذیه‌ی گزاره‌های درست و غلط به مدل‌هایی نظیر \lr{LLaMA}~\cite{touvron2023llama} و استخراج فعال‌سازی‌های لایه‌های میانی، یک دسته‌بند خطی را به‌عنوان \lr{Truth Probe} آموزش دادند. یافته‌ی کلیدی آن‌ها این بود که کاوشگر خطی می‌تواند با دقت بالایی صرفاً با نگاه به حالت‌های داخلی مدل، دروغ یا راست بودن گزاره را تشخیص دهد؛ حتی زمانی که خود مدل در خروجی متنی‌اش با لحنی مطمئن در حال توهم‌زنی و بیان یک گزاره‌ی نادرست بود. این دستاورد، علاوه بر اثبات حضور مفاهیم خطی پیچیده در شبکه، پرده از یک شکاف بنیادین میان دانش نهفته در بازنمایی‌های داخلی و خروجی متنی نهایی برداشت.

علی‌رغم قدرت بالای این روش‌ها در ایجاد پیوند میان هندسه‌ی شبکه و ادراک انسانی، ارزیابی آن‌ها در بستر زبان طبیعی محدودیت‌های بازدارنده‌ای را آشکار می‌سازد. نخست، نیاز به داده‌های نظارتی و برچسب‌دار برای هر مفهوم جدید، مقیاس‌پذیری این روش را به‌خصوص در تولید متن آزاد به‌شدت کاهش می‌دهد؛ چرا که زبان شامل بی‌نهایت مفهوم استنتاجی، نحوی و معنایی است که جمع‌آوری داده‌ی مجزا برای استخراج هر کدام عملاً ناممکن است. دوم، این چارچوب‌ها بر یک پیش‌فرض قوی استوارند که مفاهیم همواره به‌صورت خطی در لایه‌های پنهان قابل‌تفکیک هستند؛ فرضی که با توجه به درهم‌تنیدگیِ بازنمایی‌ها در شبکه‌های زبانی پیچیده، همواره صادق نیست. افزون بر این چالش‌های عملی، اتکا به این جستجوی خطی از منظر نظری نیز با انتقادات بنیادینی روبه‌رو شده است. رِج (۲۰۲۴)~\cite{rege2024probe} با بررسی پارادایم کاوشگرها نشان داد که در فضاهای برداری با ابعاد بسیار بالا – نظیر \lr{LLM}ها – یک دسته‌بند خطی به‌دلیل محدودیت‌های پیچیدگی نمونه به‌شدت مستعد فریب خوردن است؛ به این معنا که می‌تواند همبستگی‌های کاملاً تصادفی و نویزی را به‌عنوان یک مسیر مفهومی تشخیص دهد. این نقد اثبات می‌کند که یافتن یک بردار خطی تفکیک‌کننده برای یک مفهوم، لزوماً به معنای استفاده‌ی علّی مدل از آن مفهوم نیست؛ این چالش نظری در انتهای این بخش به‌تفصیل واکاوی خواهد شد.

در این راستا، جریان جدیدی در فضای پساعملی شکل گرفت که متمرکز بر رفع معایب مدل‌های گلوگاه مفهومی (\lr{CBM}) بود. همان‌طور که پیش‌تر در بخش~\ref{sec:3-1-1-3} بررسی شد، \lr{CBM}های کلاسیک نیازمند آموزش مجدد کل شبکه از پایه بودند. از مهم‌ترین قدم‌ها برای عبور از این محدودیت، چارچوب یوکسک‌گونول و همکاران (۲۰۲۲)~\cite{yuksekgonul2022posthoc} بود که مدل‌های گلوگاه مفهومی پساعملی\enfootnote{post-hoc concept bottleneck models} (\lr{P-CBM}) را معرفی کردند. این روش، یک مدل پیش‌آموزش‌دیده را بدون نیاز به آموزش کل شبکه، به یک مدل مفهومی تبدیل می‌کند. در این چارچوب، مفاهیم از مجموعه‌ای از واژگان متنی ازپیش‌تعریف‌شده توسط انسان (مانند لیست ویژگی‌ها در یک مجموعه‌داده) فراهم می‌شود. سپس با استفاده از مدل‌های چندوجهی پایه نظیر رمزگذار متنی \lr{CLIP}~\cite{radford2021learning}، این واژگان به یک بانک غنی از تعبیه‌های مفهومی تبدیل می‌گردند. در گام بعد، بازنمایی‌های داخلی لایه‌های ماقبل آخر مدل هدف، بر روی این بردارهای مفهومی تصویر می‌شوند تا میزان فعال‌بودن هر مفهوم محاسبه گردد. در نهایت، لایه‌ی تصمیم‌گیرنده‌ی مدل اصلی دور انداخته شده و یک دسته‌بند خطی خلوت مستقیماً بر روی این امتیازات مفهومی برای انجام پیش‌بینی آموزش می‌بیند.

با این حال، برای جبران افت عملکردی که در اثر اعمال این گلوگاه پساعملی ایجاد می‌شد، نویسندگان یک نسخه‌ی ترکیبی یا هیبریدی (\lr{HP-CBM}) را ارائه دادند که در آن، یک اتصال باقیمانده مستقیماً از ویژگی‌های پنهان مدل اولیه به دسته‌بند نهایی کشیده می‌شد. این ترفند معماری، همان‌طور که پیش‌تر در بحث رویکردهای درونی و شبکه‌های \lr{CEM} نیز شاهد آن بودیم، به چالش نشت اطلاعات می‌انجامد. مدل یاد می‌گیرد اطلاعات حیاتی اما غیرتفسیرپذیر خود را از مسیر باقیمانده عبور دهد و بدین‌ترتیب، تصمیم نهایی دیگر منحصراً وابسته به مفاهیم شفاف درون گلوگاه نیست.

در رویکردی مشابه برای پیشبرد مدل‌های پساعملی و حذف کامل انسان از چرخه‌ی استخراج و برچسب‌گذاری، اویکارینن و همکاران (۲۰۲۳)~\cite{oikarinen2023labelfree} چارچوب مدل‌های گلوگاه مفهومی بدون‌برچسب\enfootnote{label-free concept bottleneck models} (\lr{Label-free CBM}) را معرفی کردند. این روش، فرایند تأمین لیست مفاهیم اولیه را که در \lr{P-CBM} وابسته به انسان و مجموعه‌داده بود، کاملاً خودکار می‌کند. در این چارچوب که همچنان برای وظایف دسته‌بندی تصویر طراحی شده است، ابتدا از یک \lr{LLM} خواسته می‌شود تا ده‌ها ویژگی و مفهوم متنی مرتبط با یک کلاس را تولید کند. سپس، مشابه روش قبل، این مفاهیم متنی تولیدشده توسط \lr{CLIP} به فضای تعبیه‌ها منتقل شده و ماتریس فعال‌سازی مدل هدف بر روی آن‌ها تصویر می‌گردد تا بدون نیاز به دخالت کاربر، امتیازات مفهومی برای دسته‌بند نهایی فراهم شود.

با این وجود، بررسی انتقادی هر دو روش \lr{P-CBM} و \lr{Label-free CBM} یک نقص ساختاری کشنده را در این پارادایم آشکار می‌سازد و آن وابستگی مطلق به مدل‌های خارجی است. مدل‌های \lr{P-CBM} برای تولید تعبیه‌های دقیق به مدل‌هایی نظیر \lr{CLIP} وابسته‌اند و روش‌های بدون‌برچسب اساساً اعتبار مفاهیم خود را از مدل‌های زبانی خارجی می‌گیرند که این امر خطر ورود توهمات، سوگیری‌ها و مفاهیم نامعتبر به گراف استدلال مدل هدف را به‌همراه دارد. دومین نقد، پیرامون ادعای پساعملی بودن آن‌هاست. اگرچه این روش‌ها کل شبکه را از پایه آموزش نمی‌دهند، اما با دور انداختن لایه‌ی نهایی و آموزش یک دسته‌بند خطی جدید بر روی گلوگاه، عملاً بخشی از مدل را دستکاری و بازآموزی می‌کنند.

\paragraph{کشف خودکار مفاهیم و تجزیه فضای پنهان.}

وابستگی مطلق به مدل‌های خارجی و حاشیه‌نویسی انسانی در روش‌های نظارت‌شده، پژوهشگران را به سمت توسعه‌ی رویکردهای بدون‌نظارت سوق داد. ایده‌ی محوری در این جریان تحقیقاتی آن است که برای جلوگیری از تحمیل سوگیری‌های خارجی و درک آنچه مدل واقعاً یاد گرفته است، مفاهیم باید به‌صورت ارگانیک و مستقیماً از توزیع داده‌ها در فضای پنهان استخراج شوند.

یکی از نخستین تلاش‌ها در این مسیر، الگوریتم «توضیحات خودکار مبتنی بر مفهوم»\enfootnote{automatic concept-based explanations} (\lr{ACE}) بود که توسط قربانی و همکاران (۲۰۱۹)~\cite{ghorbani2019towards} معرفی شد. متدولوژی \lr{ACE} از یک خط‌لوله سه‌مرحله‌ای تشکیل می‌شود. در گام نخست، داده‌های ورودی (نظیر تصاویر) به قطعات کوچک‌تر یا همان ابرپیکسل‌ها تقطیع می‌شوند. سپس، این قطعات از شبکه عبور داده شده و ماتریس فعال‌سازی آن‌ها در یک لایه‌ی پنهان استخراج می‌گردد. در نهایت، با اعمال یک الگوریتم خوشه‌بندی بر روی این فعال‌سازی‌ها، هر خوشه به‌عنوان یک مفهومِ کشف‌شده در نظر گرفته شده و اهمیت آن با استفاده از مشتقات جهت‌دار ارزیابی می‌شود.

به‌موازات خوشه‌بندی، رویکردهای متکی بر تحلیل هندسی فضای پنهان در مدل‌های مولدِ تصاویر نیز توسعه یافتند. هارکونن و همکاران (۲۰۲۰)~\cite{harkonen2020ganspace} در چارچوب \lr{GANSpace} نشان دادند که اعمالِ تحلیل مؤلفه‌های اصلی (\lr{PCA}) مستقیماً بر روی توزیع بردارهای فضای پنهان می‌تواند جهت‌های معنایی پنهان را آشکار سازد. در این روش، بردارهای ویژه با بیشترین واریانس به‌عنوان دستگیره‌های کنترلی عمل می‌کنند. به بیان ریاضی، اگر $w$ یک بازنمایی پنهان و $v_k$ یک بردار ویژه‌ی استخراج‌شده باشد، نگاشت جدید به‌صورت $w' = w + \alpha v_k$ محاسبه می‌شود. حرکت پیوسته (با تغییر پارامتر مقیاس‌دهی $\alpha$) در امتداد این بردارهای خطی، تغییرات معنایی (مانند تغییر زاویه‌ی دید یا نورپردازی) در خروجی تصویر ایجاد می‌کند.

با این وجود بردارهای استخراج‌شده صرفاً بر اساس واریانس، لزوماً از نظر معنایی مستقل نیستند. نیو و همکاران (۲۰۲۳)~\cite{niu2023disentangling} برای رفع این مشکل یک مکانیزم جریمه‌سازی را جهت رفع درهم‌تنیدگی این جهت‌های کشف‌شده پیشنهاد دادند تا تغییر یک ویژگی (مانند سن) به‌طور ناخواسته ویژگی دیگری را دستخوش تغییر نکند. در تکاملی مشابه، استاین و همکاران (۲۰۲۴)~\cite{stein2024towards} با معرفی چارچوب «استخراج مفهوم ترکیب‌پذیر»\enfootnote{compositional concept extraction} (\lr{CCE})، از تصویرسازی متعامد و خوشه‌بندی کُروی برای کشف مفاهیم بهره بردند تا تضمین کنند که جمع خطی بردارهای استخراج‌شده، ترکیب منطقی مفاهیم را در شبکه به‌درستی مدل‌سازی می‌کند.

در تلاشی دیگر برای ارتقای کیفی استخراج مفاهیم بصری، وانگ و همکاران (۲۰۲۴)~\cite{wang2024probabilistic} چارچوب «مفسرهای مفهومی احتمالی»\enfootnote{probabilistic conceptual explainers} را مشخصاً برای ترنسفورمرهای بینایی\enfootnote{Vision Transformer} (\lr{ViT})~\cite{dosovitskiy2021imageworth} ارائه دادند. آن‌ها برای ایجاد ضمانت‌های آماری نظیر استحکام در برابر نویز، از یک رویکرد بیزی متغیر\enfootnote{variational Bayesian} بهره بردند که در آن مفاهیم پنهان نه به‌صورت بردارهای قطعی، بلکه در قالب توزیع‌های احتمالی شرطی مدل‌سازی می‌شوند.

در همین راستا و برای مدل‌های نوین‌تر مولد تصویر، جین و همکاران (۲۰۲۴)~\cite{jin2024animage} با معرفی چارچوب «یادگیری پرامپت چندمفهومی»\enfootnote{multi-concept prompt learning} (\lr{MCPL})، بر چالش استخراج هم‌زمان مفاهیم مجزا از یک منظره‌ی واحد تمرکز کردند. در مسائل کلاسیک یادگیری پرامپت، هدف استخراج توکن‌ها یا کلمات ساختگی جدیدی است که هر کدام نماینده‌ی یک شیء خاص در یک تصویر باشند (به‌عنوان مثال، تخصیص کلمه $V_1$ منحصراً به یک حیوان و کلمه $V_2$ به یک شیء در کنار آن). محدودیت اصلی روش‌های پیشین در این بود که برای تفکیک این مفاهیم، نیازمند حاشیه‌نویسی دستی و برش زدن تصویر برای هر شیء بودند و در صورت تغذیه‌ی کل تصویر، مدل دچار اختلاط مفهومی می‌شد و تمامی توکن‌های جدید را به ترکیب کل اشیاء مرتبط می‌ساخت. متدولوژی \lr{MCPL} این چالش را بدون نیاز به تقطیع دستی تصویر برطرف می‌سازد. نوآوری محوری این رویکرد، ترکیب دو سازوکار منظم‌سازی مکمل است: نخست ماسک‌گذاری توجه\enfootnote{attention masking} که با میانگین‌گیری از نقشه‌های توجه متقاطع پرامپت‌ها در طول گام‌های نویززدایی و آستانه‌گذاری آن‌ها، خطای بازسازی مدل انتشاری را به نواحی مرتبط با اشیاء محدود می‌کند؛ و دوم زیان تقابلی پرامپت‌ها\enfootnote{prompts contrastive loss} بر روی تعبیه‌های پنهان پرامپت‌ها تا از تداخل و هم‌پوشانی معنایی توکن‌های مفهومی جلوگیری شود:
\begin{equation}\label{eq:chap3_mcpl}
    \ell_{i,j}^{\eta} = -\log \frac{\exp\big(\text{\lr{sim}}(v_i^\eta, v_j^\eta)/\tau\big)}{\sum_{\eta'=1}^N \sum_{j'} \exp\big(\text{\lr{sim}}(v_i^\eta, v_{j'}^{\eta'})/\tau\big)}
\end{equation}
در این رابطه، $\text{\lr{sim}}(u,v)$ شباهت کسینوسی تعبیه‌ها و $\tau$ پارامتر دمای زیان \lr{InfoNCE} است. جفت $(v_i^\eta, v_j^\eta)$ نمایانگر دو دیدگاه افزوده\enfootnote{augmented views} از تعبیه‌ی یک مفهوم یکسان ($\eta$) است که به‌عنوان نمونه‌های مثبت در نظر گرفته می‌شوند، در حالی که تعبیه‌ی سایر مفاهیم ($\eta' \neq \eta$) نمونه‌های منفی را تشکیل می‌دهند. از منظر شهودی، در حالی که ماسک‌گذاری توجه با اتکا به صفحات توجه متقاطع، هر توکن ساختگی را به مکان فیزیکی شیء در بستر تصویر قفل می‌کند، زیان تقابلی پرامپت‌ها به‌عنوان یک نیروی دافعه‌ی معنایی در فضای تعبیه عمل کرده و توکن‌های مربوط به مفاهیم مختلف (نظیر «خرس عروسکی» و «اسکیت‌برد») را از یکدیگر دور می‌سازد تا از فروپاشی بازنمایی‌ها در یک بردار منفرد جلوگیری کند. ارزیابی‌های این پژوهش اثبات کرد که با استفاده از این تکنیک می‌توان مفاهیم بصری چندگانه را با دقت بالایی از یک تصویر واحد تفکیک نمود.

اگرچه این روش‌ها ظرافت هندسی فضاهای پنهان را در دامنه‌ی بینایی به زیبایی آشکار می‌سازند، ارزیابی مکانیزم‌های آن‌ها پرده از یک تفاوت بنیادین برمی‌دارد که مفاهیم بصری غالباً از ویژگی محلی‌سازی فضایی\enfootnote{spatial localization} برخوردارند. اعمال مستقیمِ این منطق مبتنی بر جداسازی (مانند جریمه‌ی عدم‌هم‌پوشانی مکانی در متدولوژی \lr{MCPL} یا استفاده از ابرپیکسل‌ها) بر روی متون در مدل‌های زبانی با چالش‌های ساختاری مواجه می‌شود. بسیاری از مفاهیم در زبان طبیعی (مانند یک پیش‌فرض علّی، لحن، یا یک استدلال منطقی) ماهیتی به‌شدت توزیع‌شده دارند که در سراسر توکن‌های یک جمله یا پاراگراف پراکنده شده‌اند. در نتیجه، تلاش برای ایزوله کردنِ مطلق کلمات از طریق جریمه‌های سخت‌گیرانه‌ی مکانی می‌تواند به تخریب ساختار نحوی و فروپاشی منطق گزاره‌ای متن منجر شود.

برای غلبه بر این چالش‌های ذاتی متون و استخراج مفاهیم قابل‌درک، پژوهشگران به پارادایم یادگیری دیکشنری\enfootnote{dictionary learning} و اصل خلوتی روی آوردند. ایده‌ی محوری یادگیری دیکشنری این است که داده‌های پیچیده و متراکم را می‌توان به‌عنوان ترکیب خطی مجموعه‌ای از عناصر پایه‌ی ساده‌تر (اتم‌ها یا مفاهیم) بیان کرد. سابرامانیان و همکاران (۲۰۱۸)~\cite{subramanian2018spine} با تکیه بر این چارچوب، روش \lr{SPINE} را معرفی کردند. چالش اصلی در این حوزه آن بود که پس از اتمام آموزش مدل‌های پایه‌ی زبانی (نظیر \lr{Word2Vec}~\cite{mikolov2013efficient})، تعبیه‌های متراکمِ تولیدشده با وجود غنای معنایی بالا، به‌دلیل توزیع اطلاعات در تمامیِ ابعاد عملاً غیرقابل‌خوانش هستند. لذا \lr{SPINE} برای پیاده‌سازی عصبی یادگیری دیکشنری، از یک خودرمزگذار نویزگیرِ $k$-خلوت\enfootnote{denoising $k$-sparse autoencoder} بهره برد. این شبکه‌ی ثانویه آموزش می‌بیند تا بردارهای متراکم را بازسازی کند، اما با این قید ریاضی که در فضای میانی خود، همواره تنها $k$ نورون فعال داشته باشد. این معماریِ مبتنی بر شبکه‌ی عصبی، بازنمایی‌ها را وادار می‌کند تا مفاهیم درهم‌تنیده را جدا کرده و هر بُعد خلوت‌شده را با یک مفهوم مجزا هم‌تراز نماید.

در کنار این رویکرد، حل مسئله‌ی یادگیری دیکشنری از طریق الگوریتم‌های کلاسیک جبر خطی نیز مورد توجه قرار گرفت. در حالی که \lr{SPINE} از محدودیت‌های صریح شبکه‌ی عصبی برای خلوت‌سازی استفاده می‌کرد، فل و همکاران (۲۰۲۳)~\cite{fel2023holistic} با بهره‌گیری از الگوریتم تجزیه‌ی ماتریس نامنفی\enfootnote{non-negative matrix factorization} (\lr{NMF})، یک چارچوب جامع برای کشف ارگانیک مفاهیم از داخل شبکه‌های عمیق ارائه دادند. در این ساختار جبری، ماتریس فعال‌سازی‌های درونی شبکه ($A \in \mathbb{R}^{d \times n}$) با ابعاد $d$ نورون و $n$ نمونه، به حاصل‌ضرب ماتریس پایه‌های مفهومی ($U \in \mathbb{R}^{d \times k}$) که نقش کلمات دیکشنری را دارند، و ماتریس ضرایب فعال‌سازی این مفاهیم ($V \in \mathbb{R}^{k \times n}$) تجزیه می‌شود. فرایند بهینه‌سازی از طریق کمینه‌سازی تابع هزینه‌ی زیر انجام می‌پذیرد:
\begin{equation}\label{eq:chap3_nmf}
    \min_{U,V \ge 0} ||A - UV||_F^2
\end{equation}
از منظر شهودی، تفاوت کلیدی \lr{NMF} با روش‌های فاکتورگیری استاندارد (نظیر \lr{PCA} یا \lr{SVD}) در شرط نامنفی بودن هر دو ماتریس ($U, V \ge 0$) نهفته است. در روش‌های خطی کلاسیک، پایه‌ها می‌توانند مقادیر منفی بگیرند و بازسازی یک حالت داخلی با جمع و تفریق‌های پیچیده و جبران‌سازی‌های تخریبی صورت می‌گیرد که برای ذهن انسان نامفهوم است. اما قید نامنفی بودن، عمل تفریق را قدغن می‌کند؛ به این معنا که شبکه مجاز به ساخت هیچ بازنمایی از طریق کم کردن یک ویژگی از ویژگی دیگر نیست. در نتیجه، هر نمونه تنها از طریق جمع افزایشی اجزای سازنده (یک بازنمایی مبتنی بر اجزا\enfootnote{parts-based representation}، نظیر ترکیب افزایشی قطعات چهره در بینایی) شکل می‌گیرد که با منطق ادراک انسان کاملاً هم‌خوانی دارد و بردارهایی خلوت و خوانا تولید می‌کند. پس از استخراج این مفاهیم پایه‌ای (ستون‌های $U$)، فل و همکاران اهمیت هر مفهوم در خروجی مدل را با استفاده از شاخص‌های سوبول\enfootnote{Sobol indices} ارزیابی کردند.

با این حال، الگوریتم کلاسیک \lr{NMF} در مواجهه با مدل‌های زبانی بزرگ چندوجهی (\lr{MLLM}) به یک مانع جبری برخورد کرد. در معماری‌های نوین، توابع فعال‌سازی نظیر \lr{GeLU} و \lr{SwiGLU} جایگزین \lr{ReLU} شده‌اند و فعال‌سازی‌های داخلی لایه‌ها ($X \in \mathbb{R}^{d \times n}$) مقادیر منفی متعددی دارند که اعمال \lr{NMF} کلاسیک (که مستلزم نامنفی بودن $X$ و $U$ است) را ناممکن می‌سازد. پارخ و همکاران (۲۰۲۴)~\cite{parekh2024conceptbased} برای حل این بن‌بست، چارچوب فاکتورگیری ماتریس نیمه‌نامنفی\enfootnote{semi-nonnegative matrix factorization} (\lr{Semi-NMF}) را به‌کار گرفتند:
\begin{equation}\label{eq:chap3_semi_nmf}
    \min_{U \in \mathbb{R}^{d \times k}, V \in \mathbb{R}_{+}^{k \times n}} ||X - UV||_F^2
\end{equation}
از منظر شهودی، این معادله یک مصالحه‌ی ریاضی هوشمندانه است: قید نامنفی بودن از روی ماتریس پایه‌ها برداشته می‌شود ($U \in \mathbb{R}^{d \times k}$) تا مفاهیم بتوانند مختصات منفی فضای پنهان ناشی از \lr{GeLU} را نیز بازنمایی کنند؛ اما قید نامنفی بودن ضرایب ترکیب ($V \ge 0$) با قاطعیت حفظ می‌شود. این یعنی با وجود اینکه ابعاد بردار مفهوم در فضای پنهان می‌تواند مثبت یا منفی باشد، فعال‌سازی هر مفهوم در یک ورودی مشخص صرفاً می‌تواند به‌صورت یک ضریب مثبت (دستور ترکیب افزایشی) ظاهر شود. بدین‌ترتیب، اصل شهودی عدم‌تفریق مفاهیم و ماهیت مبتنی بر اجزای تفسیرپذیری بدون تضاد با هندسه‌ی لایه‌های مدرن پابرجا می‌ماند.

علی‌رغم موفقیت ریاضی یادگیری دیکشنری در تجزیه فضای پنهان، تمامی این رویکردهای بدون‌نظارت با یک بن‌بست بنیادین به‌نام شکاف معنایی\enfootnote{semantic gap} روبه‌رو هستند. اتم‌ها و بردارهای استخراج‌شده در این الگوریتم‌ها صرفاً موجودیت‌هایی آماری و انتزاعی‌اند که فاقد هرگونه برچسب زبان طبیعی هستند. تخصیص یک نام یا گزاره‌ی انسانی به این اعداد همچنان نیازمند بازرسی چشمی و تفسیر ذهنی محقق است که اعتبار عینی آن‌ها را تنزل می‌دهد. مهم‌تر از آن، کشف مجموعه‌ای از بردارهای متعامد یا خلوت به خودی خود ثابت نمی‌کند که مدل واقعاً بر اساس این بردارها استدلال می‌کند. فقدان تضمین برای وجود یک رابطه‌ی علّی قدرتمند میان این موجودیت‌های ریاضی و رفتار نهایی مدل، پژوهشگران را ملزم ساخت تا برای این مفاهیم خام، مکانیزم‌های ارزیابی معنایی و اثبات علیت طراحی کنند.

\paragraph{ارزیابی معنایی و علیت در مفاهیم.}

برای پر کردن شکاف معنایی باقی‌مانده از رویکردهای کشف بدون‌نظارت، پژوهشگران نیازمند مکانیزم‌هایی بودند تا بتوانند به بردارهای هندسی استخراج‌شده معنا ببخشند، برچسب‌های انسانی به آن‌ها اختصاص دهند و در نهایت ثابت کنند که این بردارها واقعاً علت پنهان در پسِ تصمیمات شبکه هستند.

در این مسیر، نخستین گام استفاده از شبکه‌های ارزیاب خارجی برای ترجمه‌ی جهات ریاضی به مفاهیم ادراکی بود. در حوزه‌ی مدل‌های مولد تصویر، گوتشالک و همکاران (۲۰۱۹)~\cite{goetschalckx2019ganalyze} متوجه شدند که برخی ویژگی‌های شناختیِ پیچیده در انسان (مانند اینکه یک تصویر چقدر برای ذهن «به‌یادماندنی» است) فرمول ریاضی یا تعریف پیکسلی مشخصی ندارند. آن‌ها برای کنترل و تفسیر این ویژگی‌ها در خروجی شبکه‌های مولد تخاصمی، چارچوب \lr{GANalyze} را معرفی کردند. در این متدولوژی، یک شبکه‌ی ارزیاب خارجی که از پیش برای سنجش میزان به‌یادماندنی بودن تصاویر آموزش دیده است، به خروجی مدل مولد متصل می‌شود. الگوریتم سپس در فضای پنهان مدل مولد، به دنبال یک جهت یا بردار خطی می‌گردد؛ شرط صحت این بردار آن است که اگر مختصات پنهان را در امتداد آن حرکت دهیم، تصاویر تولیدشده به‌صورت تدریجی به‌یادماندنی‌تر شوند و این افزایش کیفی، مستقیماً توسط شبکه‌ی ارزیاب و به همان نسبت جابه‌جایی تأیید شود. بدین‌ترتیب، یک مفهوم گنگ انسانی به یک بردار قابل‌کنترل در فضای پنهان گره می‌خورد.

با این وجود، چارچوب‌هایی نظیر \lr{GANalyze} نیازمند آموزش یک ارزیاب اختصاصی برای هر مفهومِ منفرد بودند. برای رفع این محدودیت و ایجاد یک مفسر همه‌منظوره، ژو و همکاران (۲۰۲۵)~\cite{zhu2025latentexplainer} چارچوب \lr{LatentExplainer} را توسعه دادند. این پژوهش از \lr{MLLM}ها به‌عنوان یک مدل خارجی جامع استفاده می‌کند. در این روش، ابعاد پنهان مدل مولد به‌صورت سیستماتیک دستکاری می‌شوند و تصاویرِ تولیدشده در اثر این تغییرات، پشت‌سرهم در قالب یک دنباله‌ی بصری (نظیر یک ویدئو از تغییر تدریجی تصویر) به مدل مفسر تغذیه می‌گردند. سپس از این مدل خارجی خواسته می‌شود تا با تکیه بر درک چندوجهی خود، مشاهده کند که چه چیزی در تصاویر در حال تغییر است و آن را به زبان طبیعی توصیف کند (مثلاً بنویسد که «این بُعد پنهان، میزان روشنایی پس‌زمینه را کنترل می‌کند»). اگرچه این ایده‌ی مبتکرانه نیاز به حاشیه‌نویسی و ساخت ارزیاب‌های دستی را به‌شدت کاهش می‌دهد، اما اعتبار توضیح نهایی وابستگی مطلقی به دقت و استدلال مدل ارزیاب پیدا می‌کند و هرگونه کوری مفهومی یا توهم در این مدل خارجی، به‌عنوان منطق مدل اصلی ثبت می‌گردد.

پس از آنکه مفاهیم پنهان شناسایی و نام‌گذاری شدند، پرسش بنیادین دیگری در ادبیاتِ ارزیابی مطرح گردید، مبنی بر اینکه آیا مجموعه‌ی مفاهیم کشف‌شده، کل داستان را روایت می‌کنند و آیا واقعاً همان متغیرهایی هستند که مدل بر اساس آن‌ها تصمیم نهایی خود را می‌گیرد.

در راستای پاسخ به پرسش اول، یه و همکاران (۲۰۲۰)~\cite{yen2020completeness} با تمرکز بر وظایف بینایی ماشین و شبکه‌های دسته‌بند تصویر، مسئله‌ی کامل‌بودن را مطرح کرده و چارچوب \lr{ConceptSHAP} را توسعه دادند. ایده‌ی محوری آن‌ها این است که یک مجموعه‌ی مفهوم ($V$)، تنها زمانی کامل و جامع است که بتوان خروجی پیچیده‌ی مدل اصلی ($f(x)$) را صرفاً با استفاده از امتیاز فعال‌سازی این مفاهیم ($v(x)$) و بدون نیاز به داده‌های اصلی با دقت بالایی پیش‌بینی کرد. به بیان ساده‌تر، باید بتوان یک مدل جایگزین بسیار ساده ($g$) آموزش داد که تنها با نگاه کردن به مقادیر این مفاهیم، دقیقاً همان جوابی را بدهد که مدل اصلی داده است، یعنی $g(v(x)) \approx f(x)$.

پس از تأیید کامل‌بودن مجموعه، نوبت مرحله تعیین میزان اهمیت هر مفهوم مجزا در تصمیم‌گیری است. نویسندگان برای توزیع عادلانه‌ی این اهمیت، از مقادیر شاپلی\enfootnote{Shapley values} استفاده کردند (منطق ریاضی و فرمول‌بندی دقیق مقادیر شاپلی در ادامه‌ی این فصل و در بخش رویکردهای جعبه‌سیاه به‌تفصیل تشریح خواهد شد). در این متدولوژی، برخلاف روش‌های سنتی که اهمیت را بین ویژگی‌های خام (مثل پیکسل‌های تصویر) تقسیم می‌کنند، اهمیت بین خود مفاهیم تقسیم می‌شود؛ به این صورت که تأثیر یک مفهوم با بررسی میزان افت عملکرد مدل جایگزین در هنگام حذف آن مفهوم از گروه‌های مختلف، محاسبه می‌گردد. با وجود استحکام ریاضی این چارچوب، محاسبه‌ی دقیق مقادیر شاپلی برای $K$ مفهوم دارای پیچیدگی نمایی ($\mathcal{O}(2^K)$) است. افزون بر آن، اتکا به یک مدل جایگزین برای اثبات کامل‌بودن، همچنان تخمینی غیرمستقیم از رفتار شبکه‌ی اصلی است.

در نهایت، برای اثبات اینکه مفاهیم کشف‌شده واقعاً محرک خروجی مدل هستند، باید از تله‌ی همبستگی عبور کرد. بسیاری از روش‌های پیشین تنها نشان می‌دادند که حضور یک مفهوم با یک خروجی خاص همبستگی دارد (به این معنا که هر دو معمولاً با هم در داده‌ها ظاهر می‌شوند)، اما این هم‌زمانی اثبات نمی‌کند که مفهوم استخراج‌شده علت بروز آن خروجی بوده است. برای کشف اثرگذاری و عبور از دامنه‌ی تصویر به سمت پردازش متن، ژائو و همکاران (۲۰۲۴)~\cite{zhao2024explaining} چارچوب مفاهیم پرنفوذ\enfootnote{high-impact concepts} (\lr{HI-Concept}) را منحصراً برای مدل‌های زبانی ارائه دادند. در این متدولوژی با تجزیه‌ی مدل به دو بخش استخراج ویژگی و پیش‌بینی ($f = \psi \circ \phi$)، یک شبکه‌ی گلوگاه خودرمزگذار ($g_\theta$) تعبیه می‌شود تا با تبدیل فعال‌سازی‌های میانی $\phi(x)$ به احتمال وقوع مفاهیم ($p_{\mathcal{C}}(x)$)، بازنمایی‌ها را به‌صورت $g_\theta(p_{\mathcal{C}}(x)) \approx \phi(x)$ بازسازی کند. نوآوری کلیدی آن‌ها، تمایز میان ویژگی‌های صرفاً همبسته با ویژگی‌های واقعاً علّی از طریق وضع یک زیان علّیت\enfootnote{causality loss} مبتنی بر مداخله‌ی حذفی است:
\begin{equation}\label{eq:chap3_high_impact_concepts}
    \mathcal{L}_{\text{\lr{cau}}}(\theta, \mathcal{C}) = -\sum_{c_i \in \mathcal{C}} \sum_{x \in \mathcal{D}} \Big| \psi\big(g_\theta(p_{\mathcal{C}}(x) \mid c_i = 0)\big) - \psi\big(g_\theta(p_{\mathcal{C}}(x) \mid c_i = c_i)\big) \Big|
\end{equation}
از منظر شهودی، این معادله یک آزمون استرس علّی مستقیم است: الگوریتم با صفر کردن اجباری یک مفهوم خاص ($c_i = 0$) در گلوگاه بازنمایی، بررسی می‌کند که خروجی لایه‌های نهایی ($\psi$) چه مقدار دچار شوک و افت پیش‌بینی می‌شود. با کمینه‌سازی این زیان (بیشینه‌سازی قدرمطلق افت)، شبکه‌ی گلوگاه وادار می‌گردد مفاهیمی را در بردارهای $\mathcal{C}$ کدگذاری کند که حضورشان برای پیش‌بینی نهایی تعیین‌کننده است؛ بنابراین مفاهیمی که صرفاً به‌دلیل همبستگی نویزی داده‌ها در لایه‌های میانی ظاهر شده‌اند و حذفشان تغییری در تصمیم نهایی ایجاد نمی‌کند، کنار گذاشته می‌شوند.

با وجود اینکه این چارچوب در تفکیک علیت از همبستگی دستاورد مهمی محسوب می‌شود، اتکا به یک مدل خودرمزگذار جانشین ($g_\theta$) برای میانجی‌گری بازنمایی‌ها، خطای تقریب به‌همراه دارد. از سوی دیگر، آزمودن حذف مفاهیم به‌صورت ایزوله، از مدل‌سازی برهم‌کنش‌های غیرخطی پیچیده و پدیده‌ی انطباق چندگانه میان مفاهیم در طول تولید متن آزاد ناتوان است. این چالش‌ها نشان می‌دهد که ارزیابی مفاهیمِ پیچیده‌ی زبانی نیازمند کاوش مستقیم در دینامیک نظری و ساختار هندسی درونی خود مدل‌های مولد است.

\paragraph{دینامیک نظری و هندسه‌ی درونی مدل‌های مولد.}

برای انجام این بازنگری، باید مفروضات پایه‌ای که تمامی روش‌های استخراج پیشین بر آن‌ها استوارند را کالبدشکافی کرد. یکی از مهم‌ترین این موارد، باوری است که بیان می‌دارد فضای پنهان شبکه‌های عصبی دارای یک ساختار هندسی خطی است که در آن مفاهیم به‌صورت بردارهای مجزا قابل‌کاوش هستند. پژوهش‌های نظری اخیر اثبات کرده‌اند که اتکا به این هندسه در برخورد با ذات درهم‌تنیده‌ی زبان، پیچیدگی‌های تئوریک عمیقی به‌همراه دارد که واکاوی آن‌ها، آخرین قطعه از پازل پارادایم مبتنی بر مفهوم را شکل می‌دهد.

هسته‌ی مرکزی این بحث، فرضیه‌ی بازنمایی خطی\enfootnote{linear representation hypothesis} است. این فرضیه بیان می‌کند که شبکه‌های عصبی، مفاهیم پیچیده‌ی انتزاعی و معنایی را به ساده‌ترین شکل هندسیِ ممکن یعنی در قالب خطوط راست و بردارهای یک‌بعدی در فضای متراکم پارامترهای خود ذخیره می‌کنند. جیانگ و همکاران (۲۰۲۴)~\cite{jiang2024origins} در یک واکاوی نظری کلیدی، منشأ شکل‌گیری این ساختار خطی را در مدل‌های زبانی با استفاده از مفهوم بردارهای هدایت‌گر\enfootnote{steering vectors} بررسی کردند. اگر بردار هدایت‌گر برای مفهوم $i$ام در فضای تعبیه‌زدایی\enfootnote{unembedding space} به‌صورت تفاضل بازنمایی در حالت فعال و غیرفعال بودن آن مفهوم تعریف شود ($\Delta_{c,i} = g(c^{(i \to 1)}) - g(c^{(i \to 0)})$)، بازنمایی خطی زمانی محقق می‌گردد که این بردارهای هدایت‌گر در تمامی بافت‌ها و جملات مختلف با یکدیگر موازی باشند ($\Delta_{c,i} \parallel \Delta_{c',i}$). آن‌ها در قضیه‌ی اصلی خود اثبات کردند که اگر مدل زبانی در فرایند پیش‌بینی توکن بعدی بتواند نسبت لگاریتم شانس\enfootnote{log-odds} داده‌ها را برای مفهوم $C_i$ در بافت متنی $d$ به‌درستی فراگیرد:
\begin{equation}\label{eq:chap3_linear_representation_hypothesis}
    \ln \frac{\hat{p}(c^{(i \to 1)} \mid d)}{\hat{p}(c^{(i \to 0)} \mid d)} = \ln \frac{p(C_i = 1)}{p(C_i = 0)}
\end{equation}
آن‌گاه تمامی بردارهای هدایت‌گر متناظر با مفهوم $C_i$ الزاماً موازی یکدیگر خواهند شد. از منظر شهودی، این معادله بیان می‌کند که فعال شدن یک مفهوم مشخص، مستقل از بافت پیرامونی جمله، همواره باید اثر افزایشی ثابتی بر لگاریتم احتمال خروجی بگذارد؛ این تقارن آماری شبکه را مجبور می‌سازد تا آن مفهوم را صرفاً در قالب یک راستای برداری یکتا و موازی کدگذاری کند. افزون بر این، نویسندگان نشان دادند که سوگیری ضمنی روش گرادیان کاهشی در بهینه‌سازی تابع زیان آنتروپی متقاطع (که مشابه دسته‌بندهای با بیشینه‌ی حاشیه عمل می‌کند)، محرک اصلی این هم‌راستاسازی خطی در فضای پنهان است و اثبات می‌کند که خطی بودن مفاهیم، نتیجه‌ی هندسی گریزناپذیر یادگیری پیش‌بینی توکن بعدی است.

با این وجود، اتکای بی‌چون‌وچرا به فرضیه‌ی بازنمایی خطی و استفاده از ابزارهایی نظیر کاوشگرها برای اثبات یادگیری یک مفهوم، از منظر ریاضی با نقدهای بنیادینی مواجه است. رِج (۲۰۲۴)~\cite{rege2024probe} با ارائه‌ی یک چارچوب نظری جامع، قابلیت‌اطمینان پارادایم کاوشگرها را در مدل‌های مولد مورد ارزیابی انتقادی قرار داد. مسئله‌ی اصلی در این پژوهش، رفتار هندسی فضاهای برداری با ابعاد بسیار بالا است. رِج اثبات می‌کند که در فضاهای پنهان مدل‌های زبانی مدرن که هزاران بُعد دارند، یک دسته‌بند خطی به‌دلیل محدودیت‌های ناشی از پیچیدگی نمونه به‌شدت مستعد فریب خوردن است. از نظر ریاضی، اگر تعداد داده‌های استفاده‌شده برای آموزش کاوشگر (تعداد نمونه‌ها)، به نسبتِ ابعاد فضا به اندازه‌ی کافی عظیم نباشد، کاوشگر می‌تواند به‌راحتی یک ابرصفحه‌ی جداکننده برای الگوهای نویزی و کاملاً تصادفی پیدا کند و محقق را به این اشتباه بیندازد که مدل یک مفهوم معنایی را یاد گرفته است؛ در حالی که کاوشگر صرفاً یک همبستگی آماری پوچ را در آن فضای وسیع شناسایی کرده است. این شکنندگی تئوریک نشان می‌دهد که کشف موفقیت‌آمیز یک مسیر خطی، به‌هیچ‌وجه تضمین‌کننده‌ی وابستگی علّی مدل به آن مفهوم در فرایند استنتاج نیست.

دقیقاً در واکنش به همین نارسایی‌های هندسی و برای عبور از تله‌ی همبستگی‌های تصادفی، نحوه‌ی اندازه‌گیری و تفسیر این فضاها مورد بازبینی بنیادین قرار گرفت. پارک و همکاران (۲۰۲۳)~\cite{park2023thelinear} استدلال کردند که استفاده از معیارهای هندسی کلاسیک نظیر فاصله‌ی اقلیدسی یا شباهت کسینوسی برای یافتن مفاهیم، لزوماً با رفتار واقعی و علّی مدل سازگار نیست. آن‌ها برای اصلاح این دیدگاه از چارچوب انتزاع علّی\enfootnote{causal abstraction} استفاده کردند. بر اساس این چارچوب، ما تنها زمانی می‌توانیم ادعا کنیم یک بردار خاص ($v_c$) نمایانگر مفهوم جنسیت است که اگر به‌صورت دستی مقداری از این بردار را به بازنمایی پنهان شبکه اضافه کنیم ($h \leftarrow h + \alpha v_c$)، خروجی نهایی مدل دقیقاً به همان شکلی تغییر کند که انگار کلمات مربوط به جنسیت را در متن ورودی تغییر داده‌ایم (یک مداخله‌ی پادواقعی). برای سنجش دقیق این تطابق، نویسندگان استفاده از ضرب داخلی علّی\enfootnote{causal inner product} را پیشنهاد دادند. در این معیار شباهت دو بردار بر اساس زاویه یا فاصله‌ی خام آن‌ها در فضا سنجیده نمی‌شود، بلکه بر این اساس ارزیابی می‌گردد که دستکاری و جابه‌جایی در امتداد آن بردارها، تا چه حد تغییر و تأثیر مشابهی بر احتمالات خروجی شبکه ایجاد می‌کند. آن‌ها بدین صورت اثبات کردند که فضای پنهان مدل‌های زبانی یک فضای همگن اقلیدسی نیست، بلکه هندسه‌ی آن به‌شدت توسط گراف محاسباتی و مسیر عبور اطلاعات در شبکه تعیین می‌شود.

بُعد دیگر این پیچیدگی نظری، توزیع مکانی مفاهیم در طول معماری شبکه است. جین و همکاران (۲۰۲۵)~\cite{jin2025exploring} با معرفی سنجه‌ی عمق مفهوم، توزیع لایه‌ای مفاهیم را واکاوی کردند. بسیاری از روش‌های پیشینِ استخراج مفهوم نیازمند آن بودند که محقق یک لایه‌ی پنهان ثابت را به‌عنوان منبع استخراج ویژگی‌ها انتخاب کند. یافته‌های جین و همکاران نشان داد که شکل‌گیری مفاهیم در \lr{LLM}ها یک فرایند سلسله‌مراتبی و پویاست. مفاهیم نحوی و واژگانی ساده در لایه‌های کم‌عمق اولیه قابل‌رمزگشایی هستند، در حالی که دانش معنایی پیچیده و مفاهیم انتزاعی در لایه‌های عمیق‌تر به بلوغ می‌رسند. این پویایی لایه‌ای ثابت می‌کند که استخراج بازنمایی‌ها از یک لایه‌ی منفرد، تصویری به‌شدت ناقص از استدلال مدل ارائه می‌دهد و تحلیل جامع نیازمند رهیافت‌های کل‌نگری است که تکامل یک مفهوم را به‌صورت پیوسته در طول لایه‌های متعدد ارزیابی کنند.

در نهایت، پویایی زمانی در طول فرایند آموزش نیز ابعاد شگفت‌انگیزی از نحوه‌ی درونی‌سازی مفاهیم را نمایان ساخته است. بررسی‌های پارک و همکاران (۲۰۲۴)~\cite{park2024emergence} که بر روی دامنه‌ی گسترده‌ای از مدل‌های مولد (شامل ویژگی‌های بصری نظیر شکل و رنگ در کنار مفاهیم متنی) انجام شد، پدیده‌ی قابلیت‌های پنهان را معرفی نمود. آن‌ها برای درک نحوه‌ی یادگیری، یک فضای مفهومی تعریف کردند که در آن هر محور نمایانگر یک کیفیت مستقل در داده‌هاست. سپس آن‌ها با توجه به این فضا، مسیر یادگیری مدل را در طول گام‌های به‌روزرسانی گرادیان ردیابی کردند. یافته‌ی کلیدی آن‌ها وقوع چرخش‌های ناگهانی در این مسیر بود؛ نقاطی که در آن‌ها شبکه هندسه‌ی یک مفهوم را در فضای پنهان خود شکل می‌دهد.

شگفتی این پژوهش در آن است که مدل این مفاهیم خطی را مدت‌ها پیش از آنکه اثرشان در رفتار بیرونی (توزیع احتمالات خروجی) ظاهر شود، در لایه‌های داخلی خود یاد می‌گیرد. نویسندگان با استفاده از مداخلات فضای پنهان اثبات کردند که شبکه درک کاملی از مفهوم پیدا کرده است، اما پرامپت‌دهی ساده قادر به استخراج و نمایش آن در خروجی نیست. این یافته ادعای مهمی را ثابت می‌کند مبنی بر اینکه تکیه صرف بر خروجی‌های متنی، دانش واقعی و ظرفیت‌های استدلالی مدل را به‌شدت دست‌کم می‌گیرد. از این رو استخراج مفاهیم نیازمند ابزارهایی است که بتوانند میان آنچه مدل واقعاً می‌فهمد و آنچه تصمیم می‌گیرد در خروجی بروز دهد، تمایز قائل شوند.


\subsubsection{تفسیرپذیری مکانیکی و مهندسی معکوس}\label{sec:3-1-2-3}

در حالی که رویکردهای انتساب اعتبار (بخش~\ref{sec:3-1-2-1}) بر یافتن میزان اهمیت ورودی‌ها تمرکز داشتند و متدولوژی‌های استخراج مفهوم (بخش~\ref{sec:3-1-2-2}) به دنبال کشف موجودیت‌های معنایی در فضاهای پنهان بودند، هیچ‌یک به‌طور دقیق مشخص نمی‌کردند که شبکه چگونه و طی چه الگوریتمی این اطلاعات را پردازش می‌کند. برای پاسخ به این پرسش بنیادین، پارادایم تفسیرپذیری مکانیکی شکل گرفت. هدف اصلی در این حوزه، مهندسی معکوس شبکه‌های عصبی است؛ روندی که در آن گراف‌های محاسباتی عظیم و ماتریس‌های درهم‌تنیده‌ی وزن‌ها، به الگوریتم‌ها و مدارهای خوانا برای انسان تجزیه می‌شوند تا استدلال مدل در سطح کدهای کامپایل‌شده درک گردد.

\paragraph{مبانی تفسیرپذیری مکانیکی.}

ریشه‌های تاریخیِ مداخله در اجزای داخلی شبکه برای کشف الگوریتم‌های پنهان، به حوزه‌ی بینایی ماشین بازمی‌گردد. یکی از اساسی‌ترین پژوهش‌ها در این زمینه توسط باو و همکاران (۲۰۱۹)~\cite{bau2018visualizing} انجام شد که در آن چارچوب «کالبدشکافی \lr{GAN}» را برای درک نحوه بازنمایی مفاهیم در واحدهای داخلی (کانال‌ها یا نقشه‌های ویژگی) شبکه‌های مولد معرفی کردند. 

متدولوژی آن‌ها شامل یک ارزیابی سه‌مرحله‌ای بود. در گام نخست از یک شبکه‌ی تقطیع معنایی\enfootnote{semantic segmentation network} خارجی برای شناسایی مفاهیم در تصاویر تولیدشده استفاده می‌شد. سپس میزان هم‌پوشانی مکانی میان فعال‌سازی یک واحد داخلی خاص با نقشه‌ی تقطیع آن مفهوم محاسبه می‌گردید. در گام نهایی، علیت این ارتباط از طریق مداخله اثبات می‌شد؛ بدین صورت که با غیرفعال کردن (صفر کردن) یا فعال‌سازی اجباری آن واحد در گراف محاسباتی، یک شیء مشخص (نظیر «درخت» یا «پنجره») از خروجی حذف شده یا به آن اضافه می‌گردید. با این وجود، اگرچه این روش برای شبکه‌های پیچشیِ\enfootnote{convolutional network} دارای محلی‌سازی فضایی موفقیت‌آمیز بود، انتقال منطق تحلیل در سطح واحد منفرد به پردازش زبان طبیعی با شکست مواجه شد. در متون، مفاهیم دارای ماهیتی به‌شدت توزیع‌شده هستند و به‌ندرت در یک واحد یا نورون منفرد محصور می‌شوند که این امر نیازمند یک بازنگری بنیادین در معماری تفسیرپذیری بود.

این بازنگری ساختاری با پژوهش محوریِ الهیج و همکاران (۲۰۲۱)~\cite{elhage2021mathematical} رقم خورد که بنیان‌گذار مدرن پارادایم تفسیرپذیری مکانیکی برای معماری ترنسفورمر محسوب می‌شود. در این چارچوب نظری، مدل‌های زبانی دیگر به‌عنوان یک جعبه‌سیاهِ یکپارچه دیده نمی‌شوند، بلکه به‌عنوان شبکه‌ای از مدارهای محاسباتی تقلیل می‌یابند. نویسندگان مفهوم جریان باقیمانده را به‌عنوان یک شاهراه اطلاعاتی مرکزی در ترنسفورمر فرمول‌بندی کردند. در این مدل‌سازی، ابعاد جریان باقیمانده به‌عنوان حافظه‌ی مدل عمل کرده و هر لایه صرفاً اطلاعات جدیدی را از طریق جمع خطی به این جریان می‌افزاید. به بیانی دیگر:
\begin{equation}\label{eq:chap3_residual_stream}
    x_l = x_{l-1} + \text{\lr{Layer}}_l(x_{l-1})
\end{equation}
در این رابطه، $x_l$ بردار حالت جریان باقیمانده در لایه‌ی $l$، $x_{l-1}$ وضعیت انباشته‌شده از لایه‌ی پیشین، و $\text{\lr{Layer}}_l(x_{l-1})$ بردار به‌روزرسانی افزایشی\enfootnote{additive update} است که توسط لایه‌ی توجه یا شبکه‌ی پیش‌خور محاسبه می‌شود. به بیان شهودی، با بسط بازگشتی این معادله به‌صورت $x_L = x_0 + \sum_{l=1}^L \text{\lr{Layer}}_l(x_{l-1})$، جریان باقیمانده شبیه به یک گذرگاه ارتباطی مشترک خطی (شبیه به یک باس داده) عمل می‌کند؛ هیچ لایه‌ای کل بازنمایی پیشین را بازنویسی نمی‌کند، بلکه هر مؤلفه اطلاعات مورد نیاز خود را از این گذرگاه خوانده و پس از پردازش، بردار حاصل را مستقیماً روی آن جمع می‌زند. خطی بودن این تجمیع سبب می‌شود که اطلاعات پردازش‌شده در لایه‌های ابتدایی، بدون نیاز به گذر اجباری و دگرگونی در تبدیل‌های غیرخطی لایه‌های میانی، مستقیماً برای لایه‌های بعدی و سر رمزگشای نهایی در دسترس باقی بماند. بر پایه‌ی این دیدگاه، سرهای مکانیزم توجه به‌عنوان عملگرهای مستقلی تعریف شدند که اطلاعات را از جریان باقیمانده می‌خوانند و پس از پردازش، نتایج را دوباره روی آن می‌نویسند. این پژوهش برای نخستین بار تعاملات سرها را به دو مدار تفکیک‌شده‌ی مجزا تجزیه کرد: مدار پرسمان-کلید (\lr{QK}) که از طریق مقایسه‌ی پرسمان‌ها و کلیدها تعیین می‌کند هر توکن به کدام توکن‌های دیگر توجه کند (نقشه‌ی هدایت توجه)، و مدار خروجی-ارزش (\lr{OV}) که مشخص می‌سازد چه بردار محتوایی از موقعیت مبدأ خوانده شده و به جریان باقیمانده افزوده شود. افزون بر این، با کشف موجودیت‌هایی نظیر سرهای استقرایی\enfootnote{induction heads} به‌عنوان مکانیسم اصلی یادگیری درون‌بافتاری\enfootnote{in-context learning}، اثبات گردید که رفتارهای کلانِ زبانی مستقیماً حاصل مهندسی معکوس همین تعاملات مداری هستند.

برای درک انسجام و تکامل این قلمروی نوظهور، رای و همکاران (۲۰۲۵)~\cite{rai2025practicalreview} یک مرور ساختاریافته بر ادبیات تفسیرپذیری مکانیکی در مدل‌های ترنسفورمر ارائه دادند. این پژوهش، فرایند مهندسی معکوس را در سه محور اصلی دسته‌بندی می‌کند که به‌ترتیب شامل شناسایی جهت‌های معنادار در فضای فعال‌سازی، تحلیل مدارها برای ردیابی جریان اطلاعات میان اجزای مدل (مانند تعاملات سرها و لایه‌های پیش‌خور)، و نهایتاً کنترل رفتار مدل از طریق مداخلات علّی بر روی این متغیرهای داخلی است. با این حال، اهمیت این واکاوی جامع در برجسته‌سازی صریح چالش‌های مقیاس‌پذیریِ این پارادایم نهفته است. نویسندگان استدلال می‌کنند که استخراج گراف‌های محاسباتی و تحلیل مدارها برای مدل‌های مرزی با ده‌ها میلیارد پارامتر، با موانع محاسباتی فلج‌کننده‌ای روبه‌رو است. رهیافت‌های مکانیکی، علی‌رغم ظرافت تئوریک، مستلزم ارزیابی ترکیبی میلیون‌ها یال محاسباتی هستند و تخصیص معنا به اجزای کشف‌شده همچنان نیازمند مداخلات دستی و تقریب‌های سنگین آماری است؛ موانعی که محققان را برای کشف مدارها در مقیاس‌های کلان، به‌سوی توسعه‌ی روش‌های خودکارتر سوق داده است.

\paragraph{کشف مدارها و ردیابی مسیرهای محاسباتی.}

پس از بسط چارچوب نظری مدارها در معماری ترنسفورمر، چالش عملی اصلی برای پژوهشگران، یافتن کارآمد این زیرشبکه‌ها در میان میلیاردها یال محاسباتی بود. تکنیک استاندارد برای این منظور، وصله‌زنی فعال‌سازی\enfootnote{activation patching} است؛ فرایندی که در آن فعال‌سازی یک جزء خاص (مانند یک سر توجه) در حین پردازش یک پرامپت، با فعال‌سازی همان جزء در یک پرامپت دیگر جایگزین می‌شود تا اثر علّی آن بر خروجی سنجیده شود. با این حال، اجرای این روش به‌صورت جامع (تست تک‌تک یال‌ها برای یافتن مدار)، دارای پیچیدگی ترکیبیاتی است و نیازمند هزاران عبور پیش‌رو در شبکه است که بار محاسباتی فلج‌کننده‌ای به‌همراه دارد. 

برای غلبه بر این فاجعه‌ی محاسباتی، سید و همکاران (۲۰۲۴)~\cite{syed2024attribution} متدولوژی «وصله‌زنی انتسابی»\enfootnote{attribution patching} را پیشنهاد دادند. این چارچوب به‌جای جایگزینی فیزیکی و اجرای مجدد مدل برای هر یال، از یک تقریب خطی مبتنی بر بسط تیلور مرتبه اول استفاده می‌کند تا اثر علّی وصله‌زنی را تخمین بزند. از نظر ریاضی، اگر $a_i$ فعال‌سازیِ یال $i$ در ورودی اصلی (سالم) و $a_i^*$ فعال‌سازی همان یال در ورودی مبدأ (پادنمونه یا مخدوش) باشد، تغییر خروجی در اثر وصله‌زنی این یال ($\Delta L$) به‌صورت زیر تقریب زده می‌شود:
\begin{equation}\label{eq:chap3_attribution_patching}
    \Delta L \approx {\left( a_i^* - a_i \right)}^T \nabla_{a_i} L
\end{equation}
در این معادله، بردار تفاضلی $(a_i^* - a_i)$ نشان‌دهنده‌ی بردار جابه‌جایی یا آشفتگی فعال‌سازی است و $\nabla_{a_i} L$ گرادیان تابع خروجی (یا معیار زیان هدف $L$) نسبت به فعال‌سازی یال $a_i$ است که جهت بیشترین حساسیت خروجی را مشخص می‌سازد. از منظر شهودی، حاصل‌ضرب داخلی در این معادله، تصویرِ بردار آشفتگی را بر روی راستای حساسیت مدل اندازه می‌گیرد؛ به این معنا که اگر جابه‌جایی ناشی از پادنمونه در جهتی رخ دهد که خروجی مدل به آن حساس است، $\Delta L$ بزرگی پدید می‌آید که نشان‌دهنده‌ی اهمیت علّی و حیاتی آن یال در مدار محاسباتی است، در حالی که جابه‌جایی‌های متعامد بر گرادیان عملاً تأثیری بر خروجی نگذاشته و یال‌های متناظر می‌توانند هرس شوند. مزیت انقلابی این متدولوژی آن است که تمامی گرادیان‌های $\nabla_{a_i} L$ در کل شبکه تنها طی یک عبور پس‌رو (تفاضل‌گیری خودکار) به‌صورت موازی محاسبه می‌شوند؛ بنابراین با تنها دو عبور پیش‌رو (برای به‌دست آوردن $a_i$ و $a_i^*$) و یک عبور پس‌رو، اهمیت تمامی یال‌های گراف محاسباتی هم‌زمان تقریب زده شده و مدار نهایی هرس می‌گردد. با این وجود، اتکای این روش به یک تقریب خطی از سیستمی که به‌شدت غیرخطی است، باعث می‌شود تا تعاملات پیچیده‌تر و چندجانبه میان اجزای شبکه، در فرایند استخراج مدار نادیده گرفته شوند.

علی‌رغم محدودیت‌های ناشی از تقریب‌های خطی، ابزارهای استخراج مدار بستری را فراهم کردند تا دینامیکِ تکاملِ مدل‌های زبانی در طول فازهای مختلف آموزش مورد واکاوی قرار گیرد. یکی از پرسش‌های اساسی این بود که یک مدل پایه‌ی پیش‌آموزش‌دیده (مانند \lr{LLaMA}~\cite{touvron2023llama}) در فرایند تنظیم مبتنی بر دستورالعمل\enfootnote{instruction-tuning} چگونه تغییر می‌کند تا به‌جای صرفاً ادامه دادن متن، از قصد کاربر پیروی کند. وو و همکاران (۲۰۲۴)~\cite{wu2024language} با استفاده از روش‌های انتساب مبتنی بر گرادیان و بازرسی مستقیم الگوهای مکانیزم توجه و شبکه‌های پیش‌خور، این تغییر رفتار داخلی را ردیابی کردند. یافته‌های آن‌ها نشان داد که پس از تنظیم دستورالعمل، گراف محاسباتی دچار یک چرخش دانشی می‌شود؛ به این صورت که سرهای توجهِ مدل به‌شکلی تغییر مسیر می‌دهند که تمرکز خود را بیشتر معطوف به افعال دستوری (نظیر «خلاصه کن» یا «ترجمه کن») کنند. هم‌زمان، شبکه‌های پیش‌خور نیز توزیع بازنمایی‌های خود را بازآرایی می‌کنند تا دانش آموخته‌شده‌ی پیشین، در خدمت وظایف کاربردمحور قرار گیرد. با این حال، تحلیل‌های ارائه‌شده در این پژوهش غالباً در سطح ویژگی‌های نحوی و خُرد (نظیر افعال منفرد) باقی می‌ماند، در حالی که دستورالعمل‌های انسانی و استدلال‌های پیچیده، ماهیتی کلان و گزاره‌ای دارند که تقلیل آن‌ها به ردیابی چند فعل در سرهای توجه، نمایانگر تمام ظرفیت استنتاجی مدل نیست.

علاوه بر درک مکانیسم‌های یادگیری، کشف مدارها به‌عنوان ابزاری برای تأمین امنیت و عیب‌یابی ریشه‌ای در شبکه‌های عصبی نیز به‌کار گرفته شده است. گارسیا-کاراسکو و همکاران (۲۰۲۴)~\cite{ijcai2024detecting} با تلفیق پارادایم تفسیرپذیری مکانیکی و حملات متخاصم، چارچوبی فعالانه برای شناسایی آسیب‌پذیری‌های مدل ارائه دادند. متدولوژی آن‌ها بر یک خط‌لوله سه‌مرحله‌ای استوار است؛ نخست، استخراج زیرمجموعه‌ای از اجزا (یک مدار مشخص) که مسئول انجام یک وظیفه‌ی خاص هستند؛ دوم، هدایت فرایند تولید نمونه‌های متخاصم با تمرکز بر فریب همان مدار استخراج‌شده؛ و سوم، تحلیل رفتار داخلیِ مدار در مواجهه با نمونه‌ی مخرب برای کشف نقطه‌ی دقیق شکست. اگرچه این رویکرد گامی مهم در جهت عیب‌یابی ساختاری محسوب می‌شود، ارزیابی آن محدودیت مقیاس‌پذیری رویکردهای مکانیکی را به وضوح نمایان می‌سازد. نویسندگان، این چارچوب را منحصراً بر روی مدل \lr{GPT-2-Small}~\cite{radford2019language} پیاده‌سازی کرده‌اند، زیرا استخراج کامل مدارها و تحلیل یال‌به‌یال آن‌ها برای تولید نمونه‌های متخاصم در مدل‌های ده‌ها میلیارد پارامتری امروزی، نیازمند توان پردازشی نجومی است؛ گلوگاهی که مانع از عملیاتی شدن این تکنیک‌ها در ابعاد صنعتی می‌گردد.

\paragraph{خودرمزگذارهای خلوت برای انتساب مؤلفه‌ها.}

همان‌طور که پیش‌تر بررسی شد، یادگیری دیکشنری و روش‌هایی نظیر \lr{SPINE} به‌منظور باز کردن گره‌ی تعبیه‌های متراکم و رسیدن به بازنمایی‌های خلوت معرفی شدند. با ورود به پارادایم مکانیکی و تحلیل گراف‌های محاسباتی، پژوهشگران دریافتند که این درهم‌تنیدگی در لایه‌های میانیِ مدل به‌دلیل پدیده‌ای بنیادین به‌نام برهم‌نهی\enfootnote{superposition} رخ می‌دهد. مدل‌های زبانی تمایل دارند مفاهیم بسیار بیشتری را نسبت به تعداد نورون‌های خود یاد بگیرند؛ در نتیجه، شبکه‌ها دچار چندمعنایی می‌شوند و یک نورونِ منفرد ممکن است هم‌زمان برای مفاهیمِ کاملاً بی‌ربط (مثلاً برای مفهوم «ریاضیات» و «نام یک شهر») فعال شود.

برای غلبه بر این پدیده و استخراج ویژگی‌های تک‌معنایی\enfootnote{monosemantic} از درون جریان باقیمانده، بریکن و همکاران (۲۰۲۳)~\cite{bricken2023monosemanticity} استفاده از خودرمزگذارهای خلوت (\lr{SAE}) را برای مهندسی معکوسِ فعال‌سازی‌های داخلی ترنسفورمرها در مقیاس وسیع پایه‌گذاری کردند. در این متدولوژی، یک \lr{SAE} با دارا بودن یک لایه‌ی پنهانِ فراتکمیل\enfootnote{overcomplete} و بسیار گسترده‌تر از بُعد مدل اصلی آموزش می‌بیند تا بردار فعال‌سازی متراکم ($x \in \mathbb{R}^d$) را به یک بازنمایی خلوت در ابعاد بالا ($f \in \mathbb{R}^m$ که در آن $m \gg d$) نگاشت کرده و سپس آن را به‌صورت خطی بازسازی نماید ($\hat{x} \in \mathbb{R}^d$). فرایند بهینه‌سازی از طریق کمینه‌سازی خطای بازسازی به‌همراه یک جریمه‌ی خلوت‌سازی (عموماً نُرم $L_1$) انجام می‌پذیرد:
\begin{equation}\label{eq:chap3_sae}
    \mathcal{L} = ||x - \hat{x}||_2^2 + \lambda ||f||_1
\end{equation}
در این معادله، نگاشت رمزگذار به‌صورت $f = \text{\lr{ReLU}}(W_{\text{\lr{enc}}} (x - b_{\text{\lr{dec}}}) + b_{\text{\lr{enc}}})$ و بازسازی به فرم $\hat{x} = W_{\text{\lr{dec}}} f + b_{\text{\lr{dec}}}$ تعریف می‌شود که در آن ستون‌های ماتریس $W_{\text{\lr{dec}}}$ دیکشنریِ جهت‌های مفهومی را تشکیل می‌دهند. به بیان شهودی، تابع زیان در معادله‌ی~\ref{eq:chap3_sae} حاصل تعامل دو نیروی متقابل است: جمله‌ی خطای بازسازی ($\|x - \hat{x}\|_2^2$) وفاداری اطلاعات را تضمین می‌کند تا بازنمایی اصلی مخدوش نشود، در حالی که جمله‌ی جریمه‌ی خلوتی ($\lambda \|f\|_1$) به‌عنوان یک فیلتر سخت‌گیرانه عمل می‌کند؛ نُرم $L_1$ شبکه را مجبور می‌سازد که بردار فعال‌سازی چندمعنایی و متراکم $x$ را تنها با استفاده از ترکیب خطی تعداد بسیار اندکی از ویژگی‌های پنهان ($f_i > 0$) از میان صدها هزار یا میلیون‌ها جهت در دیکشنری بازسازی کند. ضریب $\lambda$ موازنه‌ی میان دقت بازسازی و درجه‌ی تک‌معنایی بودن بردارها را کنترل می‌نماید. کانینگهام و همکاران (۲۰۲۴)~\cite{huben2024sparse} با ارزیابی گسترده‌ی این ساختارها در مدل‌های زبانی اثبات کردند که \lr{SAE}ها با موفقیت توانسته‌اند گراف درهم‌تنیده‌ی شبکه را به مجموعه‌ای از ویژگی‌های به‌شدت تفسیرپذیر تجزیه کنند که هر کدام نمایانگر یک مفهوم انسانی مستقل و شفاف است. 

موفقیت استخراج مکانیکیِ ویژگی‌های تک‌معنایی به‌سرعت به مدل‌های مرزی نیز بسط یافت. گائو و همکاران (۲۰۲۵)~\cite{gao2025scaling} از شرکت \lr{OpenAI}، با اعمال \lr{SAE}ها بر روی مدل عظیم \lr{GPT-4}~\cite{openai2023gpt4}، این معماری را تا ۱۶ میلیون متغیر پنهان مقیاس‌دهی کردند. آن‌ها با استفاده از تکنیک‌هایی نظیر $k$-خلوت برای کنترل مستقیم خلوتی و جلوگیری از فروپاشی یادگیری (مشکل نورون‌های مرده)، ثابت کردند که یادگیری دیکشنری به‌صورت مکانیکی، حتی در پیچیده‌ترین شبکه‌های زبانی امروزی نیز قادر به آشکارسازی ساختار ارگانیک مفاهیم است. این منطق مکانیکی از دامنه‌ی زبان فراتر رفته و در مدل‌های مولد بینایی نیز به‌کار گرفته شد. تیناز و همکاران (۲۰۲۶)~\cite{tinaz2026emergence} با استخراج الگوهای فعال‌سازی داخلی از شبکه‌ی \lr{U-Net} در مدل‌های انتشاری و آموزش \lr{SAE}ها بر روی آن‌ها، نشان دادند که بازنمایی‌های پیچیده‌ی تصویری به ویژگی‌های قابل‌فهم (مانند حضور یک شیء یا سبک هنری) تجزیه می‌شوند. آن‌ها با مداخله در فعال‌سازی این ویژگی‌ها طی فرایند تولید تصویر، علیت آن‌ها را به‌طور تجربی اثبات نمودند.

با وجود دستاوردهای \lr{SAE}ها در یافتن ویژگی‌های تک‌معنایی محلی، بسیاری از مفاهیمِ معناییِ سطح بالا از تن دادن به این بومی‌سازی امتناع می‌ورزند و به‌صورت توزیع‌شده در سراسر شبکه رمزگذاری می‌شوند. هوانگ و همکاران (۲۰۲۴)~\cite{huang2024ravel} با معرفی مجموعه‌داده‌ی \lr{RAVEL}، نشان دادند که ابزارهای مکانیکی سنتی (نظیر کاوشگرها) در تفکیک صفات از مقادیر در این بازنمایی‌های توزیع‌شده ناکارآمدند. آن‌ها برای حل این مشکل، چارچوب پساعملیِ «جستجوی هم‌ترازی توزیع‌شده چندوظیفه‌ای»\enfootnote{multitask distributed alignment search} (\lr{MDAS}) را پیشنهاد دادند. این روش از طریق مداخلات علیِ ساختاریافته در گراف محاسباتی، جهت‌های توزیع‌شده‌ای را کشف می‌کند که دارای پیوند علّی مستحکم با یک مفهوم پیچیده هستند و بدین‌ترتیب ضعف رویکردهای متکی بر بومی‌سازی مطلق را پوشش می‌دهد.

در نهایت، تلاش برای مکان‌یابی و درک مفاهیم در پارادایم مکانیکی، از سطح تحلیل فعال‌سازی‌ها به سطح ریشه‌ای‌تر پارامترها ارتقا یافت. نوین و همکاران (۲۰۲۵)~\cite{nguyen2025unveiling} با ارائه‌ی چارچوب «انتساب مؤلفه»\enfootnote{component attribution} (\lr{CAD}) در مدل‌های انتشاری، مستقیماً بررسی کردند که کدام لایه‌ها و وزن‌های شبکه دارای تأثیر مثبت (ایجادکننده) و کدام‌یک دارای تأثیر منفی (بازدارنده) بر یک مفهوم خاص هستند. متدولوژی \lr{CAD} با مکان‌یابی دقیق در سطح ماتریس وزن‌ها، امکان ویرایش مستقیم مدل را از طریق الگوریتم‌هایی نظیر حذف یا تقویت مؤلفه‌ها در زمان استنتاج فراهم آورد. 

اگرچه خودرمزگذارهای خلوت و روش‌های مکان‌یابی توزیع‌شده یا پارامتری توانستند ابعاد بی‌سابقه‌ای از گراف محاسباتی را مهندسی معکوس کنند، ماهیت تمام این عملیات‌ها مستلزم جستجوی تکرارشونده در ابعاد نجومی شبکه، محاسبه‌ی بازنمایی‌های میانی و دستکاری مستقیم در ساختار مدل است؛ رویکردی که چالش‌های سنگین محاسباتی را با خود به‌همراه دارد.

\paragraph{خودتفسیری و ترجمه‌ی بازنمایی‌های پنهان به زبان طبیعی.}

علی‌رغم موفقیت روش‌هایی نظیر خودرمزگذارهای خلوت در تجزیه‌ی گراف محاسباتی، چالش شکاف معنایی همچنان به قوت خود باقی بود؛ به این معنا که تخصیص یک برچسب زبان طبیعی به بردارهای کشف‌شده، غالباً نیازمند دخالت انسانی یا اتکا به مدل‌های ارزیاب خارجی است. برای غلبه بر این چالش، جریان جدیدی در تفسیرپذیری مکانیکی شکل گرفت که هدف آن ترجمه‌ی حالت‌های پنهان به دایره‌ی واژگان قابل‌درک برای انسان بود.

نخستین ایده‌ی شهودی در این مسیر، تکنیکی به‌نام لنز لوجیت\enfootnote{logit lens} بود که بازنمایی پنهان لایه‌های میانی را مستقیماً در ماتریس تعبیه‌زدایی\enfootnote{unembedding matrix} ضرب می‌کرد تا مشخص شود شبکه در آن لایه، بالاترین احتمال را به کدام کلمات اختصاص داده است. با این حال، این روش در لایه‌های اولیه و میانی به‌شدت دچار خطا می‌شد، زیرا این لایه‌ها به‌دلیل وجود ابعاد دورافتاده با واریانس بالا و پدیده‌ی رانش تدریجیِ کوواریانس بازنمایی‌ها، هنوز با فضای واژگان خروجی هم‌تراز نشده بودند. برای حل این مشکل بنیادین، بلروز و همکاران (۲۰۲۳)~\cite{belrose2023eliciting} چارچوب «لنز تنظیم‌شده»\enfootnote{tuned lens} را معرفی کردند. متدولوژی آن‌ها بر این اصل استوار است که می‌توان برای هر لایه‌ی پنهان ($l$)، یک تبدیل آفین\enfootnote{affine transformation} اختصاصی آموزش داد که نقش یک لنز مترجم را ایفا کند. از نظر ریاضی، اگر $h_l$ بازنمایی پنهان در لایه‌ی $l$ و $W_U$ ماتریس خروجی باشد، توزیع احتمال کلمات پیش‌بینی‌شده ($\hat{y}$) با استفاده از ماتریس وزن لنز ($A_l$) و بردار سوگیری ($b_l$) به‌صورت $\hat{y} = \text{\lr{Softmax}}(W_U (A_l h_l + b_l))$ تقریب زده می‌شود. به بیان شهودی، این تبدیل آفین نقش یک میان‌بُر محاسباتی را بازی می‌کند؛ ماتریس $A_l$ با چرخش و تغییر مقیاس بازنمایی لایه $l$، اثر رانش بازنمایی را خنثی کرده و آن را با هندسه‌ی ورودیِ لایه‌ی نهایی هم‌راستا می‌سازد، و بردار سوگیری $b_l$ میانگین تغییرات تجمعی حاصل از لایه‌های باقیمانده را تقریب می‌زند. این لنزها با شبیه‌سازی محاسبات لایه‌های آتی، نشان می‌دهند که مدل در هر عمق از شبکه‌ی خود، با چه درجه‌ای از قطعیت به خروجی نهایی هم‌گرا شده است.

اگرچه لنز تنظیم‌شده ابزاری قدرتمند برای رمزگشایی از لایه‌های اولیه بود، آموزش ده‌ها لنز مجزا برای یک مدل بزرگ، فرایندی پرهزینه و وابسته به داده‌های آموزشی بود. در پاسخ به این محدودیت، قندهاریون و همکاران (۲۰۲۴)~\cite{ghandeharioun2024patchscopes} با معرفی چارچوب \lr{Patchscopes}، رویکرد ترجمه را از اتکا به کاوشگرها و لنزهای خارجی، به استفاده از ظرفیت تولیدی خود مدل زبانی ارتقا دادند. متدولوژی آن‌ها مبتنی بر تکنیک وصله‌زنی است. در این روش، ابتدا مدل یک پرامپت مبدأ را پردازش می‌کند و بازنمایی پنهان آن در لایه‌ی $l$ و موقعیت $i$ (یعنی $h_i^{(l)}$) استخراج می‌شود. سپس، این بردار مستقیماً به گراف محاسباتی یک پرامپت مقصد (نظیر «پایتخت کشور X برابر است با») تزریق شده و جایگزین فعال‌سازی متناظر در آن می‌گردد. در نهایت، متنی که مدل در ادامه‌ی پرامپت مقصد تولید می‌کند (مثلاً کلمه «پاریس»)، به‌عنوان توضیح زبان‌طبیعیِ محتوای آن حالت پنهان در نظر گرفته می‌شود. این چارچوبِ بدون‌نیاز به آموزش، نه‌تنها توانست روش‌های پیشین تفسیرپذیری را به‌عنوان حالت‌های خاصی از خود بازتولید کند، بلکه نشان داد بازنمایی‌ها در لایه‌های اولیه مبهم بوده و در لایه‌های عمیق‌تر معنای دقیقی به خود می‌گیرند. با این وجود، چالش \lr{Patchscopes} در طراحی پرامپت مقصد نهفته است؛ تداوم محاسبات در گراف شبکه پس از عمل وصله‌زنی می‌تواند اطلاعات جدیدی را از بافتار پرامپت مقصد به بردار پنهان اولیه القا کرده و خلوص توضیح نهایی را مخدوش سازد.

به‌طور هم‌زمان و در رویکردی بسیار مشابه، چن و همکاران (۲۰۲۴)~\cite{chen2024selfie} چارچوب \lr{SelfIE}\enfootnote{self-interpretation of embeddings} را توسعه دادند. این پژوهش نیز دقیقاً بر استخراج تعبیه‌ی پنهان و جاسازی آن درون یک پرامپت پرسشی (مانند «لطفاً پیام قبلی را خلاصه کن») استوار است تا \lr{LLM} وادار گردد محتوای معناییِ آن تعبیه را تشریح کند. با این حال، \lr{SelfIE} برای مقابله با مشکل آلودگی اطلاعات در طول عبور پیش‌رو (نقطه‌ی ضعف \lr{Patchscopes})، یک امتیاز ارتباط معرفی می‌کند. این معیار با مقایسه‌ی احتمال تولید هر توکن در دو حالتِ با و بدون تزریق تعبیه، اثر علّی خالص بردار پنهان را از تأثیرات پرامپت مقصد تفکیک و برجسته می‌سازد. علاوه بر این ظرافت ارزیابی، نوآوری متمایزکننده‌ی \lr{SelfIE} بسط دادن این تفسیرپذیری به فاز کنترل است. در این چارچوب، با استفاده از الگوریتم‌های بهینه‌سازی و محاسبه‌ی گرادیان‌ها (نظیر کمینه‌سازی خطای میانگین مربعات بردارها یا استفاده از مدل‌های ارزیاب زبانی برای تخصیص پاداش)، تعبیه‌های داخلی ویرایش می‌شوند تا رفتار مدل در تولید محتوای نامطلوب نظیر آسیب‌پذیری در برابر تزریق پرامپت یا بروز دانش مخرب به‌طور ریشه‌ای اصلاح گردد.

تلاش برای ترجمه‌ی فضای پنهان به زبان طبیعی، از تفسیر حالت‌های ایستا فراتر رفته و به تفسیر دینامیک استدلال نیز تعمیم یافت. چن و همکاران (۲۰۲۶)~\cite{chen2026latentdebate} با معرفی چارچوب مناظره پنهان\enfootnote{latent debate}، نشان دادند که فرایند تصمیم‌گیری مدل در طول یک استنتاج واحد، غالباً شامل کشمکش میان سیگنال‌های متناقض است. متدولوژی آن‌ها بر استخراج سیگنال‌های حمایت‌کننده و حمله‌کننده از لایه‌های میانی استوار است. آن‌ها این تقابل داخلی را به‌صورت یک مناظره در یک مدل جانشین فرمول‌بندی کردند و نشان دادند که حجم بالای تضاد این سیگنال‌های پنهان، دارای همبستگی شدیدی با بروز توهم در خروجی نهایی مدل است.

اگرچه استفاده از خود مدل به‌عنوان مفسر خویش، راهکاری ظریف برای دور زدن نیاز به برچسب‌گذاری انسانی است، اما این رویکرد از منظر تئوریک با محدودیت‌های بنیادینی روبه‌روست. تزریق و وصله‌زنی یک حالت پنهان از یک گراف محاسباتی (مبدأ) به یک گراف محاسباتی دیگر (مقصد)، موجب بروز یک شوک توزیعی می‌شود. در این حالت، از آنجا که بردار تزریق‌شده با بافتار جدید هم‌خوانی ارگانیک ندارد، توضیح تولیدشده توسط مدل ممکن است صرفاً یک توهم ناشی از قرارگیری در فضای خارج از توزیع (\lr{OOD}) باشد، نه بازتاب وفادارانه‌ی دانش موجود در آن لایه. افزون بر این، اتکای مطلق به خروجی مدل برای ارزیابی درونیات خودش، نیازمند پذیرش این پیش‌فرض است که مدل همواره مفسری وفادار و صادق برای بازنمایی‌های خود است؛ پیش‌فرضی که مقدمه‌ساز جدی‌ترین و کشنده‌ترین نقدهای وارد بر کل پارادایم مکانیکی گردید.

\paragraph{آسیب‌شناسی مداخلات و توهم تفسیرپذیری در مهندسی معکوس.}

علی‌رغم ظرافت تئوریک در استخراج مدارها و ترجمه‌ی حالت‌های پنهان، تمامی این رویکردهای مکانیکی بر یک پیش‌فرض بنیادین استوارند و آن اینکه درک بازنمایی‌های داخلی مستقیماً به درک رفتار خروجی مدل منجر می‌شود. فریدمن و همکاران (۲۰۲۳)~\cite{friedman2023comparing} در پژوهشی انتقادی این فرض پایه‌ای را به چالش کشیدند. آن‌ها با آموزش مدل‌های ترنسفورمر بر روی زبان‌های ساختاریافته‌ی مصنوعی (زبان‌های دایک\enfootnote{Dyck languages})، میزان تطابق میان شباهت بازنمایی (ارزیابی فضاهای فعال‌سازی) و شباهت عملکردی (پیش‌بینی‌های خروجی) را مورد سنجش قرار دادند. یافته‌ی خلاف شهود آن‌ها این بود که این دو معیار لزوماً همسو نیستند؛ مدل‌هایی که دارای فضاهای فعال‌سازی به‌شدت مشابهی هستند، می‌توانند در مواجهه با داده‌های \lr{OOD} رفتارهای تولیدی کاملاً متناقضی از خود نشان دهند. این پدیده‌ی گسست میان شباهت بازنمایی و رفتار عملکردی، هشدار می‌دهد که تحلیل‌های صرفاً متکی بر هندسه‌ی داخلی شبکه می‌تواند محقق را در خصوص منطق عملکردی واقعی مدل و تعمیم‌پذیری آن فریب دهد.

پیامد این قطع ارتباط در فازِ کنترل و دستکاری شبکه نیز خودنمایی می‌کند. یکی دیگر از مفروضات انگیزه‌بخش در مهندسی معکوس این است که کشف مکان ذخیره‌سازی یک مفهوم در شبکه، ما را قادر می‌سازد تا آن ویژگی را بهینه‌تر کنترل یا ویرایش کنیم. هیس و همکاران (۲۰۲۳)~\cite{hase2023does} این فرض را با استفاده از تکنیک‌های ردیابی علّی برای مکان‌یابی حقایق و سپس اعمال روش‌های پیشرفته‌ی ویرایش مدل در همان مکان‌ها مورد آزمون قرار دادند. نتایج آن‌ها نشان داد که عملاً هیچ همبستگی معناداری میان نقشه‌های مکان‌یابی و موفقیت ویرایش وجود ندارد. به‌عبارتی، لایه‌ای که به‌عنوان محل اصلی پردازش یک مفهوم شناسایی می‌شود، غالباً بهترین لایه برای دستکاری آن مفهوم نیست. این قطع ارتباط علّی میان بازنمایی داخلی و مداخله‌ی عملکردی، اثبات می‌کند که تکیه بر تحلیل یال‌ها و پارامترها، راهنمای قابل‌اعتمادی برای هدایت رفتار مدل نیست.

دلیل ریشه‌ای این ناکامی‌ها در ماهیت مخرب مداخلات داخلی نهفته است. تکنیک استاندارد برای ارزیابی علیت در رویکردهای مکانیکی، حذف یا دستکاری مستقیم فعال‌سازی‌ها است. لی و جانسون (۲۰۲۴)~\cite{li2024optimal} با آسیب‌شناسی دقیق این فرایند نشان دادند که در روش‌های مرسوم حذف (مانند صفر کردن یک فعال‌سازی به‌صورت $A(x) \leftarrow 0$ یا قرار دادن آن روی میانگین بی‌قیدوشرط $A(x) \leftarrow \mathbb{E}[A(X)]$ یا بازنمونه‌برداری)، افت شدید عملکرد مدل ناشی از خطای دستکاری نامتعارف و ایجاد سیگنال‌های مصنوعی موسوم به جعل\enfootnote{spoofing} است؛ بدین معنا که این مقادیر قراردادی، همبستگی درونی اجزای شبکه را نادیده گرفته و گراف محاسباتی را به حالت‌های غیرطبیعی و \lr{OOD} سوق می‌دهند، بدون آنکه افت عملکرد صرفاً بازتابی از حذف محتوای اطلاعاتی آن جزء باشد. آن‌ها برای رفع این مشکل، چارچوب حذف بهینه\enfootnote{optimal ablation} را پیشنهاد دادند که در آن مؤلفه با یک مقدار ثابت بهینه ($a^* = \arg\min_a \mathbb{E}_{X \sim \mathcal{D}} [L(M_A(X, a), M(X))]$) جایگزین می‌شود. از نظر تئوریک اثبات گردید که این مقدار، کمترین شکاف زیان ممکن را در میان تمامی روش‌های حذف کلی به‌همراه دارد و سهم جعل را کمینه‌سازی می‌کند.

با این حال، نیاز به چنین فرایند بهینه‌سازی پیچیده‌ای برای جلوگیری از تخریب، آشکار می‌سازد که فضای درونی مدل‌های زبانی به‌شدت شکننده است و هرگونه جراحی ساختاری در بردارهای فعال‌سازی، استنتاج علّی را در معرض خطاهای توزیعی سیستماتیک قرار می‌دهد.

شواهد تجربی قاطع مبنی بر ناکارآمدی عملی این مداخلات مکانیکی، در پژوهش بهالا و همکاران (۲۰۲۵)~\cite{bhalla2025unifying} ارائه گردید. آن‌ها با ارائه‌ی یک چارچوب یکپارچه، قابلیت‌روش‌های پیشرفته‌ی مکانیکی (شامل خودرمزگذارهای خلوت و لنزهای لوجیت) را در تغییر علّی خروجی مدل ارزیابی کردند. یافته‌ی کلیدی آن‌ها کشف پدیده‌ای به‌نام مصالحه‌ی انسجام-مداخله\enfootnote{coherence-intervention trade-off} بود؛ مداخلات مکانیکی در سطح بازنمایی‌های داخلی، به‌طور مداوم و سیستماتیک منجر به تخریب انسجام، روان‌بودن و منطق زبانی متن تولیدشده می‌گردند. شگفت‌انگیزتر آنکه، این پژوهش با قاطعیت اثبات نمود برای کنترل رفتار مدل، دستکاری‌های ساده در سطح ورودی (تکنیک‌های پرامپت‌دهی) عملکردی به‌مراتب امن‌تر، منسجم‌تر و مؤثرتر از مداخلات پیچیده‌ی داخلی نظیر \lr{SAE}ها دارند.

این شواهد تجربی و تئوریک در کنار یکدیگر، محدودیت‌های ذاتی مهندسی معکوس را عیان ساخته و از منظر علمی توجیه می‌کنند که چرا برای درک و کنترل معنادار رفتار مدل‌های زبانی، می‌توان همچنین از جراحی درونی پارامترها فاصله گرفت و به استنتاج‌های علّی در سطح کلان‌تر ورودی-خروجی روی آورد.


\subsubsection{رویکردهای داده‌محور و ویرایش ساختاری}\label{sec:3-1-2-4}

تمامی روش‌های جعبه‌سفیدی که تا به اینجا واکاوی شدند، متمرکز بر تحلیل رفتار مدل در حین پردازش یک ورودی خاص بودند. با این حال، جریان مهم دیگری در پاردایم جعبه‌سفید وجود دارد که صورت‌مسئله را تغییر می‌دهد؛ این رویکردها به‌جای پرسش از چگونگی پردازش ورودی، به دنبال یافتن چرایی بروز یک رفتار از طریق ردیابی ریشه‌های آن در داده‌های آموزشی، و یا تغییر دائمی آن رفتار از طریق دستکاری و تغییر مستقیم وزن‌ها و ساختار مدل پس از اتمام آموزش هستند.

\paragraph{توابع نفوذ و ردیابی داده‌های آموزشی.}

در یک تغییر پارادایم اساسی نسبت به روش‌های انتساب ویژگی، کو و لیانگ (۲۰۱۷)~\cite{koh2017understanding} رویکرد تفسیرپذیری داده‌محور را با استفاده از مفهوم کلاسیکِ توابع نفوذ\enfootnote{influence functions} به شبکه‌های عصبی عمیق معرفی کردند. مسئله‌ی محوری در این رهیافت، پاسخ به این پرسش است که: «کدام داده‌های آموزشی بیشترین تأثیر را در شکل‌گیری این پیش‌بینی خاص داشته‌اند؟». متدولوژی توابع نفوذ به‌جای حذف فیزیکی یک داده و آموزش مجدد و پرهزینه‌ی کل مدل، از یک تقریب ریاضی مبتنی بر مشتقات بهره می‌برد. از نظر ریاضی، تأثیر حذف یک نقطه‌ی آموزشی ($z$) بر روی خطای مدل در یک نقطه‌ی آزمون مشخص ($z_{\text{\lr{test}}}$)، با استفاده از معکوس ماتریس هسین\enfootnote{inverse Hessian matrix} ($H_{\hat{\theta}}^{-1}$) به‌صورت زیر تقریب زده می‌شود:
\begin{equation}\label{eq:chap3_influence_function}
    \mathcal{I}_{\text{\lr{up, loss}}}(z, z_{\text{\lr{test}}}) = -\nabla_{\theta} {L(z_{\text{\lr{test}}}, \hat{\theta})}^T H_{\hat{\theta}}^{-1} \nabla_{\theta} L(z, \hat{\theta})
\end{equation}
در این معادله، $z = (x, y)$ یک نمونه‌ی آموزشی، $z_{\text{\lr{test}}}$ نمونه‌ی آزمون مورد ارزیابی، و $\nabla_{\theta} L(z, \hat{\theta})$ گرادیان زیان نسبت به پارامترهای مدل $\theta$ روی داده‌ی آموزشی است که مشخص می‌کند پارامترها در پاسخ به آن نمونه در چه جهتی حرکت می‌کنند. ماتریس هسین $H_{\hat{\theta}}$ بیانگر انحنای موضعی سطح زیان کل در نقطه‌ی بهینه است و ضرب معکوس آن ($H_{\hat{\theta}}^{-1}$) در گرادیان، اثر تغییر پارامترها را در قالب یک گام نیوتن تقریب می‌زند؛ به‌طوری‌که جهت‌های با انحنای تند مهار شده و دره‌های مسطح فضای بهینه‌سازی برجسته می‌گردند. در نهایت، با ضرب داخلی این بردار جابه‌جایی در گرادیان نقطه‌ی آزمون ($\nabla_{\theta} L(z_{\text{\lr{test}}}, \hat{\theta})$)، میزان هم‌پوشانی اثر آن‌ها سنجیده می‌شود. از منظر شهودی، علامت این رابطه نقشی کلیدی دارد: اگر $\mathcal{I}_{\text{\lr{up, loss}}} < 0$ باشد، نشان می‌دهد که وزن‌دهی به داده‌ی آموزشی $z$ خطای نمونه‌ی آزمون را کاهش می‌دهد، بنابراین $z$ یک حامی یا نمونه‌ی مفید برای استنتاج مدل بوده است؛ در مقابل، مقدار مثبت نشان‌دهنده‌ی یک داده‌ی معارض یا مخرب است که حضور آن به سوگیری و خطای آزمون دامن می‌زند. 

با این وجود، محاسبه و ذخیره‌سازی معکوس ماتریس هسین برای شبکه‌هایی با میلیاردها پارامتر، با پیچیدگی محاسباتی $\mathcal{O}(P^3)$ (که $P$ تعداد پارامترهاست) روبه‌رو است که اجرای آن را در عمل فاجعه‌بار و غیرقابل‌حل می‌سازد. در تلاشی برجسته برای غلبه بر این بن‌بست محاسباتی، ملودوژنیچ و همکاران (۲۰۲۵)~\cite{mlodozeniec2025influence} توابع نفوذ را برای مدل‌های مولد پیشرفته – اختصاصاً مدل‌های انتشاری – مقیاس‌پذیر کردند. متدولوژی آن‌ها برای فرار از محاسبه‌ی مستقیم پرهزینه‌ی هسین، استفاده از تقریب \lr{K-FAC}\enfootnote{Kronecker-factored approximate curvature} بود. این تقریب با فرض استقلال تقریبی میان لایه‌های شبکه و بهره‌گیری از ساختار بلوک-قطریِ ماتریس‌ها، انحنای فضای پارامترها را با استفاده از ضرب کرونکر تخمین می‌زند؛ ترفندی که وارون‌سازی ماتریس را از یک مانع غیرقابل‌حل، به یک عملیات محاسباتی مقیاس‌پذیر تبدیل می‌کند.

یافته‌های این پژوهش نشان داد که با اعمال این تقریب در تابع زیان تطبیق امتیاز، می‌توان با دقت بالایی ردیابی کرد که یک خروجیِ تولیدشده (مثلاً یک تصویر خاص)، دقیقاً از کدام نمونه‌های آموزشی الهام گرفته است و کدام داده‌ها نقش مفید یا مخربی در شکل‌گیری آن خروجی داشته‌اند. با این حال، ارزیابی این دستاورد در بستر مدل‌های زبانی عظیم، محدودیت‌های بازدارنده‌ی پارادایم داده‌محور را بیش از پیش نمایان می‌سازد. نخست، تقریب \lr{K-FAC} اگرچه وارون‌سازی را ممکن می‌کند، اما همچنان نیازمند دسترسی مطلق به گراف محاسباتی و محاسبه‌ی گرادیان‌ها برای تک‌تک لایه‌هاست. دوم، این روش برای ردیابی ریشه‌های یک رفتار، نیازمند اسکن مداوم بر روی کل مجموعه‌داده‌ی آموزشی است؛ نیازی که آن‌ها را از حوزه‌ی کاربردهای سبک و سریع خارج می‌سازد. سوم، این فرمول‌بندی مشخصاً برای فضاهای پیوسته در مدل‌های انتشاری طراحی شده است و تعمیم آن به ماهیت گسسته و خودهمبسته تولید متن در \lr{LLM}ها با چالش‌های نظری عمیقی روبه‌روست که کارایی آن را برای حسابرسی زبان طبیعی به‌شدت محدود می‌سازد.

\paragraph{مداخلات پساآموزشی و ویرایش ساختاری.}

در مقابل روش‌هایی که صرفاً داده‌های مؤثر را ردیابی می‌کنند، جریان دیگری از پژوهشگران تلاش کردند تا با دستکاری، به‌روزرسانی یا ساده‌سازی ساختار و وزن‌های مدل پس از اتمام آموزش، تفسیرپذیری و رفتار استدلالی آن را به‌صورت دائمی اصلاح کنند.

یکی از برجسته‌ترین نمونه‌ها در این زمینه، چارچوب «ویرایش مدل رتبه‌یک»\enfootnote{rank-one model editing} (\lr{ROME}) است که توسط منگ و همکاران (۲۰۲۲)~\cite{meng2022locating} ارائه گردید. مسئله‌ی این پژوهش، تغییر دقیق یک دانش یا حقیقت خاص در شبکه‌ی زبانی (مثلاً تغییر پایتخت فرانسه از پاریس به رم) بدون تخریب یا دست‌خوردگی سایر اطلاعات شبکه بود. متدولوژی \lr{ROME} ابتدا با استفاده از تکنیک ردیابی علّی نشان می‌دهد که لایه‌های پیش‌خور در اعماق میانی شبکه، دقیقاً مانند یک حافظه‌ی تداعی‌گر خطی ($W k \approx v$) عمل می‌کنند؛ جایی که موضوع گزاره به‌عنوان یک کلید ($k$) و ویژگی یا حقیقت مرتبط با آن به‌عنوان یک ارزش ($v$) ذخیره می‌شود. برای ویرایش این حافظه، آن‌ها یک ماتریس اصلاحی رتبه‌-یک ($\Delta$) را محاسبه کرده و به ماتریس وزن‌های پیش‌خور ($W$) اضافه می‌کنند ($W \leftarrow W + \Delta$) تا شبکه برای کلید قدیمی، ارزش جدید ($v_*$) را برگرداند. شهود این به‌روزرسانی در آن است که ماتریس $\Delta$ به‌صورت یک ضرب خارجی میان بردار خطای ارزش و بردار کلید پروژکت‌شده محاسبه می‌شود؛ با به‌کارگیری ماتریس معکوس کوواریانس کلیدها، این دستکاری به‌گونه‌ای تنظیم می‌گردد که کمترین تداخل را با سایر جهت‌های متعامد در حافظه‌ی شبکه ایجاد کند. یافته‌های این پژوهش ثابت کرد که می‌توان حقایق را با موفقیت و پایداری بالا در درون وزن‌ها بازنویسی کرد. با این حال، محدودیت بنیادین \lr{ROME} در سادگی هدف آن است؛ این روش صرفاً برای حقایق صریح سه‌تایی (موضوع-رابطه-مفعول) کارآمد است و برای استدلال‌های پیچیده و دانش‌های ضمنی ناتوان است.

برای غلبه بر محدودیت روش‌های پیشین در مواجهه با استدلال‌های پیچیده، بالاسوبرامانیان و همکاران (۲۰۲۵)~\cite{balasubramanian2025decomposition} چارچوب \lr{DecompTune} را معرفی کردند. مسئله‌ی اصلی آن‌ها یافتن ریشه‌ی تصمیمات در وظایف پرسش‌وپاسخ چندگام بود؛ جایی که پاسخ نهایی حاصل ترکیب اطلاعات از چندین منبع پراکنده است و انتساب یک‌باره‌ی اعتبار به ورودی ناممکن است. متدولوژی آن‌ها مسئله‌ی تفسیرپذیری را از یک محاسبه‌ی صرفاً ریاضی، به یک وظیفه‌ی تولید متن تبدیل می‌کند. ایده‌ی شهودی این است که مدل را آموزش دهند تا پیش از ارائه‌ی پاسخ نهایی، استدلال خود را به گزاره‌های مستقل تجزیه کرده و منبع هر گزاره را در متن مشخص کند.

برای اجرای این ایده، مدل‌های پایه‌ای نظیر \lr{Qwen-2.5}~\cite{qwen2024qwen25} تحت یک خط‌لوله‌ی آموزش پساآموزشی دومرحله‌ای قرار می‌گیرند. در مرحله‌ی اول (تنظیم ظریف نظارت‌شده یا \lr{SFT})، مدل روی داده‌هایی که پیش‌تر توسط یک \lr{LLM} قوی‌تر (به‌عنوان معلم) تجزیه و استدلال‌نویسی شده‌اند، آموزش می‌بیند. در مرحله‌ی دوم، از الگوریتم یادگیری تقویتی \lr{GRPO} استفاده می‌شود؛ به این صورت که مدل برای یک پرسش، چندین مسیر استدلالی مختلف تولید می‌کند. هر مسیر بر اساس میزان صحت تجزیه‌ها و تطابق آن‌ها با متن مرجع، یک نمره‌ی پاداش ($R$) دریافت می‌کند و سپس سیاست تولید مدل ($\pi_\theta$) برای بیشینه‌سازی این پاداش‌ها بهینه‌سازی می‌شود. اگرچه این رویکرد خروجی‌های استدلالی و به‌شدت شفافی تولید می‌کند، اما با دو محدودیت روش‌شناختی مواجه است؛ نخست، وابستگی مطلق به کیفیت داده‌های تولیدشده توسط مدل معلم (که خطر یادگیری و درونی‌سازی توهمات آن مدل را به‌همراه دارد) و دوم، تغییر عامدانه‌ی توزیع اصلی شبکه. در حقیقت، این روش دیگر در حال تفسیر مدل اولیه‌ی خود نیست، بلکه رفتار آن را از طریق آموزش مجدد تغییر داده است تا مجبور به شفافیت شود.

در امتداد تلاش‌ها برای تغییر دائمی مدل، رویکرد دیگری مبتنی بر ساخت مدل‌های جانشین ساده‌شده از طریق تغییرات ساختاری درونی شکل گرفت. فریدمن و همکاران (۲۰۲۴)~\cite{friedman2024interpretability} یکی از بنیادی‌ترین مفروضات در این زمینه را به چالش کشیدند؛ این فرض که اگر بتوانیم شبکه را به یک نسخه‌ی ساده‌تر و شفاف‌تر تبدیل کنیم که رفتار مدل اصلی را تقریب می‌زند، منطق درونی آن را نیز کشف کرده‌ایم. متدولوژی آن‌ها شامل آموزش ترنسفورمرها بر روی زبان‌های ساختاریافته (نظیر زبان‌های پرانتزگذاری متوازن دایک) و سپس اعمال مداخلات ساختاری و ماندگار بود؛ آن‌ها با استفاده از تقلیل ابعاد خطی (نظیر \lr{SVD} بر روی ماتریس وزن‌ها)، خوشه‌بندی، و تبدیل مکانیزم توجه نرم به توجه سخت، مدل اصلی را به یک نسخه‌ی جانشین ساده‌تر تبدیل کردند. یافته‌های آن‌ها از یک شکاف تعمیم پرده برداشت؛ این جانشین‌های ساده‌سازی‌شده در داده‌های درون‌توزیع به‌شدت به مدل اصلی وفادار بودند، اما در سناریوهای \lr{OOD} کاملاً شکست خوردند و رفتاری متفاوت از مدل اصلی نشان دادند. این پدیده که تحت عنوان توهم‌های تفسیرپذیری معرفی شد، هشداری قاطع برای رویکردهای مداخله‌گر و ساده‌ساز است. این ارزیابی اثبات می‌کند که بازطراحی ساختاری و ساده‌سازی‌های ماتریسی، تنها ویژگی‌های سطحی و همبستگی‌های آماری شبکه را حفظ می‌کنند و قادر به استخراج یا حفظ مکانیسمِ علّی و واقعی استنتاج نیستند.


\subsubsection{جمع‌بندی رویکردهای پساعملی جعبه‌سفید}\label{sec:3-1-2-5}

با مرور جامع پارادایم تفسیرپذیری جعبه‌سفید، از انتساب اعتبارِ مبتنی بر سیگنال‌های پایه‌ی گرادیان و توجه (بخش~\ref{sec:3-1-2-1})، تا کشف مفاهیم ارگانیک در فضاهای پنهان (بخش~\ref{sec:3-1-2-2})، واکاوی هندسه‌ی مدارهای محاسباتی داخلی (بخش~\ref{sec:3-1-2-3}) و در نهایت ردیابی توابع نفوذ و ویرایش مستقیمِ وزن‌ها (بخش~\ref{sec:3-1-2-4})، مشخص می‌گردد که این روش‌ها به‌دلیل بهره‌مندی مستقیم از گرادیان‌ها، فعال‌سازی‌های پنهان و ساختار گراف محاسباتی، چارچوب‌هایی قدرتمند و از نظر ریاضی اصیل برای تحلیل رفتار شبکه‌ها فراهم می‌آورند. از همین رو، در شرایطی که دسترسی کامل به مدل فراهم باشد، این متدولوژی‌ها به‌عنوان خط‌پایه‌های مقایسه‌ای بنیادین و بسیار معتبری برای سنجش عملکرد سایر رویکردها عمل می‌کنند. 

با این حال، در کاربردهای عملیاتی مدرن، بسیاری از مدل‌های زبانی بزرگِ مرزی به‌صورت سیستم‌های انحصاری، تجاری و مسدود در پسِ رابط‌های برنامه‌نویسی (\lr{API}) ارائه می‌شوند که در آن‌ها دسترسی به وزن‌ها، گرادیان‌ها یا فعال‌سازی‌های درونی مسدود است. این محدودیت دسترسی، مانع از به‌کارگیری مستقیم روش‌های جعبه‌سفید در ممیزی مدل‌های تجاری بسته می‌گردد و ضرورت پژوهش پیرامون پارادایم تفسیرپذیری جعبه‌سیاه را به‌عنوان رویکردی واقع‌گرایانه در مواجهه با سیستم‌های مسدود ایجاب می‌کند.


\subsection{تفسیرپذیری پساعملی: رویکردهای جعبه‌سیاه}\label{sec:3-1-3}

در پارادایم تفسیرپذیری جعبه‌سیاه، فرض اساسی بر انسداد کامل گراف محاسباتی استوار است~\cite{guidotti2018survey}. در این رویکرد، تحلیل‌گر هیچ‌گونه دسترسی به معماری داخلی، ماتریس وزن‌ها، فضاهای فعال‌سازی یا گرادیان‌های شبکه ندارد و مدل اساساً به‌عنوان یک موجودیت کدر مدل‌سازی می‌شود که تابعی ناشناخته نظیر $f$ را برای نگاشت ورودی $X$ به خروجی $Y$ تقریب می‌زند. با این وجود، میزان دسترسی به خروجیِ این تابع در سیستم‌های مختلف، یک طیف را تشکیل می‌دهد. در حالی که رابط‌های تجاری مسلط امروزی غالباً تعامل را به یک ورودی و خروجی مطلق (برای مثال متن) محدود می‌سازند و به‌عبارتی جعبه‌سیاه مطلق هستند، طیف گسترده‌ای از متدولوژی‌های کلاسیک این حوزه که تحت عنوان روش‌های مستقل از مدل شناخته می‌شوند، برای ارزیابی خروجی نیازمند دسترسی به توزیع احتمالات یا لوجیت‌های شبکه ($P(Y \mid X)$) هستند. بر این اساس، هدف کلان در این پارادایم، استنتاج مکانیزم‌های پنهان مدل صرفاً از طریق دستکاری سیستماتیک متغیرهای ورودی ($X$) و رصد تغییرات و نوسانات در متن خروجی یا توزیع احتمالات آن است.

با توجه به اینکه امروزه سنگین‌ترین و پیچیده‌ترین مسئله‌ی جعبه‌سیاه در حوزه‌ی هوش مصنوعی معطوف به \lr{LLM}ها است، تمرکز اصلی این بخش بر روی چالش‌های مرتبط با تولید متن آزاد خواهد بود~\cite{schneider2024explainable, danilevsky2020survey}. بر این اساس، در ادامه ابتدا به بررسی شهودی‌ترین راهکار برای تفسیر سیستم‌های مدرن (با دسترسی کاملاً مسدود)، یعنی رویکردهای مبتنی بر خودتوضیحی و استخراج متنی پرداخته می‌شود. سپس، روش‌های مستقل از مدل – از جمله اعمال اختلالات سیستماتیک بر روی ورودی (که غالباً نیازمند دسترسی به احتمالات هستند) و آموزش مدل‌های جانشین – که در اصل برای وظایف ساختاریافته‌تری نظیر دسته‌بندی توسعه یافته بودند، مرور شده و چالش‌های عمیق نظری در تطبیق آن‌ها با ماهیت گسسته‌ی زبان طبیعی بررسی می‌گردد. در نهایت، با توجه به فقدان دسترسی به گراف محاسباتی و عدم‌وجود حقیقت پایه در این شبکه‌ها، بحران ارزیابی و چالش‌های اعتبارسنجیِ متدولوژی‌های تفسیری در این پارادایم به‌تفصیل واکاوی خواهد شد.


\subsubsection{رویکردهای مبتنی بر خودتوضیحی و استخراج متنی}\label{sec:3-1-3-1}

در مواجهه با سیستم‌های هوشمند جعبه‌سیاه، یکی از شهودی‌ترین رویکردها، درخواست از خود مدل برای توضیح رفتار خویش است؛ با این وجود، بایستی تأکید کرد که این متدولوژی مشخصاً و منحصراً برای \lr{LLM}ها کاربرد دارد. در این پارادایم که تحت عنوان تفسیرپذیری خودتوضیحی یا مبتنی بر پرامپت شناخته می‌شود، فرض بر این است که \lr{LLM}های مدرن، به‌دلیل برخورداری از ظرفیت‌های عظیمِ بازنمایی زبان و قابلیت‌های منحصربه‌فرد یادگیری درون‌بافتاری، توانایی درون‌نگری و بیان منطق تصمیم‌گیری خود را در قالب زبان طبیعی دارند. از منظر مکانیکی، از آنجا که این مدل‌ها بر پایه‌ی پیش‌بینی توکن بعدی عمل می‌کنند، وادار کردن آن‌ها به تولید گام‌به‌گام استدلال‌ها باعث می‌شود مدل در طول تولید، بافتار غنی‌تری برای خود بسازد. این فرایند نه تنها با شرطی‌سازی احتمالات خروجی بر روی توکن‌های استدلالی پیشین باعث بهبود چشمگیر عملکرد در حل مسئله می‌شود، بلکه تا حدی فرایند استنتاج را شفاف‌سازی کرده تا کاربر بتواند مسیر رسیدن به پاسخ را تفسیر کند. 

\paragraph{تکامل استدلال درون‌بافتاری.}

پارادایم خودتوضیحی با پژوهش جریان‌ساز وی و همکاران (۲۰۲۲)~\cite{wei2022chain} و معرفی ایده‌ی زنجیره‌ی افکار (\lr{CoT}) پایه‌گذاری شد. مسئله‌ی محوری این پژوهش، ناکارآمدی روش‌های استاندارد پرامپت‌دهی در حل مسائل پیچیده‌ی ریاضی و منطقی بود. متدولوژی \lr{CoT} با استفاده از پرامپت‌دهی چندنمونه‌ای\enfootnote{few-shot prompting}، مدل را وادار می‌کرد تا خروجی خود را در فرمتِ «ورودی $\leftarrow$ زنجیره‌ای از گام‌های استدلالی میانی $\leftarrow$ خروجی» تولید کند. با شکستن مسائل به گام‌های کوچکتر، بار محاسباتی مدل توزیع شده و احتمال خطای انباشته کاهش یافت. با این حال، تحلیل روش‌شناختی نشان می‌دهد که خروجی \lr{CoT} اساساً یک دلیل‌تراشی رو به جلو است. مدل صرفاً متنی تولید می‌کند که از نظر توزیع آماری زبان، توجیه منطقی پاسخ نهایی به نظر می‌رسد؛ اما این خروجی به‌هیچ‌وجه یک انتساب علّی رو به عقب نیست و مشخص نمی‌کند کدام پیش‌آیندها در پرامپت کاربر، از نظر محاسباتی باعث فعال‌شدن این مسیر استدلالی در گراف شبکه شده‌اند.

در امتداد این رویکرد، کوجیما و همکاران (۲۰۲۲)~\cite{kojima2022large} اثبات کردند که توانایی خودتوضیحی لزوماً منوط به ارائه‌ی نمونه‌های دست‌ساز انسانی نیست. آن‌ها با معرفی استدلال زنجیره‌ای صفر-نمونه\enfootnote{zero-shot \lr{CoT}} نشان دادند تنها با افزودن عبارات محرک نظیر «بیا گام‌به‌گام فکر کنیم»، مدل وادار به تولید یک زنجیره‌ی استدلالی خام می‌شود که از طرفی باعث افزایش دقت پاسخ در مسائل پیچیده منطقی می‌شود. این کشف، ماهیت عمیقاً آماری و شرطی این استدلال‌ها را عیان‌تر ساخت؛ تولید این توجیهات صرفاً واکنشی به یک الگوی زبانی است و لزوماً بازتاب‌دهنده‌ی یک استنتاج مستحکم در لایه‌های پنهان نیست، امری که مدل را به‌شدت مستعد تولید توهمات متنی می‌کند.

برای غلبه بر شکنندگی و توهم در مسیرهای استدلالی منفرد، ساختارهای پیشرفته‌تری توسعه یافتند. وانگ و همکاران (۲۰۲۳)~\cite{wang2023selfconsistency} تکنیک سازگاری درونی\enfootnote{self-consistency} را معرفی کردند که در آن، به‌جای استخراج حریصانه‌ی یک خروجی، ده‌ها مسیر استدلالی مجزا نمونه‌برداری شده و پاسخ نهایی از طریق رای‌گیری اکثریت استخراج می‌گردد. هم‌زمان، یائو و همکاران (۲۰۲۳)~\cite{yao2023tree} با ارائه‌ی چارچوب «درخت افکار»\enfootnote{tree-of-thoughts} (\lr{ToT})، فرایند را به یک جستجوی درختی (با استفاده از الگوریتم‌هایی نظیر \lr{BFS} یا \lr{DFS}) تبدیل کردند. وانگ و همکاران (۲۰۲۳)~\cite{wang2023planandsolve} نیز در متدولوژی \lr{Plan-and-Solve}، فاز طراحی یک نقشه‌ی کلان را پیش از اجرای محاسبات الزامی نمودند. علی‌رغم دستاوردهای چشمگیر این ساختارها در افزایش دقت و انسجام متن، تمامی آن‌ها ماهیتاً الگوریتم‌هایی برای بهینه‌سازی جستجو هستند و در ذات خود مکانیزمی بیرونی برای سنجش عملکرد و مهار واقعی توهمات استدلالی ندارند.

در پاسخ به این محدودیت، جریان تحقیقاتی جدیدی شکل گرفت که تلاش می‌کرد خود مدل را به‌عنوان یک ناظر بر خروجی‌هایش به‌کار گیرد. مدان و همکاران (۲۰۲۳)~\cite{madaan2023selfrefine} سیستم \lr{Self-Refine} را پیشنهاد دادند که در آن مدل به‌صورت چرخه‌ای نقش‌های تولیدکننده، منتقد و پالایش‌گر را برای یافتن و اصلاح خطاهای خود ایفا می‌کند. دهولیاوالا و همکاران (۲۰۲۴)~\cite{dhuliawala2024chain} نیز زنجیره‌ی اعتبارسنجی\enfootnote{chain-of-verification} (\lr{CoVe}) را توسعه دادند تا مدل پیش از ارائه‌ی پاسخ نهایی، سؤالات مستقلی را برای راستی‌آزمایی ادعاهای خود طرح و پاسخ دهد. همچنین، وانگ و ژائو (۲۰۲۴)~\cite{wang2024metacognitive} با پرامپت‌دهی فراشناختی\enfootnote{metacognitive prompting}، مدل را هدایت کردند تا میزان اطمینان و شکاف‌های دانشی خود را طی یک مونولوگ متنی بسنجد. موفقیت نسبی این روش‌ها در اصلاح خطاهای سطحی، این فرض خوش‌بینانه را تقویت کرد که سیستم‌های مولد می‌توانند نقش مفسر و داوری بی‌طرف را برای منطق خود ایفا کنند.

با این وجود، اتکا به این مونولوگ‌های متنی به‌عنوان توضیحات وفادارانه، در عمل با شکست‌های تجربی عمیقی مواجه گردید، زیرا این رویکردها فاقد یک معیار مستقل و بیرونی برای اعتبارسنجی گراف استدلال هستند. افزون بر این، باید توجه داشت که ماهیت مسئله‌ای که این چارچوب‌ها به دنبال حل آن‌اند ذاتاً متفاوت است؛ تمرکز آن‌ها صرفاً بر شفاف‌سازی فرایند استنتاج (تولید رو به جلو) است و اساساً مکانیزمی برای استخراج اهمیت ویژگی‌ها یا کشف وابستگی‌های علّی در ورودی ارائه نمی‌دهند. در ادامه، موارد مذکور به‌تفصیل بررسی می‌شوند.

\paragraph{بحران وفاداری و توهم استدلالی.}

اعتبار پارادایم خودتوضیحی و اعتماد به ناظر درونی، با واکاوی‌های تجربی پیرامون پدیده‌ی وفاداری در هم شکست. در ادبیات تفسیرپذیری، یک توضیح تنها زمانی وفادار تلقی می‌شود که به‌درستی و صداقت، بازتاب‌دهنده‌ی فرایند محاسباتی و مکانیسم علّی واقعی در درون مدل برای رسیدن به یک خروجی خاص باشد. تورپین و همکاران (۲۰۲۳)~\cite{turpin2023language} در یک مطالعه‌ی انتقادی نشان دادند که استدلال‌های متنی تولیدشده توسط مدل، غالباً پوششی فریبنده برای مکانیزم‌های تصمیم‌گیری واقعی آن هستند. آن‌ها با تزریق سوگیری‌های نامحسوس به پرامپت (نظیر مرتب‌سازی گزینه‌ها به‌گونه‌ای که پاسخ همواره گزینه‌ی الف باشد)، متوجه شدند که شبکه در لایه‌های پنهان خود دقیقاً تحت تأثیر همین سوگیری‌ها تصمیم‌گیری می‌کند؛ اما در خروجی متنی، هرگز به این علل واقعی اشاره نمی‌کند. در عوض، شبکه دچار استدلال هدفمند شده و برای توجیه تصمیمی که از پیش بر اساس سوگیری گرفته، یک زنجیره‌ی استدلالی ظاهراً منطقی اما کاملاً نامعتبر تولید می‌نماید.

این عدم‌وفاداری، علاوه بر خدشه‌دار کردن اعتبار تفسیرپذیری، پیامدهای اجتماعی خطرناکی نیز به‌همراه دارد. شیخ و همکاران (۲۰۲۳)~\cite{shaikh2023onsecondthought} با ارزیابی استدلال زنجیره‌ای نشان دادند که اعطای آزادی به یک مدل جعبه‌سیاه برای تولید گام‌به‌گام تفکراتش، به‌طور معناداری احتمال تولید خروجی‌های سمی و سوگیری‌دار را افزایش می‌دهد. آن‌ها با مقایسه‌ی پرامپت‌دهی استاندارد و صفر-نمونه بر روی مجموعه‌داده‌های حساس، مشاهده کردند که الزام مدل به تولید مسیر استدلالی، باعث می‌شود شبکه به‌صورت خودکار مفروضات و کلیشه‌های مضر را فراخوانی کرده و از آن‌ها برای نتیجه‌گیری‌های سمی استفاده کند. در واقع، فضای متنیِ اختصاص‌یافته برای توجیهات، مانند یک کاتالیزور عمل کرده و پیش‌فرض‌های نامطلوب پنهان در داده‌های آموزشی را تشدید می‌کند.

بروز این رفتارهای فریبنده و سمی، پژوهشگران را با یک پرسش اساسی مواجه ساخت؛ آیا مدل‌ها اساساً ظرفیت آن را دارند که (همان‌طور که رویکردهای مبتنی بر پرامپت نظیر \lr{Self-Refine} ادعا می‌کردند) خروجی‌ها و استدلال‌های خود را به‌عنوان یک منتقد ارزیابی و اصلاح کنند؟ پاسخ نظری و قاطع به این پرسش در مطالعه‌ی شارما و همکاران (۲۰۲۴)~\cite{sharma2024acritical} پیرامون خط‌لوله‌های یادگیری از بازخورد هوش مصنوعی\enfootnote{learning from AI feedback} (\lr{LAIF}) ارائه شد. این پژوهشگران برای بررسی این ظرفیت، از یک مدل زبانی قدرتمند (به‌عنوان معلم) برای هم‌ترازی یک مدل پایه‌ی کوچکتر استفاده کرده و دو سناریو را با هم مقایسه نمودند؛ سناریوی نخست که در آن مدل معلم نقش ارزیاب و منتقد را ایفا کرده و با مقایسه‌ی خروجی‌ها، داده‌های ترجیحی لازم برای یادگیری تقویتی را تولید می‌کند (\lr{RLAIF})، و سناریوی دوم که مدل معلم صرفاً نقش تولیدکننده را داشته و مدل پایه مستقیماً بر روی پاسخ‌های ایده‌آل آن تنظیم ظریف می‌شود (\lr{SFT}). 

نتایج این آزمایش‌های کنترل‌شده نشان داد که عملکرد نهایی مدل پایه در رویکرد مبتنی بر منتقد، برتری معناداری نسبت به رویکرد تولید مستقیم ندارد. این یافته‌ی تجربی، خط بطلانی بر فرضیه‌ی شکاف تولیدکننده-ارزیاب\enfootnote{generator-evaluator gap} کشید و اثبات کرد که توانایی انتقادی و ارزیابانه‌ی حتی پیشرفته‌ترین مدل‌ها، از توانایی تولیدی آن‌ها فراتر نمی‌رود. به بیان دیگر، تخصیص نقش منتقد به یک مدل (چه برای ارزیابی خودش و چه مدل‌های دیگر)، معجزه‌ای در درک منطقی آن ایجاد نمی‌کند؛ مدل در جایگاه داور نیز درگیر همان سوگیری‌ها، توهمات و نقاط کوری است که در زمان تولید متن به آن‌ها مبتلاست. این واقعیت ساختاری ثابت می‌کند که سیستم‌های زبانی فاقد یک توانایی مستقل برای اعتبارسنجی خروجی‌ها هستند و واگذاری وظیفه‌ی کشف حقیقت به مکانیزم‌های خودارزیابی، رویکردی ذاتاً نامعتبر است.

با درک این بن‌بست که مدل از ارزیابی درونی خود ناتوان است، تلاش‌ها به سمت استفاده از معیارهای مستقل و بیرونی معطوف گردید. برای نمونه، یئو و همکاران (۲۰۲۴)~\cite{weijie2024howinterpretable} با ارزیابی شکنندگی توضیحات مبتنی بر پرامپت، چارچوب \lr{SEA-CoT} را پیشنهاد دادند. در این متدولوژی، مدل چندین استدلال میانی و پاسخ نهایی تولید می‌کند، اما سپس با استفاده از فیلترهای استلزام متنی (که توسط یک سیستم یا معیار مستقل خارجی بررسی می‌شود)، تنها مسیرهایی پذیرفته می‌شوند که از نظر منطق صوری فاقد تناقض باشند. با وجود ارتقای چشمگیر انسجام ظاهری متن با تکیه بر این قیدهای خارجی، یک محدودیت بنیادین نمایان می‌گردد؛ تضمین عدم‌تناقض یک متن با خودش، هرگز معادل با اثبات وفاداری آن استدلال به گراف محاسباتی شبکه نیست و چالش یافتن پیش‌نیازهای علّی واقعی را در یک جعبه‌سیاه برطرف نمی‌سازد.

\paragraph{تقابل ماهوی بهینه‌سازی هدف‌محور و انتساب پساعملی رفتار.}

با توسعه‌ی ابزارهای تعامل با مدل‌های جعبه‌سیاه، جریان دیگری از پژوهش‌ها شکل گرفت که مرز تئوریک میان مهندسی پرامپت و تفسیرپذیری را کم‌رنگ کرد. روش‌هایی نظیر چارچوب مهندسی پرامپت خودکار (\lr{APE})~\cite{zhou2023large} برای یافتن دستورالعمل‌هایی با بالاترین دقت، و چارچوب \lr{PreciseDebias}~\cite{clemmer2024precisedebias} برای هدایت توزیع جمعیتی در مدل‌های مولد تصویر توسعه یافتند. اگرچه این روش‌ها ارتباط میان تغییرات ورودی و خروجی را فرمول‌بندی می‌کنند، اما تحلیل آن‌ها نشان‌دهنده‌ی یک تفاوت ماهوی است؛ این چارچوب‌ها استراتژی‌هایی برای هدایت مدل و بهینه‌سازی رو به جلو هستند. یافتن یک پرامپت بهینه برای رسیدن به یک هدف صریح، معادل با کشف علیت نیست و پاسخی به این پرسش نمی‌دهد که چرا یک پرامپت دست‌نخورده، منجر به تولید یک خروجی مبهم در مدل شده است. از طرفی وفق این روش‌ها برای این وظیفه جدید دچار چالش‌های ماهوی است.

این مسئله در رویکردهای مبتنی بر حملات متخاصم نظیر الگوریتم جستجوی گرادیان مختصات حریصانه\enfootnote{greedy coordinate gradient} (\lr{GCG})~\cite{zou2023universal} با وضوح بیشتری آشکار می‌گردد. موفقیت الگوریتم‌های جستجوی متخاصم بر پایه‌ی کمینه‌سازی تابع زیان برای یک توالی هدفِ صریح و ازپیش‌مشخص‌شده ($y_{\text{\lr{target}}}$) استوار است. به بیانی دیگر:
\begin{equation}\label{eq:chap3_gcg}
    \min_{x_{\text{\lr{adv}}}} -\log P(y_{\text{\lr{target}}} \mid x_{\text{\lr{prompt}}}, x_{\text{\lr{adv}}})
\end{equation}
در این معادله، $x_{\text{\lr{prompt}}}$ دستورالعمل اولیه‌ی کاربر و $x_{\text{\lr{adv}}}$ یک پسوند متنی (دنباله‌ای از توکن‌ها) است که الگوریتم به دنبال کشف آن است. تابع هدف در اینجا کمینه‌سازی لگاریتم منفی احتمال (همان زیان آنتروپی متقاطع استاندارد) است که از منظر شهودی و ریاضی، دقیقاً با بیشینه‌سازی احتمال شرطی تولید رشته‌ی هدف ($P(y_{\text{\lr{target}}} \mid \dots)$) برابری می‌کند. الگوریتم با ثابت نگه داشتن پرامپت اولیه‌ی کاربر ($x_{\text{\lr{prompt}}}$)، فضای گسسته‌ی توکن‌ها را به کمک سیگنال گرادیان جستجو می‌کند تا کلمات پسوند متخاصم ($x_{\text{\lr{adv}}}$) را به‌شکلی بازنویسی کند که عبارت هدف ($y_{\text{\lr{target}}}$، نظیر جمله‌ی تأییدی «بله، آموزش ساخت بمب به این شرح است») با بالاترین احتمال ممکن به‌عنوان ادامه‌ی متن توسط مدل انتخاب شود.

این مکانیزم صرفاً یک نمونه از قاعده‌ای ساختاری در الگوریتم‌های جستجو است؛ همان‌طور که در سایر روش‌های خودکارسازی نظیر \lr{APE}~\cite{zhou2023large} نیز صادق است، فرایند بهینه‌سازی در فضای ورودی ماهیتاً نیازمند یک هدف قطعی و ازپیش‌تعیین‌شده برای ارزیابی تابع امتیاز یا تقریب گرادیان است. در مقابل، مسئله‌ی انتساب اثرگذاری در استنتاج، ذاتاً فاقد چنین حقیقت پایه‌ای برای سنجش وفاداری است. هدف در تفسیرپذیری بهینه‌سازی یک خروجی جدید نیست، بلکه تشخیص علت تولید یک پاسخ خودکامه و ازپیش‌تولیدشده است. در غیاب یک توالی هدف مشخص، نمی‌توان پرامپت‌ها را مستقیماً برای کشف علیت بهینه‌سازی کرد و در نتیجه، روش‌های مبتنی بر جستجو کارایی تحلیلی خود را در یافتن چرایی رفتار مدل از دست می‌دهند.

در واکنش به این محدودیت‌های ذاتی در استخراج علیت، بخش مهمی از ادبیات این حوزه، صورت‌مسئله‌ی انتساب را به‌شکلی کاربردی‌تر بازتعریف کرد. راشکین و همکاران (۲۰۲۳)~\cite{rashkin2023measuring} با معرفی چارچوب «انتساب‌پذیر به منابع مشخص» (\lr{AIS})، مفهوم انتساب را به‌عنوان اعتبارسنجی واقعیت در نظر گرفتند. در این معماری، هدف کشف مکانیسم درونی شبکه نیست، بلکه بررسی می‌شود که آیا تک‌تک گزاره‌های تولیدشده توسط مدل، قابلیت‌انتساب و ردیابی به یک سند منبع خارجی و موثق را دارند یا خیر. اگرچه این رویکرد در مقابله با توهم اطلاعاتی و برای استقرار ایمن سیستم‌ها یک دستاورد کلیدی است، اما بررسی آن در اینجا مرز میان حقیقت‌سنجی و تفسیرپذیری علّی را روشن می‌سازد. یک خروجی متنی می‌تواند کاملاً مستند و منطبق بر واقعیت بیرونی باشد، اما چارچوب \lr{AIS} پاسخی به این پرسش نمی‌دهد که مدل با تکیه بر کدام مفاهیم در پرامپت ورودی کاربر، تصمیم به فراخوانی و تولید آن دانش خاص گرفته است.

بنابراین، اگرچه رویکردهای مبتنی بر پرامپت و خودتوضیحی در نگاه نخست سرراست‌ترین ابزار برای تعامل با مدل‌های مولد امروزی به نظر می‌رسند، اما اتکای آن‌ها به خروجی‌های سطحی و فقدان وفاداری به مکانیزم‌های علّی درون شبکه، آن‌ها را از منظر تفسیرپذیری دقیق نامعتبر می‌سازد. این ناکامی، ضرورت بازگشت به روش‌های ساختاریافته‌تر و متقدم‌تر ارزیابی جعبه‌سیاه و وفق آن‌ها به تولید متن آزاد را ایجاب می‌کند؛ روش‌هایی که وابستگی‌های علّی را نه از طریق پرسش از مدل، بلکه از طریق دستکاری سیستماتیک ورودی‌های آن استخراج می‌کنند.


\subsubsection{انتساب ویژگی از طریق اختلال در ورودی}\label{sec:3-1-3-2}

در غیاب دسترسی به گراف محاسباتی و معماری داخلی، متقدم‌ترین راهبرد برای کشف روابط علّی در یک مدل جعبه‌سیاه، استفاده از تکنیک‌های مبتنی بر اختلال است~\cite{guidotti2018survey}. این خانواده از روش‌ها که ریشه در وظایف کلاسیک دسته‌بندی و متمایزگر دارند، یک رویکرد کاملاً رفتارگرایانه را اتخاذ می‌کنند. در این رهیافت، مدل به‌عنوان یک اوراکل ریاضی ($f$) در نظر گرفته می‌شود و تحلیل‌گر تلاش می‌کند با اعمال تغییرات سیستماتیک بر روی نمونه‌ی ورودی ($x \rightarrow x'$) و رصد نوسانات متناظر در خروجی ($f(x')$)، رفتار سیستم را مهندسی معکوس کند. با ایجاد این نگاشت آماری میان ورودی‌های مخدوش‌شده و خروجی‌ها، می‌توان سهم علّی هر ویژگی (مانند کلمات یا پیکسل‌ها) را تخمین زد.

با گسترش \lr{LLM}ها، تلاش‌های متعددی برای تعمیم این متدولوژی‌های کلاسیک به دامنه‌ی تولید متن صورت گرفت. با این حال، همان‌گونه که در ادامه واکاوی می‌شود، گذار از داده‌های ساختاریافته به متن پیوسته و گسسته، این رویکردها را با موانع نظری عمیقی مواجه می‌سازد.

\paragraph{تقریب‌های محلی و پایه‌های ریاضیاتی.}

\normalfootnotes{}
ریبیرو و همکاران (۲۰۱۶)~\cite{ribeiro2016lime} با معرفی متدولوژی \lr{LIME}\enfootnote{local interpretable model-agnostic explanations} و ارزیابی آن بر روی وظایف کلاسیکی نظیر دسته‌بندی متن (با مدل‌هایی چون جنگل تصادفی) و دسته‌بندی تصویر (با شبکه‌های عصبی پیچشی)، پایه‌گذار رویکرد انتساب ویژگی در محیط جعبه‌سیاه محسوب می‌شوند. مسئله‌ی محوری در این پژوهش، غیرخطی بودن و پیچیدگی شدید مرز تصمیم‌گیری در شبکه‌های عصبی بود که استخراج یک توضیح سراسری خطی را ناممکن می‌ساخت. روش \lr{LIME} برای غلبه بر این محدودیت، از تقریب محلی بهره برد؛ به این معنا که رفتار پیچیده‌ی مدل را صرفاً در همسایگی یک نقطه‌ی داده ($x$) تفسیر می‌کند.
\twocolumnfootnotes{}

در این الگوریتم، ابتدا با حذف یا ماسک کردن تصادفی بخش‌هایی از داده‌ی ورودی (مانند توکن‌ها در متن)، مجموعه‌ای از نمونه‌های مصنوعی و اختلال‌یافته ($x'$) تولید شده و خروجی مدل هدف ($f(x')$) برای تک‌تک آن‌ها استعلام می‌گردد. سپس، یک تابع مجاورت ($\pi_x(x')$) برای وزن‌دهی به این نمونه‌ها بر اساس فاصله‌ی هندسی آن‌ها تا نمونه‌ی اصلی تعریف می‌شود. در گام نهایی، یک مدل جانشینِ ذاتاً شفاف ($g$) از یک خانواده‌ی قابل‌تفسیر ($G$) – نظیر یک رگرسیون خطی خلوت – آموزش می‌بیند تا پیش‌بینی‌های مدل اصلی را در آن همسایگی محدود تقلید کند. تابع خطای این فرایند به‌صورت زیر فرمول‌بندی می‌شود:
\begin{equation}\label{eq:chap3_lime}
    \text{\lr{Explanation}}(x) = \arg\min_{g \in G} \mathcal{L}(f, g, \pi_x) + \Omega(g)
\end{equation}
این رابطه‌ی بهینه‌سازی، تعادلی میان دو نیروی متضاد اما ضروری در فرایند تفسیر برقرار می‌سازد: وفاداری محلی در برابر سادگی ادراکی انسان.
مؤلفه‌ی نخست، یعنی $\mathcal{L}(f, g, \pi_x)$، خطای عدم‌وفاداری مدل جانشینِ ساده ($g$) نسبت به مدل پیچیده‌ی جعبه‌سیاه ($f$) را اندازه می‌گیرد؛ تابع مجاورت $\pi_x(x')$ که مانند یک هسته‌ی فاصله‌سنج نمایی عمل می‌کند، به نمونه‌های مختل‌شده‌ای که فاصله‌ی کمتری تا نمونه‌ی اصلی $x$ دارند، وزن به‌مراتب بالاتری اختصاص می‌دهد. از منظر شهودی، این جزء تضمین می‌کند که مدل ساده، رفتار مرز تصمیم‌گیری پیچیده را صرفاً در همسایگی بلافصل نقطه‌ی $x$ با بیشترین دقت تقلید کند، بدون آنکه نگران ناتوانی ذاتی یک مدل خطی در توصیف کل فضای تصمیم‌گیری سراسری باشد.
مؤلفه‌ی دوم، یعنی $\Omega(g)$، جریمه‌ی پیچیدگی یا عدم‌شفافیت مدل جانشین است که خوانایی آن را تضمین می‌کند؛ برای نمونه در یک رگرسیون خطی، $\Omega(g)$ جریمه‌ای برای محدود کردن تعداد ضرایب غیرصفر (خلوت‌بودن مدل) است تا کاربر انسانی بتواند مهم‌ترین عوامل تصمیم را به‌آسانی تفکیک کند. بنابراین، وزن‌های یادگرفته‌شده در مدل بهینه‌ی $g^*$ مستقیماً به‌عنوان سهم یا اهمیت خطی هر ویژگی در اتخاذ آن تصمیم محلی تفسیر می‌شوند. با وجود انعطاف‌پذیری بالای \lr{LIME}، این روش به‌دلیل حساسیت شدید به نحوه‌ی تعریف تابع مجاورت و نمونه‌برداری، فاقد ضمانت‌های تئوریک برای توزیع عادلانه و پایدار سهم ویژگی‌ها در حضور متغیرهای همبسته است.

برای رفع این نقصان تئوریک و ایجاد یک بنیان ریاضی مستحکم، لوندبرگ و لی (۲۰۱۷)~\cite{lundberg2017unified} چارچوب جامع \lr{SHAP}\enfootnote{Shapley additive explanations} را ارائه دادند که در ابتدا بر روی داده‌های ساختاریافته‌ی جدولی و وظایف دسته‌بندی تصویر ارزیابی و اثبات شد. این پژوهش، با استفاده از اصول نظریه‌ی بازی‌های همکارانه\enfootnote{cooperative game theory}، روش‌های تقریب محلی را یکپارچه ساخت. در این مدل‌سازی، فرایند تولید خروجی در یک مدل جعبه‌سیاه، معادل یک بازی گروهی در نظر گرفته می‌شود؛ خروجی مدل (یا میزان انحراف آن از میانگین) به‌عنوان پاداش و ویژگی‌های ورودی (توکن‌ها) به‌عنوان بازیکنانی فرض می‌شوند که برای کسب این پاداش با یکدیگر ائتلاف کرده‌اند. ارزش شاپلی ($\phi_i$)، که سهم عادلانه‌ی ویژگیِ $i$ را تعیین می‌کند، با محاسبه‌ی میانگین مشارکت حاشیه‌ای آن ویژگی در تمامی ائتلاف‌های ممکن از سایر ویژگی‌ها به‌دست می‌آید:
\begin{equation}\label{eq:chap3_shapley}
    \phi_i = \sum_{S \subseteq F \setminus \{i\}} \frac{|S|! (|F| - |S| - 1)!}{|F|!} [f(S \cup \{i\}) - f(S)]
\end{equation}
در این فرمول‌بندی، $F$ مجموعه‌ی تمامی ویژگی‌های ورودی و $S$ یک زیرمجموعه (ائتلاف) از آن‌هاست که ویژگی هدف ($i$) در آن غایب است.
شهود ملموس این فرمول را می‌توان با تمثیل ورود تصادفی بازیکنان به بازی دریافت: فرض کنید تمام ویژگی‌های ورودی با یک ترتیب کاملاً تصادفی وارد میدان می‌شوند. عبارت تفاضلی $f(S \cup \{i\}) - f(S)$ دقیقاً مشارکت حاشیه‌ای\enfootnote{marginal contribution} ویژگی $i$ را نشان می‌دهد؛ یعنی میزان جهش یا افتی که در خروجی مدل، دقیقاً در لحظه‌ی اضافه شدن ویژگی $i$ به ائتلاف حاضرین ($S$) حاصل می‌شود.
ضریب ترکیبیاتی $\frac{|S|! (|F| - |S| - 1)!}{|F|!}$، احتمال ریاضی دقیق آن است که در میان تمامی $|F|!$ جایگشتِ محتمل برای ترتیب ورود ویژگی‌ها، دقیقاً اعضای ائتلاف $S$ پیش از ویژگی $i$ و مابقی ویژگی‌ها پس از آن وارد شوند. بنابراین، ارزش شاپلی ($\phi_i$) در حقیقت امید ریاضیِ سهم حاشیه‌ای ویژگی $i$ در تمام ترتیب‌های ممکنِ ورود است. % chktex 40
از منظر تئوریک، استحکام بنیادین \lr{SHAP} در این است که تنها روش توزیع اهمیتی است که چهار اصل موضوعه‌ی عدالت در نظریه‌ی بازی‌ها را به‌طور هم‌زمان ارضا می‌کند: دقت محلی (تضمین اینکه مجموع سهم‌های اختصاص‌یافته برابر با تفاوت خروجی کل با خروجی پایه‌ی میانگین باشد)، فقدان (اختصاص سهم صفر به ویژگی‌های بی‌اثر و غایب)، تقارن (تخصیص ارزش یکسان به ویژگی‌هایی که سهم حاشیه‌ای یکسانی دارند)، و سازگاری (تضمین اینکه با افزایش تأثیر واقعی یک ویژگی در مدل، ارزش شاپلی آن هرگز کاهش نیابد).

تعمیم رویکردهای انتساب محلی نظیر \lr{LIME} و \lr{SHAP} از مدل‌های متمایزگر کلاسیک به مدل‌های زبانی مولد، با دو چالش بنیادین مواجه بود. نخست آنکه محاسبه‌ی مقادیر شاپلی دارای پیچیدگی ترکیبیاتیِ $\mathcal{O}(2^N)$ است، که در رویارویی با پرامپت‌های طولانی متشکل از هزاران توکن، از نظر محاسباتی کاملاً فلج‌کننده و غیرعملی است. چالش دوم، ماهیت تابع خروجی است؛ این روش‌ها اساساً برای انتساب اعتبار به یک خروجی اسکالر (مانند احتمال تعلق به یک کلاس خاص) فرمول‌بندی شده‌اند. اما در مدل‌های زبانی، خروجی یک توالی متغیر و طولانی از کلمات است و انتساب یک امتیاز واحد شاپلی ($\phi_i$) به فرایند تولید یک متن آزاد، از نظر ریاضی بدتعریف است. 

برخی پژوهش‌ها برای فرار از پیچیدگی انتساب اعتبار به یک توالی متنی طولانی، صورت‌مسئله را با تقلیل خروجی شبکه به یک امتیاز اسکالر ساده‌سازی کردند. برای نمونه، هاکمن و همکاران (۲۰۲۴)~\cite{hackmann2024wordimportance} متدولوژی‌ای ارائه دادند که در آن، به‌جای بررسی احتمال تولید توکن‌ها، ویژگی‌های کلان متن تولیدشده توسط توابع امتیازدهی ازپیش‌تعریف‌شده (نظیر ارزیابی سطح خوانایی، میزان پرحرفی یا لحن) سنجیده می‌شود. با ماسک کردن کلمات پرامپت و رصد تغییرات این امتیازات در متن جدید، سهم هر کلمه در شکل‌گیری آن ویژگی محاسبه می‌گردد. اگرچه این رویکرد ابزاری کارآمد برای درک تأثیر پرامپت بر سبک و فرم خروجی فراهم می‌آورد، اما علیت منطقی را به نوسان آماری سطحی تقلیل می‌دهد و از کشف مکانیسم استدلالی و پیش‌نیازهای اطلاعاتی مدل برای رسیدن به خود پاسخ عاجز است.

برای غلبه بر این محدودیت‌ها به‌شکلی اصولی‌تر، انوئن و همکاران (۲۰۲۴)~\cite{enouen2024textgenshap} چارچوب \lr{TextGenSHAP} را مشخصاً برای تفسیر خروجیِ \lr{LLM}ها در وظایفی با پرامپت‌های بسیار طولانی (نظیر سیستم‌های پرسش‌وپاسخ تقویت‌شده با بازیابی و سیستم‌های بازیابی سند) توسعه دادند. این الگوریتم برای حل مشکل انفجار محاسباتی، از یک روش گروه‌بندی سلسله‌مراتبی بهره می‌برد؛ به این معنا که متن ورودی را به‌جای توکن‌های مجزا، به واحدهای معنایی بزرگ‌تر (مانند جملات یا پاراگراف‌ها) تقسیم می‌کند تا فضای نمونه‌برداری به‌شدت کاهش یابد. نوآوری اصلی در \lr{TextGenSHAP}، بازتعریف تابع ارزش است. از آنجا که خروجی مدل یک دنباله‌ی متنی تولیدشده است ($Y = y_1, y_2, \dots, y_m$)، تابع ارزش به لگاریتم احتمال تولید کل این دنباله‌ی هدف، مشروط بر حضور زیرمجموعه‌ی $S$ از پرامپت بازتعریف می‌شود:
\begin{equation}\label{eq:chap3_textgenshap}
    v(S) = \sum_{j=1}^{m} \log P(y_j \mid y_{<j}, x_S)
\end{equation}
در این فرمول‌بندی، $x_S$ نشان‌دهنده‌ی پرامپت ورودی است که تنها قطعات موجود در ائتلاف $S$ در آن حفظ شده و مابقی ماسک یا حذف شده‌اند.
شهود ریاضی نهفته در این فرمول‌بندی، حل معضل حاصل‌ضرب احتمالات در مدل‌های زبانی خودهمبسته است؛ زیرا احتمال توأم تولید کل توالی (به‌صورت ریاضی یعنی $P(Y \mid x_S) = \prod_{j=1}^m P(y_j \mid y_{<j}, x_S)$) با افزایش طول متن ($m$) به‌سرعت افت کرده و به‌دلیل پدیده‌ی زیرسیلان عددی\enfootnote{numerical underflow}، تفاوت میان ائتلاف‌ها را محو می‌سازد. اعمال لگاریتم این حاصل‌ضرب را به یک مجموع جمع‌پذیر از لگاریتم احتمالات شرطی تبدیل کرده و یک تابع ارزش اسکالر و پایدار ($v(S)$) می‌سازد. از منظر شهودی، $v(S)$ سنجه‌ای یکپارچه از «میزان اطمینان کلی مدل در تولید پیوسته‌ی تمامی توکن‌های پاسخ» در حضور ائتلاف $S$ است. با جایگذاری این تابع ارزش تجمیعی به‌جای تابع $f(S)$ در معادله‌ی شاپلی، \lr{TextGenSHAP} قادر است با استعلام لگاریتم احتمالات از رابط جعبه‌سیاه، سهم علّی دقیق هر جمله‌ی ارائه‌شده در پرامپت را در شکل‌گیری متن نهایی محاسبه کند.

با وجود نوآوری‌های محاسباتی در روش‌هایی نظیر \lr{TextGenSHAP}، تعمیم رویکردهای اختلالی و مستقل از مدل (اعم از \lr{LIME} و \lr{SHAP}) برای ارزیابی سیستم‌های جعبه‌سیاه زبانی با دو چالش ساختاری و بازدارنده مواجه است که کارایی آن‌ها را نقض می‌کند. نخستین چالش در پیش‌نیازهای استعلام این متدولوژی‌ها نهفته است. فرمول‌بندی توابع ارزش و محاسبه‌ی مجاورت در تمامی این چارچوب‌ها، وابستگی مطلقی به استخراج لگاریتم احتمالات دقیق (لوجیت‌ها) برای تک توکن‌های خروجی ($P(y \mid x)$) دارد. از آنجا که \lr{API}ها (به‌عنوان سیستم‌های جعبه‌سیاه مطلق) غالباً دسترسی به توزیع احتمالات لایه‌ی آخر را مسدود کرده و صرفاً رشته‌ی متنی نهایی را ارائه می‌دهند، پیاده‌سازی این رویکردها بر روی مسلط‌ترین مدل‌های امروزی عملاً غیرممکن می‌گردد. 

دومین و مهم‌ترین محدودیت این روش‌ها، کوری معنایی در ارزیابی متون باز است. توابع ارزش در رهیافت‌های اختلالی، صرفاً احتمال تولید متوالی توکن‌های اولیه را در توالی مخدوش‌شده ردیابی می‌کنند. این در حالی است که در فرایند تولید متن آزاد (برخلاف دسته‌بندی کلاسیک)، کاهش احتمال تولید یک کلمه‌ی خاص لزوماً معادل با حذف آن مفهوم از خروجی نیست. با انتخاب پرامپت به‌گونه‌ای که یک توکن خاص از خروجی حذف شود، مدل ممکن است به‌دلیل افزونگی زبانی، همان منطق و بار معنایی را با استفاده از مجموعه‌ی کلمات و ساختارهای نحوی متفاوتی (نظیر پارافریز) بازتولید کند. با این وجود، روش‌های اختلالی این افت احتمال خُرد را به‌عنوان فروپاشی موفق معنا در خروجی تفسیر کرده و به اشتباه، سهم علّی بالایی را به متغیرهایی اختصاص می‌دهند که تنها در نحوه‌ی بیان نحوی مؤثر بوده‌اند، نه در شالوده‌ی استدلالی پاسخ.

تلاش برای استخراج ساختارهای منطقی‌تر از اختلال، در پژوهش شن و همکاران (۲۰۲۳)~\cite{shen2023inferencelogic} با بهره‌گیری از نظریه بازی‌ها دنبال شد. متدولوژی آن‌ها بر این فرض استوار است که مدل‌های زبانی استنتاج خود را نه بر روی توکن‌های منفرد، بلکه بر مفاهیم گسسته بنا می‌کنند. آن‌ها با ماسک کردن ترکیبات مختلفی از کلمات پرامپت و ثبت تغییرات در احتمال خروجی، تعاملات چندجانبه میان توکن‌ها را محاسبه کردند. کلماتی که دارای تعاملات هم‌افزای قدرتمندی در حفظ خروجی بودند، گروه‌بندی شده و به‌عنوان یک مفهوم نمادین استخراج گردیدند. با این وجود، محاسبه‌ی این تعاملات برای کشف ساختارهای منطقی، نیازمند جستجو در فضای تمامی ائتلاف‌های ممکن از کلمات است که مجدداً این رهیافت را با سد محاسباتی و پیچیدگی نمایی ($\mathcal{O}(2^N)$) مواجه می‌سازد.

این بن‌بست ترکیبیاتی در مرور جامع استپین و همکاران (۲۰۲۱)~\cite{stepin2021survey} بر روش‌های تقابلی و خلاف واقع نیز با وضوح مشهود است. رویکردهای خلاف واقع به دنبال یافتن حداقل تغییرات ممکن در ورودی برای تغییر صریح خروجی هستند (پاسخ به پرسش «چرا $P$ به‌جای $Q$؟»). اگرچه این رویکرد از منظر تئوری برای اثبات و استخراج روابط علّی بسیار قدرتمند است، اما تولید تغییرات خلوت و معنادار در فضاهای با ابعاد بالا به‌شدت چالش‌برانگیز است و همیشه معنادار نیست. در دامنه‌ی تولید متن باز، تلاش برای یافتن این زیرمجموعه‌ی بهینه‌ی کلمات از طریق اختلالات جعبه‌سیاه، نیازمند فراخوانی‌های نمایی از رابط برنامه‌نویسی است که جستجوی کورکورانه برای یافتن گزاره‌ی خلاف واقع مطلوب را در پرامپت‌های طولانی عملاً غیرقابل‌اجرا می‌کند. در نهایت، فارغ از این هزینه‌ی بالا، اعمال هرگونه اختلال در این روش‌ها به تخریب ساختار گرامری ورودی و بروز خطای شوک توزیعی نیز منجر می‌گردد؛ چالشی بنیادین که در ادامه‌ی این بخش بررسی خواهد شد.

\paragraph{چالش شوک توزیعی در زبان طبیعی.}

اعمال متدولوژی‌های مبتنی بر اختلال و تقریب محلی، در مواجهه با ماهیت گسسته‌ی زبان طبیعی، پژوهشگران را با یک چالش بنیادین روبه‌رو ساخت. لی و همکاران (۲۰۱۷)~\cite{li2017understanding} با معرفی روش‌های مبتنی بر فرسایش یا حذف بازنمایی، یکی از پایه‌ای‌ترین رویکردهای انتساب ویژگی در پردازش زبان طبیعی را شکل دادند. متدولوژی آن‌ها بر این اصل استوار بود که با حذف سیستماتیک ابعاد بردارهای تعبیه یا کلمات ورودی، میزان تغییرات در خروجی شبکه سنجیده شود. در این پژوهش، از طریق یک رویکرد جستجوی تقویتی تلاش شد تا کوچک‌ترین مجموعه از کلمات یافته شود که حذف آن‌ها، تصمیم نهایی مدل را در وظایف دسته‌بندی معکوس (فلیپ) می‌کند. اگرچه این رویکرد به‌عنوان یک استاندارد در ارزیابی اهمیت کلمات پذیرفته شد، اما بررسی دقیق آن یک نارسایی معنایی جدی را نمایان می‌سازد؛ حذف تصادفی یا ساختاریافته‌ی توکن‌ها، انسجام نحوی و گرامری متن را متلاشی می‌کند. در نتیجه، مشخص نیست که آیا تغییر در خروجی مدل به‌دلیل از دست رفتن یک پیش‌نیاز منطقی بوده است، یا صرفاً واکنشی به یک جمله‌ی درهم‌ریخته و فاقد معنا.

این پدیده که همان رویارویی با داده‌های \lr{OOD} است، در پژوهش وفا و همکاران (۲۰۲۱)~\cite{vafa2021rationales} به‌شکل عمیق‌تری واکاوی شد. آن‌ها در متدولوژی خود به دنبال استخراج یک توجیه (زیرمجموعه‌ای کوچک از توکن‌های ورودی) بودند که برای حفظ پیش‌بینی متوالی مدل کافی باشد. برای این کار توکن‌ها به‌تدریج با استفاده از یک رویکرد جستجوی حریصانه از ورودی حذف می‌شدند. نویسندگان در ارزیابی این روش متوجه یک تناقض ساختاری شدند؛ متون پاره‌پاره‌ای که در اثر حذف توکن‌ها ایجاد می‌شدند، به‌شدت از توزیع داده‌های زبانی طبیعی که مدل روی آن‌ها آموزش دیده بود، فاصله داشتند. مدل‌های زبانی در مواجهه با این داده‌های \lr{OOD} دچار افت شدید عملکرد و تولید خروجی‌های نامربوط می‌شدند. برای غلبه بر این فروپاشی، نویسندگان ناگزیر شدند مدل اصلی را بر روی توزیع متون مخدوش‌شده تنظیم ظریف کنند تا مدل توانایی پردازش آن‌ها را به‌دست آورد. این الزام، ذات جعبه‌سیاه را نقض می‌کند؛ چرا که در ارزیابی مدل‌های تجاری و مسدود، امکان دستکاری وزن‌ها یا تنظیم ظریف شبکه اساساً وجود ندارد. این امر ثابت می‌کند که اتکا به حذف گسسته‌ی کلمات برای ممیزی مدل‌های ایستا، رویکردی فاقد کارایی عملی است.

در حوزه‌ی بینایی ماشین، آگاروال و نگوین (۲۰۲۰)~\cite{agarwal2020explaining} برای حل معضل مشابهی (تولید تصاویر غیرطبیعی در اثر سیاه کردن پیکسل‌ها)، استفاده از شبکه‌های ترمیم‌کننده\enfootnote{inpainters} را پیشنهاد دادند تا جای خالی پیکسل‌های حذف‌شده با پس‌زمینه‌ای واقع‌گرایانه (مبتنی بر مدل‌های مولد) پر شود. با این وجود، تعمیم این ایده به پردازش زبان طبیعی با بن‌بست منطقی روبه‌روست. در متون زبانی، استفاده از یک مدل ثانویه برای ترمیم یا جایگزینی یک کلمه در میانه متن، صرفاً یک اصلاح پیکسلی یا هندسی نیست؛ بلکه بار منطقی، لحن و جهت‌گیری کل گزاره را تغییر می‌دهد. در نتیجه، تغییر خروجی مدل هدف دیگر قابل‌انتساب به توکن اولیه‌ی حذف‌شده نخواهد بود، بلکه ناشی از معنای جدیدی است که کلمه‌ی ترمیم‌شده به جمله القا کرده است.

با درک این محدودیت‌های توزیعی و معنایی، ژائو و شان (۲۰۲۴)~\cite{zhao2024reagent} چارچوب \lr{ReAGent} را برای مدل‌های زبانی مولد توسعه دادند تا بدون تخریب گرامر، میزان اهمیت ویژگی‌ها را به‌عنوان پیش‌آیند انتساب دهند. متدولوژی آن‌ها به‌جای حذف خام کلمات، از یک رویکرد جایگزینی آگاه از بافتار بهره می‌برد؛ بدین صورت که توکن هدف ($x_i$) در پرامپت، با مجموعه‌ای از کلمات محتمل ($\tilde{x}_i$) که توسط یک مدل زبانی ماسک‌شده (نظیر \lr{RoBERTa}~\cite{liu2021roberta}) پیش‌بینی شده‌اند، جایگزین می‌گردد. سپس، اهمیت توکن اصلی در تولید توکن هدفِ $y_t$، از طریق امید ریاضی افت احتمال در خروجی مدل جعبه‌سیاه سنجیده می‌شود:
\begin{equation}\label{eq:chap3_reagent}
    R(x_i \rightarrow y_t) = \mathbb{E}_{\tilde{x}_i} \left[ P(y_t \mid x) - P(y_t \mid x_{\setminus i}, \tilde{x}_i) \right]
\end{equation}
در این معادله، $x_{\setminus i}, \tilde{x}_i$ نشان‌دهنده‌ی پرامپت دستکاری‌شده است که در آن کلمه‌ی اصلی $x_i$ با جایگزین محتمل $\tilde{x}_i$ تعویض شده است.
شهود این فرمول، استوار بر مقایسه‌ی خروجی مدل در جهان واقعی با جهان‌های جایگزینِ محتمل و معتبر زبانی است. جمله‌ی نخست ($P(y_t \mid x)$)، احتمال تولید توکن هدف را در حضور جمله‌ی دست‌نخورده نشان می‌دهد. جمله‌ی دوم ($P(y_t \mid x_{\setminus i}, \tilde{x}_i)$)، این احتمال را زمانی می‌سنجد که کلمه‌ی $x_i$ با واژه‌ای هم‌خوان با ساختار دستوری جمله تعویض شده باشد. تفاضل این دو عبارت، میزان افت اطمینان مدل را منعکس می‌سازد و عملگر امید ریاضی ($\mathbb{E}$) اثر این تغییر را بر روی توزیعی از واژگان محتمل میانگین‌گیری می‌کند. از منظر شهودی، اگر جایگزین کردن $x_i$ با سایر کلمات محتمل زبان موجب سقوط معنادار احتمال تولید $y_t$ شود (تفاضل مثبت بزرگ)، اثبات می‌گردد که مدل مشخصاً به محتوای معناییِ منحصربه‌فرد این کلمه تکیه داشته است، نه صرفاً به نقش نحوی آن در اسکلت جمله. روش \lr{ReAGent} این محاسبات اختلالی را به‌صورت بازگشتی انجام می‌دهد تا توزیع اهمیت توکن‌ها در طول فرایند تولید گام‌به‌گام به‌روزرسانی شود.

در رویکردی موازی و برای فرار کامل از نیاز به مدل‌های ثانویه در حفظ انسجام متن، سیفکا و لیوتکوس (۲۰۲۳)~\cite{cifka2023blackbox} متدولوژی غیرمخربِ کاوش طول بافتار\enfootnote{context-length probing} را معرفی کردند. مسئله‌ی اصلی در این پژوهش نیز ارائه‌ی یک متدولوژی انتساب مستقل از مدل بدون تولید داده‌های \lr{OOD} بود. متدولوژی آن‌ها به‌جای ایجاد حفره در متن یا جایگزینی کلمات، بر مکانیزم افزودن تدریجی و توکن‌به‌توکن پیشوند ورودی استوار است؛ استدلال نویسندگان این است که پیشوندهای یک متن طبیعی، خود جملاتی طبیعی (هرچند ناتمام) هستند و مدل‌های زبانی خودهمبسته ذاتاً برای پردازش این پیشوندها آموزش دیده‌اند. در این الگوریتم، برای سنجش سهم یک توکن ورودی ($x_i$) در پیش‌بینی کلمه‌ی هدف ($y_t$)، مدل با دو طول بافتار متوالی استعلام می‌شود: یک بار با متنی که دقیقاً قبل از آن توکن پایان می‌یابد ($x_{1:i-1}$) و یک بار با افزودن خود آن توکن ($x_{1:i}$). میزان اثرگذاری ناشی از افزوده شدن توکن $x_i$ به‌عنوان پیش‌آیند، مستقیماً از طریق تفاضل لگاریتمِ احتمالات خروجی محاسبه می‌گردد:
\begin{equation}\label{eq:chap3_context_length_probing}
    I(x_i \rightarrow y_t) = \log P(y_t \mid x_{1:i}) - \log P(y_t \mid x_{1:i-1})
\end{equation}
شهود ریاضی نهفته در این فرمول، اندازه‌گیری افزایش اطلاعات\enfootnote{information gain} یا جهش در لگاریتم شانس پیش‌بینی در اثر افزوده شدن یک نشانه به بافتار است. در این دیدگاه، مدل در دو وضعیت ترتیبی ارزیابی می‌شود: ابتدا بدون مشاهده‌ی توکن $x_i$ (صرفاً با تکیه بر تاریخچه‌ی پیشین $x_{1:i-1}$)، اطمینان مدل نسبت به توکن هدف سنجیده می‌شود؛ سپس با اضافه شدن توکن $x_i$ به افق دید مدل ($x_{1:i}$)، میزان تغییر در لگاریتم احتمال رصد می‌گردد. اگر مقدار $I$ مثبت و بزرگ باشد، نشان می‌دهد که کلمه‌ی $x_i$ حاوی سیگنال اطلاعاتی کلیدی بوده که موجب جهش در اطمینان مدل شده است. در مقابل، مقادیر نزدیک به صفر گویای خنثی بودن اطلاعاتی کلمه در آن جایگاه بافتاری است و مقادیر منفی نشان می‌دهند که ورود این کلمه حتی پیش‌بینی مدل را نسبت به هدف تضعیف کرده است. یافته‌های این پژوهش نشان داد که این روش بسیار سبک، بدون نیاز به هزاران نمونه‌برداری تصادفی، دقتی رقابتی با روش‌های پیچیده‌ی جعبه‌سفید (نظیر \lr{IG}) ارائه می‌دهد و فروپاشی توزیعی را خنثی می‌کند. با این وجود، متدولوژی آن‌ها با یک محدودیت ساختاری مهم روبه‌روست؛ اتکای مطلق به ارزیابی چپ‌به‌راست و ترتیبی باعث می‌شود این روش تنها تأثیر مستقیم کلمات را در همان نقطه‌ی حضورشان بسنجد و از مدل‌سازی تعاملات ترکیبی و غیرخطی میان کلماتی که با فاصله از هم در متن قرار دارند، عاجز بماند.

اگرچه متدولوژی‌هایی نظیر \lr{ReAGent} و کاوش طول بافتار موفق شده‌اند با دور زدن اختلالات مخرب، خطای ناشی از شوک توزیعی (\lr{OOD}) را مهار کنند، اما حل این معضل به‌تنهایی برای ممیزی مدل‌های زبانیِ مدرن کفایت نمی‌کند. همان‌طور که معادلات~\ref{eq:chap3_reagent} و~\ref{eq:chap3_context_length_probing} به وضوح نشان می‌دهند، هسته‌ی محاسباتیِ این رویکردهای نوین همچنان بر استخراج مستقیم لگاریتم احتمالات خروجی ($P$) استوار است. در نتیجه، این متدولوژی‌ها نیز وارث همان بن‌بست دسترسی در \lr{API}های تجاری مسدود هستند که پیش‌تر در نقد چارچوب‌های اختلالی کلاسیک مورد بحث قرار گرفت.

فراتر از موانع دسترسی، این راهکارها کماکان در حصار تقلیل‌گرایی و کوری معنایی گرفتارند. برای نمونه، اگرچه \lr{ReAGent} تلاش کرده است با میانگین‌گیری سهم توکن‌های داخل یک توالی با طول معین در خروجی، ارزیابی را از سطح خرد به توالی فراتر ببرد، اما این تجمیع ریاضیاتی نمی‌تواند ماهیت منعطف و چندوجهی زبان را رمزگشایی کند. همان‌گونه که استدلال شد، قابلیت‌ذاتی مدل‌ها در جبران ساختار از طریق پارافریز، نوسانات آماریِ توزیع توکن‌ها را به شاخصی نامعتبر برای سنجشِ علیت تبدیل می‌کند. متدولوژی‌هایی که اساسِ داوری خود را بر این نوسانات بنا می‌کنند – حتی با وجود حفظ کامل گرامر – همچنان تغییر در قالب بیان را با فقدان یا فروپاشی استدلال اشتباه می‌گیرند. این ناکامی ساختاری ثابت می‌کند که برای حسابرسی منطقی یک متن آزاد، صرف رفع خطای گرامری ورودی کافی نیست؛ بلکه واحد تحلیل باید از اساس به سطوح بالاتری ارتقا یابد.

برای پر کردن این شکاف معنایی و فرار از تقلیل‌گرایی توکن‌محور، کافانیا و همکاران (۲۰۲۳)~\cite{cafagna2023interpreting} استدلال کردند که اعمال مستقیم الگوریتم‌های انتساب بر روی ویژگی‌های خام (مانند پیکسل‌ها یا نشانه‌های مجزای متنی) لزوماً به استخراج توضیحاتی منسجم و قابل‌درک برای انسان منجر نمی‌شود. متدولوژی آن‌ها برای حل این مسئله، بر ادغام پیش‌فرض‌های بصری-معنایی\enfootnote{semantic-visual priors} استوار است. در این ساختار، به‌جای ارزیابی نقطه‌به‌نقطه، ابتدا از مدل‌های تقطیع معنایی برای دسته‌بندی و یکپارچه‌سازی ویژگی‌های خام به مفاهیم کلان (نظیر اشیاء فیزیکی در یک تصویر) استفاده می‌شود. سپس، الگوریتم‌هایی نظیر \lr{SHAP} به‌جای محاسبه‌ی ارزش حاشیه‌ای برای توکن‌های منفرد، سهم انتسابیِ این مفاهیم شیء‌محور را در خروجی مدل محاسبه می‌کنند. یافته‌های آن‌ها نشان داد که انتقال واحد تحلیل از سطح ویژگی خام به سطح شیء، وفاداری و خوانایی توضیحات را ارتقا می‌دهد. با این وجود، انتقال این راهبرد از دامنه‌ی چندوجهی به پردازش متون انتزاعی با بن‌بست ساختاری همراه است؛ در تصویر، مرز فیزیکی یک شیء به‌وضوح تعریف شده است، اما در زبان طبیعی، یک پیش‌فرض منطقی یا مفهوم انتزاعی ماهیتی پیچیده دارد و تقطیع قطعی آن نیازمند تعاریف محدودکننده‌ی مفهوم است.

\paragraph{بهره‌برداری عملی از نتایج انتساب.}

در رویکردی متفاوت که هدف آن پاسخ به چگونگی استخراج این نمرات و بهره‌برداری عملی از خروجی همین الگوریتم‌های اختلالی بود، کریشنا و همکاران (۲۰۲۳)~\cite{krishna2023post} نشان دادند که توضیحات پساعملی می‌توانند برای بهبود عملکرد خود مدل زبانی به‌کار روند. آن‌ها چارچوب \lr{AMPLIFY} را معرفی کردند که بر یک حلقه‌ی بازخورد استوار است. در این روش، ابتدا میزان اثرگذاری ویژگی‌های پرامپت با روش‌های انتساب محلی محاسبه می‌شود. سپس، این نمرات عددی، توسط یک مؤلفه‌ی میانی مبتنی بر قواعد به الگوهای متنیِ زبان طبیعی (برای مثال جمله‌ای نظیر «مدل برای این پاسخ بیشترین توجه را به کلمات $X$ و $Y$ داشته است») ترجمه می‌شوند. در نهایت، این توضیحات متنی به‌عنوان بافتار تکمیلی در یادگیری درون‌بافتاری مجدداً به پرامپت مدل تزریق می‌گردند تا فرایند استنتاج بعدی را هدایت کنند. نتایج این پژوهش نشان داد که آگاه‌سازی مدل از بخش‌های اثرگذار در استنتاج‌های پیشین، می‌تواند دقت آن را در وظایف بازشناسی بهبود بخشد. با این حال، محدودیت کلیدی این روش در مکانیزم ترجمه‌ی نمرات به متن نهفته است؛ تبدیل نمرات انتساب در سطح توکن به زبان طبیعی، وابستگی شدیدی به قالب‌های متنی صلب و اکتشافی دارد که انعطاف‌پذیری سیستم را کاهش داده و تعمیم‌پذیری آن را به وظایف باز و متون آزاد محدود می‌سازد.

\paragraph{انتساب به داده‌های آموزشی در رژیم جعبه‌سیاه.}

علاوه بر انتساب تصمیمات مدل به مؤلفه‌ها و ویژگی‌های پرامپت ورودی، بُعد تکمیلی دیگری از انتساب پساعملی در سیستم‌های جعبه‌سیاه، ردیابی نقش و سهم داده‌های آموزشی در شکل‌گیری خروجی است. در حالی که روش‌های کلاسیک نفوذ داده‌ها (نظیر توابع نفوذ واکاوی‌شده در بخش~\ref{sec:3-1-2-4}) با اتکا به فرضیات جعبه‌سفید، نیازمند دسترسی کامل به گرادیان‌ها هستند، انسداد کامل مدل‌های تجاری به‌کارگیری آن‌ها را ناممکن می‌سازد. برای غلبه بر این محدودیت، آیوالیس و همکاران (۲۰۲۵)~\cite{aivalis2025enhancing} رویکردی متکی بر تحلیل نفوذ مبتنی بر جستجو را به‌عنوان یک روش مستقل از مدل برای هوش مصنوعی مولد معرفی کردند. آن‌ها با اتخاذ دیدگاه تفسیرپذیری مشاهده‌ای\enfootnote{observational interpretability}، مسئله را در یک چارچوب دومرحله‌ای فرمول‌بندی می‌کنند: نخست بازیابی نمونه‌های آموزشی مرتبط بر اساس پرامپت ورودی، و سپس تحلیل تطبیقی خروجی تولیدشده با نمونه‌های بازیابی‌شده با بهره‌گیری از معیارهای شباهت در سطح داده‌ی خام و فضای تعبیه. این متدولوژی نشان می‌دهد که حتی در شرایط انسداد مطلق پارامترها و در غیاب گرادیان‌های درونی، می‌توان بدون نیاز به آموزش مجدد مدل، زیرمجموعه‌های اثرگذار داده‌های آموزشی را در تولیدات مدل شناسایی و ممیزی کرد.


\subsubsection{یادگیری مدل‌های جانشین مفسر}\label{sec:3-1-3-3}

با آشکار شدن بن‌بست‌های محاسباتی و نارسایی‌های توزیعی در روش‌های اختلالی زمان استنتاج، جریان دیگری از پژوهش‌ها به سمت استفاده از مدل‌های جانشین متمایل گردید~\cite{guidotti2018survey, danilevsky2020survey}. در این پارادایم، تلاش می‌شود به‌جای اعمال تقریب‌های محلی و پرهزینه به ازای هر ورودی مجزا، یک مدل ثانویه آموزش ببیند تا منطق شبکه‌ی اصلی را در قالبی ساختاریافته‌تر استخراج و تقریب بزند.

چن و همکاران (۲۰۱۸) با ارائه‌ی چارچوب \lr{L2X}\enfootnote{learning to explain}، این استراتژی را با انتقال بار محاسباتی از فاز استنتاج، به فاز آموزش یک شبکه‌ی مفسر مجزا پایه‌گذاری کردند~\cite{chen2018learning}. این پژوهش، پارادایم آموزش سراسری یک مفسر محلی\enfootnote{globally-trained local interpreter} را شکل داد؛ به این معنا که برخلاف روش‌هایی نظیر \lr{LIME} که در آن برای هر نمونه‌ی ورودی یک مدل خطی محلی و ایزوله آموزش می‌بیند، \lr{L2X} یک مفسر واحد را به‌صورت سراسری روی کل توزیع داده‌ها بهینه‌سازی می‌کند تا برای هر نمونه، توضیحات محلی ارائه دهد. در این متدولوژی که مشخصاً برای وظایف دسته‌بندی (نظیر تحلیل احساسات متون یا دسته‌بندی تصاویر) طراحی شده است، هدف استخراج یک زیرمجموعه‌ی کوچک و گسسته از ویژگی‌های ورودی ($X_S$) با طول ثابت $k$ است که بیشترین اطلاعات را درباره‌ی خروجی مدل جعبه‌سیاه ($Y$) در بر داشته باشند. به بیان ریاضی، شبکه‌ی مفسر ($\mathcal{IN}$) آموزش می‌بیند تا اطلاعات متقابل میان ویژگی‌های انتخاب‌شده و متغیر پاسخ را بیشینه سازد. از آنجا که بهینه‌سازی مستقیم اطلاعات متقابل دشوار است، \lr{L2X} از یک کران پایین متغیر بهره برده و یک خانواده از توزیع‌های شرطی تقریب‌زننده ($\mathbb{Q}_S(Y \mid X_S)$) را معرفی می‌کند که وظیفه دارند احتمال وقوع خروجی $Y$ را صرفاً بر اساس زیربردار منتخب ($X_S$) تخمین بزنند:
\begin{equation}\label{eq:chap3_l2x_opt}
    \max_{\mathcal{IN}, \mathbb{Q}} \mathbb{E} \big[ \log \mathbb{Q}_S(Y \mid X_S) \big] \quad \text{\lr{s.t.}} \quad S \sim \mathcal{IN}(X)
\end{equation}
شهود نظری نهفته در این فرمول ریشه در مبانی نظریه‌ی اطلاعات دارد. بر اساس تعریف، اطلاعات متقابل به‌صورت $I(X_S; Y) = H(Y) - H(Y \mid X_S)$ بیان می‌شود. از آنجا که آنتروپی تصمیم‌های مدل ($H(Y)$) مقداری ثابت نسبت به مفسر است، بیشینه‌سازی اطلاعات متقابل دقیقاً هم‌ارز با کمینه‌سازی آنتروپی شرطی $H(Y \mid X_S)$ خواهد بود؛ این به معنای به حداقل رساندن جهل و نااطمینانی نسبت به خروجی $Y$ در شرایطی است که ناظر تنها زیرمجموعه‌ی فشرده‌ی $X_S$ را مشاهده کند. از آنجا که توزیع شرطی واقعی خروجی نامعلوم است، درست‌نمایی شرطی مدل جانشین ($\mathbb{E}[\log \mathbb{Q}_S(Y \mid X_S)]$) یک کران پایین متغیر و محاسبات‌پذیر بر اطلاعات متقابل فراهم می‌سازد. از منظر شهودی، این سازوکار مانند یک بازی همکارانه میان دو مؤلفه است: مفسر $\mathcal{IN}$ با نگاشت فضای ویژگی به توزیعی روی زیرمجموعه‌ها، فشرده‌ترین و گویاترین کلمات ورودی را استخراج می‌کند، و مدل جانشین $\mathbb{Q}$ می‌کوشد صرفاً بر مبنای همین کلمات برگزیده، تصمیم مدل جعبه‌سیاه را با بیشترین درست‌نمایی بازسازی کند. برای عبور از مشکل غیرمشتق‌پذیر بودن انتخاب گسسته‌ی ویژگی‌ها در شبکه‌ی $\mathcal{IN}$، این چارچوب از تقریب پیوسته‌ی گامبل-سافت‌مکس استفاده می‌کند تا توزیع احتمال انتخاب کلمات را طی فرایند پس‌انتشار به‌صورت پیوسته آموزش دهد.

اگرچه \lr{L2X} با پایه‌گذاری این چارچوب، هزینه‌ی استخراج ویژگی‌ها را در زمان استنتاج به‌شدت کاهش داد، اما فرمول‌بندی آن برای فضای خروجی محدود (مانند چند کلاس مشخص) بهینه‌سازی شده است. در مواجهه با مدل‌های زبانی مولد، خروجی ($Y$) یک توالی از کلمات با طول متغیر است؛ تقریب توزیع شرطی تقریب‌زننده و محاسبه‌ی اطلاعات متقابل در چنین فضای عظیمی (با پیچیدگی نمایانی از دایره‌ی واژگان)، از نظر محاسباتی غیرقابل‌حل است و اعمال آن برای انتخاب کلمات، به‌دلیل ایجاد حفره در پرامپت، مجدداً به تولید متون تکه‌تکه و بروز خطای شوک توزیعی منجر می‌شود.

با درک محدودیت‌های استفاده از اطلاعات متقابل برای خروجی‌های ساختاریافته (نظیر تولید توالی یا برچسب‌گذاری توالی)، سیدصالحی و همکاران (۲۰۲۴) چارچوب \lr{SOInter} را با الهام از منطق شبکه‌ی مفسر، اما مبتنی بر مدل‌های انرژی ارائه دادند~\cite{seyyedsalehi2024sointer} (پیشینه‌ی استفاده مدل‌های انرژی برای خروجی‌های ساختاریافته در بخش~\ref{sec:3-2-2} به‌تفصیل بررسی خواهد شد). این پژوهش با تکیه بر همان پارادایم آموزش سراسری مفسر محلی که در \lr{L2X} معرفی شد، ساختار جدیدی را برای خروجی‌های دارای وابستگی درونی پیشنهاد داد. در این روش، ابتدا یک شبکه‌ی انرژی سراسری ($\mathcal{E}_{sb}(x, y; \theta)$) بر روی تعداد زیادی از ورودی‌ها و خروجی‌های شبکه‌ی جعبه‌سیاه آموزش می‌بیند. این تابع پیوسته، یاد می‌گیرد که وابستگی‌های متقابل و میزان انطباق‌پذیری میان کل ورودی و کل ساختار خروجی را در سطح سراسری درک کند و سپس این دانش ساختاری برای هدایت مفسر محلی در یافتن بهینه‌ترین ویژگی‌ها مورد استفاده قرار می‌گیرد.

سپس، در فاز تفسیر و با تمرکز بر یک بُعد خاص از خروجی به‌عنوان هدف ($y_t^{sb}$)، یک تابع مفسر ($\mathcal{IN}_t(x; \alpha)$) مبتنی بر شبکه‌های عصبی و واحد گامبل-سافت‌مکس آموزش می‌بیند تا برداری \lr{$k$-hot} از مهم‌ترین ویژگی‌های ورودی را استخراج کند. تابع بهینه‌سازی این مفسر به دنبال کمینه‌سازی تفاوت انرژی میان پیش‌بینی جعبه‌سیاه روی ورودی مختل‌شده ($\tilde{y}$) و حالتی است که مقدار هدف با خروجی اصلی جعبه‌سیاه ($y_t^{sb}$) جایگزین شده باشد. معادله‌ی بهینه‌سازی مفسر به فرم زیر تعریف می‌شود:
\begin{equation}\label{eq:chap3_sointer}
    \arg\min_{\alpha} \mathbb{E}_{p(x)} \Big[ \max \big\{0, \mathcal{E}_{sb}(x \odot \mathcal{IN}_t(x; \alpha), [y_t^{sb}, \tilde{y}_{-t}]; \theta) - \mathcal{E}_{sb}(x \odot \mathcal{IN}_t(x; \alpha), \tilde{y}; \theta) \big\} \Big]
\end{equation}
در این رابطه‌ی بهینه‌سازی، عبارت $x \odot \mathcal{IN}_t$ نشان‌دهنده‌ی ورودی مختل‌شده‌ای است که در آن تنها ویژگی‌های منتخب مفسر باقی مانده و سایرین حذف شده‌اند. $\tilde{y}$ خروجی مدل برای این ورودی ناقص است و بردار $[y_t^{sb}, \tilde{y}_{-t}]$ نسخه‌ی اصلاح‌شده‌ای را نشان می‌دهد که در آن مقدار بُعد هدف به اجبار با خروجی اصلی مدل جایگزین شده است.
شهود فیزیکی این معادله از طریق مفهوم انرژی قابل‌تبیین است: در مدل‌های انرژی، مقدار انرژی کمتر بیانگر انطباق‌پذیری بالاتر میان ورودی و خروجی است و پیش‌بینی بهینه‌ی جعبه‌سیاه در نقطه‌ی کمینه‌ی انرژی رخ می‌دهد.
در این رابطه، $\tilde{y}$ پیش‌بینی واقعی مدل برای ورودی ناقص است؛ بنابراین مقدار انرژی $\mathcal{E}_{sb}(\dots, \tilde{y})$ ذاتاً کمینه‌ی انرژی ممکن برای آن ورودی است. در مقابل، بردار دستکاری‌شده‌ی $[y_t^{sb}, \tilde{y}_{-t}]$ خروجی همان ورودی ناقص است که در آن مقدار مؤلفه‌ی $t$ به پاسخ اولیه‌ی جعبه‌سیاه ($y_t^{sb}$) مقید شده است.
شهود محوری در این نکته نهفته است: اگر مفسر ویژگی‌های واقعاً تعیین‌کننده را برای بُعد $t$ در ورودی نگه داشته باشد، جعبه‌سیاه حتی با وجود حذف سایر کلمات، کماکان همان پاسخ اصلی را برای بُعد هدف پیش‌بینی می‌کند ($\tilde{y}_t = y_t^{sb}$). در این شرایط مطلوب، دو بردار یکسان شده و تفاوت انرژی آن‌ها دقیقاً به صفر می‌رسد. اما چنانچه مفسر ویژگی‌های بی‌اثر را برگزیده باشد، جعبه‌سیاه برای بُعد هدف مقدار دیگری تولید می‌کند؛ در این حالت، تحمیل دستی مقدار $y_t^{sb}$ موجب ناسازگاری ساختاری شده و انرژی آن نسبت به $\tilde{y}$ جهش می‌یابد. تابع زیان مفصلی با ساختار $\max\{0, \Delta \mathcal{E}\}$ دقیقاً این شکاف و جریمه‌ی انرژی را خنثی می‌سازد. در نتیجه، مفسر وادار می‌شود زیرمجموعه‌ای از ویژگی‌ها را گزینش کند که پیش‌بینی بُعد هدف را بدون وارد کردن شوک انرژی به سایر ابعاد خروجی پایدار نگه دارد.

این فرمول‌بندی نشان می‌دهد که در نظر گرفتن اطلاعات ساختاری و وابستگی میان متغیرهای خروجی حین آموزش مفسر با یک شبکه‌ی انرژی سراسری، انعطاف‌پذیری و دقت تفسیر را برای خروجی‌های دارای وابستگی درونی به‌شدت افزایش می‌دهد. با این حال، تعمیم مستقیمِ \lr{SOInter} به مدل‌های تولید متن باز با چالش روبه‌روست. برخلاف مسائلی نظیر برچسب‌گذاری توالی که ورودی و خروجی طول متناظر و ساختار هندسی شفافی دارند و این جایگذاری را ممکن می‌سازند، تولید متن یک فرایند متغیر و بدون قالب است. در نتیجه، تعریف یک تابع انرژی واحد که انطباق‌پذیری یک پرامپت مخدوش را با یک متن آزاد با طول متغیر بسنجد، نیازمند تغییر معماری روش است.

\normalfootnotes{}
برتری تئوریک رویکردهای مبتنی بر جانشین‌های سراسری نسبت به روش‌های محلی، در پژوهش لی و همکاران (۲۰۲۵)~\cite{li2025fromexplainability} نیز مورد تأیید قرار گرفت. آن‌ها در حوزه‌ی یادگیری تقویتی، چارچوبی به‌نام \lr{SILVER}\enfootnote{Shapley value-based interpretable policy via explanation regression} ارائه دادند تا شکاف میان توضیح‌پذیری محلی و تفسیرپذیری سراسری را پر کنند. این روش بر پایه‌ی مقادیر شاپلی بنا شده است؛ با این تفاوت که به‌جای استفاده‌ی ایزوله از مقادیر شاپلی برای توجیه یک تصمیم واحد، این مقادیر را به‌عنوان سیگنال‌های نظارتی در یک فرایند رگرسیون به‌کار می‌گیرد. به بیان دقیق‌تر، \lr{SILVER} با محاسبه‌ی مقادیر شاپلی در سراسر فضای حالات، یک مجموعه‌داده‌ی غنی از منطق تصمیم‌گیری مدل جعبه‌سیاه می‌سازد و سپس یک مدل جانشینِ ذاتاً شفاف (مانند درختان تصمیم) را روی این توضیحات آموزش می‌دهد. نویسندگان اثبات کردند که این انتقال از سیگنال‌های محلی (که غالباً پرنویز و ناپایدارند) به یک سیاست سراسری، نه تنها عملکرد اصلی سیستم را حفظ می‌کند، بلکه درک مستحکم‌تر و پایدارتری از رفتار عامل ارائه می‌دهد. این یافته ثابت می‌کند که تجمیع توضیحات محلی در قالب یک مدل جانشین سراسری، به‌شکل بنیادینی از محدودیت‌های تقریب‌های نقطه‌ای مصون است.
\twocolumnfootnotes{}

در مسیر تعمیم مدل‌های جانشین برای استخراج منطق در دامنه‌ی تولید متون طولانی، رویکردهایی نظیر چارچوب \lr{TriSum}~\cite{jiang2024trisum} توسعه یافتند. جیانگ و همکاران در این پژوهش با بهره‌گیری از تقطیر دانش، یک \lr{LLM} عظیم را به‌عنوان مدل معلم برای استخراج توجیهات متنی از داده‌ها به‌کار گرفتند. در وظیفه‌ی خلاصه‌سازی متن، مدل معلم استدلال‌های خود را به‌صورت صریح در قالب سه‌گانه‌های معنایی و ساختاریافته (شامل جنبه-موجودیت-رابطه) فرمول‌بندی می‌کند. سپس، یک مدل زبانی کوچک‌تر و کارآمدتر (جانشین یا دانش‌آموز) آموزش می‌بیند تا معماری استنتاج خود را تغییر دهد؛ این مدل پیش از تولید خلاصه‌ی نهایی، موظف است این سه‌گانه‌های استدلالی را تولید کند. استفاده از این قالب، خروجی‌های مدل جانشین را به‌شدت شفاف و منطق گزاره‌ای آن را برای ناظر انسانی قابل‌ردیابی می‌سازد.

با این وجود، ارزیابی معماری‌هایی نظیر \lr{TriSum} به‌عنوان ابزار ممیزی، با یک نقد نظری روبه‌روست. همان‌گونه که در بررسی رویکردهای جعبه‌سفید واکاوی شد، ساده‌سازی مدل‌های پیچیده و تقطیر دانش آن‌ها به جانشین‌های تقلیل‌یافته، خطر بروز توهم تفسیرپذیری و عدم‌تطابق سازوکارهای محاسباتی درونی را در مواجهه با داده‌های خارج از توزیع به‌همراه دارد~\cite{friedman2024interpretability}. این واگرایی در جانشین‌های مولد به‌مراتب شدیدتر است؛ چرا که فضای وسیع تولید متن آزاد به مدل جانشین امکان می‌دهد تا توجیهاتی (نظیر سه‌گانه‌های معنایی) تولید کند که صرفاً برای گراف محاسباتی خودش معتبرند، نه بازتاب‌دهنده‌ی سازوکار مدل جعبه‌سیاهِ اصلی. این چالش نشان می‌دهد که استقرار مدل‌های جانشینِ وفادار نیازمند تدابیر ساختاری ویژه‌ای است؛ به‌گونه‌ای که سیستم مفسر ترجیحاً از تولید متن آزاد نامحدود پرهیز کرده و با بهره‌گیری از تمهیداتی نظیر بهینه‌سازی انطباقی، به‌طور پیوسته خود را با پویایی مدل اصلی هم‌تراز سازد. ضرورت احتراز از این توهم تبیینی و نیاز به سنجش وفاداری مفسرها، مسیر پژوهش‌ها را به‌سوی بازنگری بنیادین در معیارهای ارزیابی سوق می‌دهد.


\subsubsection{بحران ارزیابی و فقدان حقیقت پایه}\label{sec:3-1-3-4}

از آنجا که در سیستم‌های جعبه‌سیاه به گراف محاسباتی و متغیرهای درونی مدل دسترسی مستقیم وجود ندارد، چگونه می‌توان ادعای وفاداری این الگوریتم‌های مفسر را راستی‌آزمایی کرد؟ در این بخش نشان داده می‌شود که متدولوژی‌های ارزیابی موجود به‌دلیل ماهیت منعطف زبان و فقدان حقیقت پایه‌ی عینی، با چالش‌های نظری عمیقی مواجه هستند و در ارزیابیِ اتکاپذیری انتساب‌ها دچار تله‌های روش‌شناختی و قانون گودهارت می‌شوند.

\paragraph{معیارهای فرسایشی و تله‌های منطقی در ارزیابی انتساب.}

برای استانداردسازی ارزیابی‌ها در پردازش زبان طبیعی، دی‌یانگ و همکاران (۲۰۲۰)~\cite{deyoung2020eraser} محک جامع \lr{ERASER} را معرفی کردند. این پژوهش دو معیار اصلی برای سنجش وفاداری ارائه داد که شامل جامعیت\enfootnote{comprehensiveness} و کفایت است. با فرض آنکه $x$ متن ورودی، $r$ توجیه استخراج‌شده توسط مفسر (زیرمجموعه‌ای از توکن‌ها) و $f(x)$ تابع پیش‌بینی مدل باشد، این دو معیار به‌صورت زیر محاسبه می‌شوند:
\begin{align}
    \text{\lr{Comprehensiveness}} &= f(x) - f(x \setminus r) \label{eq:chap3_comprehensiveness} \\
    \text{\lr{Sufficiency}} &= f(x) - f(r) \label{eq:chap3_sufficiency}
\end{align}
شهود منطقی حاکم بر این دو رابطه، بر پایه‌ی دو سنجه‌ی متضاد و مکمل برای راستی‌آزمایی توجیهات بنا شده است:
۱. معیار جامعیت (معادله‌ی~\ref{eq:chap3_comprehensiveness}): این سنجه پاسخ به این پرسش است که «آیا تمام شواهد کلیدی تصمیم در توجیه استخراج‌شده ($r$) گردآوری شده‌اند؟» اگر کلمات موجود در $r$ به‌راستی همه‌ی علل اصلی تصمیم باشند، حذف آن‌ها از بافتار ورودی ($x \setminus r$) باید شبکه را از اطلاعات حیاتی محروم ساخته و منجر به سقوط شدید احتمال پیش‌بینی ($f(x \setminus r)$) شود. بنابراین، هر اندازه تفاضل $f(x) - f(x \setminus r)$ بزرگ‌تر و مثبت‌تر باشد، نشان‌دهنده‌ی جامعیت بالاتر توجیه در پوشش تمام دلایل تصمیم است.
۲. معیار کفایت (معادله‌ی~\ref{eq:chap3_sufficiency}): این سنجه به دنبال پاسخ به این پرسش است که «آیا توجیه برگزیده ($r$) به‌تنهایی و بدون اتکا به سایر کلمات بافتار، برای متقاعد ساختن مدل جهت اتخاذ همان تصمیم کفایت می‌کند؟» در صورتی که مدل صرفاً با دریافت زیرمجموعه‌ی $r$ بتواند با همان اطمینان اولیه پیش‌بینی قبلی را بازتولید کند ($f(r) \approx f(x)$)، تفاضل $f(x) - f(r)$ به صفر نزدیک خواهد شد. از این رو، مقادیر کوچک‌تر و نزدیک به صفر (و حتی منفی در صورت افزایش اطمینان) بیانگر کفایت بالاتر توجیه انتخابی هستند.
افزون بر این، برای ارزیابی مفسرهایی که نمرات اهمیت پیوسته تولید می‌کنند، میانگین این تغییرات احتمالاتی در طول منحنی حذف تدریجی توکن‌ها در قالب معیار سطح روی منحنی اختلال یا \lr{AOPC}\enfootnote{area over the perturbation curve} محاسبه می‌گردد.

با این وجود، استفاده از این معیارهای مبتنی بر اختلال با نقدهای بنیادینی مواجه شد. جو و همکاران (۲۰۲۲)~\cite{ju2022logictraps} با واکاوی دقیق متدولوژی‌های ارزیابی، از وجود سه تله‌ی منطقی پرده برداشتند که اعتبار ارزیابی‌ها را مخدوش می‌کند.

نخستین تله‌ی منطقی، در استفاده از توضیحات انسانی به‌عنوان حقیقت پایه نهفته است. شبکه‌های عصبی غالباً بر اساس همبستگی‌های آماری تصمیم می‌گیرند که با منطق انسان متفاوت است. جو و همکاران با یک آزمایش کنترل‌شده نشان دادند که در مجموعه‌داده‌ی پرسش‌وپاسخ \lr{RACE}~\cite{lai2017race}، یک مدل درک مطلب حتی با حذف کامل صورت سوال، همچنان در ۶۴ درصد از مواقع پیش‌بینی درست خود را با کمترین افت در سطح اطمینان حفظ می‌کند. این پدیده ثابت می‌کند که مکانیسم تصمیم‌گیری مدل با استدلال انسانی تفاوت ماهوی دارد؛ در نتیجه، عدم‌تطابق خروجی یک مفسر با منطق انسانی، لزوماً به معنای عدم‌وفاداری آن مفسر به گراف محاسباتی شبکه‌ی عصبی نیست.

تله‌ی دوم به استفاده از رویکردهای مبتنی بر اختلال معنادار برای مقایسه و امتیازدهی به مفسرها بازمی‌گردد. معیارهایی نظیر جامعیت و \lr{AOPC}، خود ماهیتاً یک روش انتساب اعتبار هستند که بر اساس اختلال عمل می‌کنند. پژوهش‌ها اثبات کردند که نتیجه‌ی ارزیابی در این معیارها کاملاً به استراتژی اختلال (حذف مستقیم یا جایگزینی با نمونه‌برداری) بستگی دارد و قابل‌دستکاری است. به‌عنوان نمونه، زمانی که استراتژی اختلال در \lr{AOPC} بر پایه‌ی حذف کامل توکن‌ها تنظیم شود، مفسر کلاسیک حذف یکی\enfootnote{leave-one-out} بهترین عملکرد را در محک نشان می‌دهد. اما با تغییر استراتژی ارزیابی به جایگزینی توکن‌ها با استفاده از توزیع احتمالات یک مدل زبانی، مفسری نظیر \lr{Marg}~\cite{kim2020interpretation} که خود بر اساس همین منطق حاشیه‌سازی توسعه یافته است، پیروز ارزیابی می‌شود. این مسئله نشان می‌دهد که معیارهای یادشده صرفاً میزان شباهت رفتار مفسر به استراتژی اختلال خود را می‌سنجند، نه وفاداری واقعی به مدل جعبه‌سیاه را.

تله‌ی سوم در ارزیابی استحکام مفسرها از طریق بررسی ثبات نمرات انتساب در نمونه‌های متخاصم یا مشابه معنایی رخ می‌دهد. مسئله‌ی محوری این است که این رویکردها فرض می‌کنند «اگر دو ورودی از نظر معنایی مشابه باشند و خروجی یکسانی از مدل دریافت کنند، مفسر نیز باید نمرات انتساب مشابهی به آن‌ها بدهد» و در صورت تغییر نمرات، الگوریتم تفسیر را شکننده و نامعتبر می‌خوانند. اما جو و همکاران~\cite{ju2022logictraps} با استدلال‌های دقیق تجربی اثبات کردند که این پیش‌فرض از اساس مغالطه‌آمیز است. تغییر در نمرات انتساب ممکن است به این دلیل باشد که فرایند استدلال مدل جعبه‌سیاه واقعاً تغییر کرده است، نه اینکه روش تفسیر غیرقابل‌اعتماد باشد. آن‌ها با طراحی یک چارچوب دومرحله‌ای (شامل استخراج‌کننده و دسته‌بند) نشان دادند که شبکه‌های عصبی در مواجهه با نمونه‌های متخاصمی که خروجی یکسانی با نمونه‌ی اصلی دارند، ممکن است به زیرمجموعه‌های کاملاً متفاوتی از کلمات توجه کنند (که در برخی موارد صفر درصد با کلمات استدلالی جمله‌ی اصلی اشتراک دارند). بنابراین، تغییر در خروجی مفسر نشانه‌ی نقص آن نیست، بلکه بازتاب وفادارانه‌ی تغییر رفتار و منطق پنهان مدل در مواجهه با ورودی جدید است. این تله‌ی منطقی باعث می‌شود مفسرهایی که تغییرات استدلال مدل را به‌درستی بازتاب می‌دهند، در سیستم‌های ارزیابی موجود به اشتباه تنبیه شوند.

\paragraph{بحران داده‌های خارج از توزیع و قانون گودهارت.}

ریشه‌ی اصلی تله‌های ارزیابی مبتنی بر اختلال، در بروز خطای داده‌های \lr{OOD} قرار دارد. مدسن و همکاران (۲۰۲۲)~\cite{madsen2022evaluating} با فرمول‌بندی دقیق این مشکل نشان دادند که اعمال ماسک و حذف کلمات کلیدی، گرامر جمله را متلاشی می‌کند. در نتیجه، افت احتمال مدل ناشی از قطع زنجیره‌ی استنتاج نیست، بلکه واکنش آماری شبکه به یک متن غیرطبیعی و مخدوش است. آن‌ها استدلال کردند که برای رفع این مشکل، تنها راه نظری استفاده از رویکرد روش حذف و آموزش مجدد\enfootnote{remove and retrain} (\lr{ROAR}) است تا مدل روی مجموعه‌داده‌ای از متون مخدوش دوباره آموزش ببیند و افت عملکرد صرفاً ناشی از فقدان اطلاعات باشد. با این حال، اجرای این راهکار برای مدل‌های زبانی عظیم جعبه‌سیاه غیرعملی است و حتی در صورت انجام، مدلی که بازآموزی شود، دیگر همان مدل اولیه‌ی مورد ارزیابی نخواهد بود.

این شکنندگی ساختاری تنها به روش‌های انتساب توکن محدود نمی‌شود. آتاناسووا و همکاران (۲۰۲۳)~\cite{atanasova2023faithfulness} با طراحی آزمون‌های وفاداری برای رویکردهای مبتنی بر تبیین‌های مبتنی بر زبان طبیعی\enfootnote{natural language explanations} (\lr{NLE})، فرضیه‌ی درون‌نگری مدل‌ها را به چالش کشیدند. آن‌ها با اعمال اختلالات جزئی و بی‌اثر در ورودی (تغییراتی که پیش‌بینی نهایی مدل را ثابت نگه می‌دارد)، نشان دادند که توجیهات متنی تولیدشده توسط مدل به‌شدت تغییر می‌کنند. این ناپایداری تأیید می‌کند که گزاره‌های خودتوضیحی نیز لزوماً بازتاب سازوکار درونی تصمیم‌گیری نیستند و ممکن است صرفاً توجیهاتی پساعملی به شمار آیند.

اتکا به این معیارهای ارزیابی دارای نقص، پیامد مخربی در جهت‌دهی به پژوهش‌ها به‌همراه داشته است. هسیا و همکاران (۲۰۲۴)~\cite{hsia2024goodhart} با بررسی محک‌های تفسیرپذیری اثبات کردند که «قانون گودهارت» به‌طور کامل بر این حوزه سایه افکنده است. بر اساس این قانون در آمار، وقتی یک معیار (نظیر نمره‌ی جامعیت) به هدفِ بهینه‌سازی تبدیل شود، اعتبار خود را از دست می‌دهد. هنگامی که پژوهشگران مدل‌های جانشینِ مفسر را به‌گونه‌ای آموزش می‌دهند تا مقادیر جامعیت و کفایت را مستقیماً بیشینه کنند، این مدل‌ها یاد می‌گیرند که محک را هک کنند. مفسر جانشین به‌جای کشف الگوی اثرگذاری واقعی ویژگی‌ها، کلماتی را استخراج و حذف می‌کند که بیشترین شوک توزیعی (یعنی رخداد خارج از توزیع) را به مدل جعبه‌سیاه وارد سازند. در نتیجه، نمره‌ی ارزیابی در محک به‌صورت کاذب به ۱۰۰ درصد نزدیک می‌شود، در حالی که این توکن‌های انتخاب‌شده اساساً فاقد ارزشِ معرفت‌شناختی و تفسیری هستند.

\paragraph{نتیجه‌گیری.}

مجموعه‌ی این پژوهش‌ها یک چالش روش‌شناختی بنیادین را در ارزیابی مفسرها آشکار می‌سازد. در ممیزی سیستم‌های جعبه‌سیاه، اگرچه استفاده از روش‌های مبتنی بر اختلال امری ناگزیر است، اما این روش‌ها همواره در معرض خطای ناشی از ورود به ناحیه‌ی داده‌های خارج از توزیع قرار دارند. از سوی دیگر، در مدل‌های شفاف و منبع‌باز، وجود دسترسی مستقیم به گرادیان‌ها و وزن‌ها امکان تعریف خط‌پایه‌های مقایسه‌ای معتبر و دقیقی (نظیر روش‌های جعبه‌سفید) را فراهم می‌آورد؛ هرچند این روش‌ها نیز فرضیات و ویژگی‌های خاص خود را دارند. با این حال، در غیاب دسترسی به گراف محاسباتی در سیستم‌های تجاری مسدود، سنجش وفاداری توضیحات بدون دستکاری مخرب ورودی و بدون نقض اصول ارزیابی، همچنان به‌عنوان یکی از مسائل باز و چالش‌های حل‌نشده در این پارادایم مطرح است.



% =========================================================
% بدنه دوم
% =========================================================



\section{پایه‌های نظری چارچوب پیشنهادی}\label{sec:3-2}

با آشکار شدن بن‌بست‌های محاسباتی و نارسایی‌های معنایی در متدولوژی‌های رایج جعبه‌سیاه، ضرورت گذار به رهیافتی نوین برای ارزیابی خروجی \lr{LLM}ها اجتناب‌ناپذیر می‌نماید. در پاسخ به این ضرورت، بخش حاضر به مرور پیشینه‌ی ابزارهایی می‌پردازد که شالوده‌ی محاسباتی چارچوب پیشنهادی این پژوهش را شکل می‌دهند. بدین منظور، واکاوی ادبیات در این بخش در سه گام متوالی ساختاردهی شده است: نخست، پیشینه‌ی ارتقای واحد تحلیل از سطح خُرد به سطح جمله با هدف حفظ یکپارچگی گزاره‌ای تبیین می‌شود؛ سپس، مبانی مدل‌های مبتنی بر انرژی به‌عنوان ارزیابان یکپارچه در مسائل ساختاریافته معرفی می‌گردند؛ و در نهایت، سیر تطبیق این مدل‌های کلان‌نگر در حوزه‌ی زبان طبیعی مورد تحلیل قرار می‌گیرد.


\subsection{ارتقای واحد معنایی به جمله}\label{sec:3-2-1}

همان‌طور که در بخش‌های پیشین مشاهده شد، تقلیل‌گرایی به سطح توکن یکی از ریشه‌های بنیادین ناکامی در تحلیل خروجی مدل‌های زبانی است. زبان طبیعی برخلاف تصاویر، ماهیتی گسسته و ترکیبیاتی دارد که در آن معنا نه در نشانه‌های منفرد، بلکه در ساختار گزاره‌ای و شبکه‌ی روابط میان آن‌ها شکل می‌گیرد. این دیدگاه که ریشه در اصل بافتار فرگه~\cite{frege1991foundations} دارد، بیان می‌کند که کلمات تنها در بستر یک گزاره‌ی کامل دارای معنای قطعی هستند. بر این اساس، ارتقای واحد تحلیل از توکن به جمله، به‌عنوان یک راهکار پایدار و منطقی برای حفظ ساختار گرامری و معنایی متن مطرح می‌شود؛ رویکردی که ریشه‌های مستحکمی در تکامل ادبیات پردازش زبان طبیعی و علوم شناختی دارد.


\paragraph{پیشینه‌ی بازنمایی جملات و استخراج مفاهیم.}

در تاریخچه‌ی پردازش زبان طبیعی، در نظر گرفتن جملات به‌عنوان اشیای معنایی مستقل، سیر تکاملی روشنی داشته است. در معماری‌های پایه‌ای نظیر \lr{BERT}~\cite{devlin2019bert}، برای درک روابط منطقیِ کلان میان بخش‌های مختلف یک متن، وظیفه‌ی پیش‌بینی جمله‌ی بعدی\enfootnote{next sentence prediction} تعریف شد. در این مکانیزم پیش‌آموزش، مدل به‌جای پردازش منفرد کلمات، جفت‌جملاتی به فرم $(A, B)$ را دریافت می‌کند و باید به‌عنوان یک مسئله‌ی دسته‌بندی دوتایی تشخیص دهد که آیا جمله $B$ در متن اصلی بلافاصله پس از جمله $A$ آمده است یا یک جمله‌ی تصادفی است. مدل این فرایند را با استفاده از بردار خروجی توکن ویژه \lr{[CLS]} (با نماد $C \in \mathbb{R}^H$) و عبور آن از یک لایه‌ی خطی استنتاجی با ماتریس وزن $W \in \mathbb{R}^{2 \times H}$ انجام می‌دهد تا احتمال وقوع توالی، یعنی $P(\text{\lr{IsNext}} \mid A, B) = \text{\lr{Softmax}}(CW^\top)$، را محاسبه کند. از منظر شهودی، توکن ویژه‌ی \lr{[CLS]} در طول لایه‌های خودتوجهی به‌عنوان یک خازن اطلاعاتی عمل می‌کند که برهم‌کنش‌های معناییِ متقابل میان تمامی کلمات دو جمله را در بردار فشرده‌ی $C$ خلاصه می‌نماید. حاصل‌ضرب خطی $CW^\top$ این چکیده‌ی معنایی را به دو نمره‌ی خام (لوجیت) برای دو کلاس متوالی و نامرتبط نگاشت کرده و تابع بیشینه‌ساز هموار (\lr{Softmax}) آن را به یک توزیع احتمالاتی تبدیل می‌کند. این فرمول نشان داد که یادگیری هم‌رخدادی کلمات به‌تنهایی کافی نیست و مدل حتی برای تصمیم‌گیری‌های کلان در سطح متن و درک مفاهیمی نظیر استلزام، نیازمند متراکم‌سازی معنای یک گزاره‌ی کامل در یک بازنمایی واحد است.

در امتداد همین نیاز، معماری \lr{Sentence-BERT}~\cite{reimers2019sentencebert} توسعه یافت. همان‌طور که فرمول‌بندی و مکانیزم تجمیع این معماری پیش‌تر در بخش~\ref{sec:2-3-4} بررسی شد، شبکه‌های سنتی برای یافتن میزان شباهت دو جمله، نیازمند پردازش هم‌زمان با بار محاسباتی سنگین بودند. مدل \lr{Sentence-BERT} با بهره‌گیری از شبکه‌های سیامی، ثابت کرد که می‌توان بار معنایی یک گزاره‌ی کامل را به‌عنوان یک واحد محاسباتیِ مستقل و هندسی در یک فضای پیوسته فرمول‌بندی کرد تا ارزیابی قرابت معنایی جملات از طریق معیارهایی نظیر شباهت کسینوسی امکان‌پذیر گردد. به بیان شهودی، به‌جای مقایسه‌ی پرهزینه‌ی کلمه‌به‌کلمه در یک گراف محاسباتی حجیم، هر جمله به نقطه‌ای مستقل در یک فضای برداری پیوسته نگاشت می‌شود؛ جایی که نزدیکی هندسی دو بردار (فاصله‌ی زاویه‌ای یا کسینوسی)، مستقیماً بازتاب‌دهنده‌ی قرابت معنایی آن دو اندیشه است.

این گذار ساختاری از توکن به جمله، مرزهای مدل‌سازی موضوعی را نیز درنوردید. جامعه‌ی پردازش زبان به‌تدریج دریافت که استخراج مفاهیم واقعی متون، نیازمند عبور از آمار سطح کلمه است. برای نمونه، وانگ و همکاران (۲۰۲۳)~\cite{wang2023prompting} با معرفی چارچوب \lr{PromptTopic} نشان دادند که مدل‌های سنتی کشف موضوع که صرفاً بر تواتر و هم‌رخدادی توکن‌ها استوارند، در استخراج مفاهیم از استدلال‌های پیچیده و متون کوتاه شکست می‌خورند. آن‌ها برای غلبه بر این محدودیت، استخراج مفاهیم را از سطح توکن به سطح جمله ارتقا دادند؛ در این متدولوژی، از ظرفیت \lr{LLM}ها استفاده می‌شود تا موضوعات را از گزاره‌های متنی استخراج کرده و سپس این مفاهیم جمله‌محور را با تکنیک‌هایی نظیر \lr{c-TF-IDF} پالایش کنند. این رویکردها تأیید می‌کنند که در پیشرفته‌ترین متدولوژی‌های \lr{NLP} نیز، گره‌گشایی از معنای متن نیازمند دست‌کشیدن از توکن‌ها و اتکا به سطح جمله است.


\paragraph{بازتعریف واحد استدلال در مدل‌های مفهومی.}

تلاش برای تطبیق مکانیزم‌ها با این ماهیت گزاره‌ای زبان، منجر به توسعه‌ی معماری‌هایی در سطح مدل‌های بنیادین گردید. همان‌طور که پیش‌تر در بخش تفسیرپذیری درونی (بخش~\ref{sec:3-1-1-3}) اشاره شد، تیم \lr{LCM} در شرکت متا (۲۰۲۴)~\cite{lcmteam2024lcm} با معرفی مدل‌های مفهومی بزرگ، پردازش متوالی در سطح توکن را به‌نفع استدلال در سطح جملات کنار گذاشتند. این مدل‌ها به‌جای محاسبه‌ی توزیع احتمال روی توکن‌های منفرد، متن را مستقیماً در یک فضای مفهومی پیوسته مدل کرده و فرایند تولید خودهمبسته را بر روی بردارهای سطح جمله ($\mathbf{C}_t \in \mathbb{R}^d$) و به فرم $P(\mathbf{C}_t \mid \mathbf{C}_{<t})$ انجام می‌دهند. از منظر شهودی، تفاوت بنیادین این مدل‌سازی با مدل‌های زبانی رایج در این است که در مدل‌های توکن‌محور، پیش‌بینی بر روی یک توزیع احتمالاتیِ گسسته از میان ده‌ها هزار کلمه‌ی مجزا انجام می‌شود ($P(w_t \mid w_{<t})$)؛ در حالی که در مدل‌های مفهومی، متغیر هدف $\mathbf{C}_t$ یک بردار پیوسته در فضای بازنمایی معنایی گزاره‌ها (نظیر فضای تعبیه \lr{SONAR}) است. بدین‌ترتیب، رابطه‌ی خودهمبسته‌ی $P(\mathbf{C}_t \mid \mathbf{C}_{<t})$ نمایانگر جریانی از ایده‌ها و افکار متوالی است؛ یعنی مدل در هر گام، چارچوب مفهومی گزاره‌ی بعدی را بر اساس توالی ایده‌های پیشین شکل می‌دهد و ترجمه‌ی این مفهوم به کلمات ملموس به یک ماژول رمزگشای جداگانه محول می‌گردد. اگرچه در بخش‌های پیشین محدودیت‌های این معماری مورد نقد قرار گرفت، اما شکل‌گیری این ساختار از زاویه‌ی پردازش زبان طبیعی حائز اهمیت است.

دلیل توسعه‌ی چنین معماری‌هایی آن است که تولید حریصانه و کلمه‌به‌کلمه، با شیوه‌ی استدلال سلسله‌مراتبی و ایده‌محور انسان در تضاد است. تغییر واحد پیش‌بینی از توکن به جمله در یک فضای بازنمایی، این واقعیت را برجسته می‌سازد که جمله، یک قالب طبیعی و کارآمد برای تقطیع یک مفهوم و فکر کامل در زبان طبیعی است. هنگامی که تحلیل و تولید در سطح جمله انجام می‌گیرد، نه تنها گرامر و ساختار نحوی متن از هم نمی‌پاشد، بلکه مفهوم منطقیِ نهفته در آن نیز برای انسان و ماشین به‌صورت یکپارچه حفظ می‌شود. این ارتقای واحد محاسباتی تضمین می‌کند که دستکاری یا تحلیل سیستم، ساختار گرامری را مخدوش نکرده و ارزیابی‌ها مستقیماً بر روی واحدهای حامل منطق انجام پذیرد. بنابراین، پذیرش جمله به‌عنوان واحد یا مفهوم پایه‌ی معنایی یک انتخاب سلیقه‌ای نیست؛ بلکه رویکردی است که پشتوانه‌ی مستحکمی در روند تکاملی معماری‌های هوش مصنوعی داشته و می‌تواند به‌عنوان بستری پایدار برای تحلیل‌های معنایی و استخراج استدلال‌ها مورد استفاده قرار گیرد.


\subsection{مدل‌های انرژی در پیش‌بینی ساختاریافته}\label{sec:3-2-2}

پیش‌بینی ساختاریافته به دسته‌ای از مسائل یادگیری ماشین اطلاق می‌شود که در آن‌ها متغیرهای خروجی دارای وابستگی‌های متقابل هستند (مانند برچسب‌گذاری پیکسل‌های یک تصویر یا تخصیص هم‌زمان چند برچسب به یک متن). در این مسائل، تصمیم‌گیری مستقل برای هر متغیر به نتایج متناقض منجر می‌شود. مدل‌های انرژی به‌دلیل توانایی در تخصیص یک امتیاز سراسری به کل ساختار خروجی، ابزاری کارآمد برای حل این وابستگی‌ها به شمار می‌آیند. این بخش به مرور سیر تکامل شبکه‌های انرژی در این دامنه می‌پردازد تا روشن سازد چرا مکانیسم‌های جستجو و استنتاج کلاسیک در رویارویی با خروجی‌های نامحدود زبان طبیعی در زمان تولید آزاد که همچنان دارای ساختاری درونی است، شکست می‌خورند.


\paragraph{شبکه‌های انرژی عمیق و چالش استنتاج در فضاهای گسسته.}

همان‌طور که در بخش~\ref{sec:2-4-3} بررسی شد، در پیشینه‌ی مدل‌های انرژی برای پیش‌بینی ساختاریافته، معماری «شبکه‌های انرژی پیش‌بینی ساختاریافته» (\lr{SPEN})~\cite{belanger2016spen} یک نقطه‌ی عطف محسوب می‌شود. پیش از این پژوهش، ارزیابی خروجی‌های پیچیده و ساختاریافته غالباً با استفاده از گراف‌های احتمالی (نظیر میدان‌های تصادفی مارکوف\enfootnote{Markov random field}) انجام می‌شد که برای حفظ حل‌پذیری ریاضی، ناگزیر به محدود کردن وابستگی‌ها در سطح همسایگی‌های محلی بودند. مدل \lr{SPEN} با استفاده از یک شبکه‌ی عصبی عمیق این محدودیت را کنار زد؛ به این صورت که شبکه، یک تابع انرژی ($E_w(x, y)$) را یاد می‌گیرد تا با یک نگاه سراسری به کل ورودی ($x$) و کل خروجی پیشنهادی ($y$)، میزان انطباق‌پذیری آن‌ها را به‌صورت یکپارچه داوری کند. شهود فیزیکی حاکم بر این فرمول‌بندی آن است که انرژی را همانند پتانسیل یک سیستم در نظر می‌گیرد؛ وضعیت‌های منطبق و منسجم خروجی در قعر دره‌های انرژی قرار دارند و پیش‌بینی‌های متناقض انرژی بالایی دریافت می‌کنند. با این وجود، این توانایی برتر، چالش‌های عملیاتی در مسئله‌ی استنتاج به‌همراه دارد. برای یافتن بهترین خروجی در فاز استنتاج، سیستم باید در فضای خروجی‌ها جستجو کرده و گزینه‌ی دارای کمترین انرژی ($\arg\min_y E_w(x, y)$) را استخراج کند. استنتاج برخلاف شبکه‌های پیش‌خور معمولی که خروجی را با یک عبور مستقیم تولید می‌کنند، به یک مسئله‌ی بهینه‌سازی پویا تبدیل می‌شود که در آن خروجی $y$ مانند ذره‌ای در یک چشم‌انداز ناهموار با گام‌های گرادیانی در جهت شیب تند ($\nabla_y E_w$) می‌لغزد تا پایدارترین پیکربندی حاصل شود؛ جستجویی که در غیاب ساختارهای گرافیِ ساده، نیازمند اجرای تکرارشونده‌ی الگوریتم گرادیان کاهشی مستقیماً بر روی متغیرهای خروجی ($y$) است.

آموزش توابع انرژی عمیق در ابتدا بسیار ناپایدار بود. برای حل این معضل، بلانجر و همکاران (۲۰۱۷)~\cite{belanger2017endtoend} روشی تحت عنوان آموزش انتها‌به‌انتها را برای \lr{SPEN} معرفی کردند. ایده‌ی کلیدی آن‌ها تغییر نگاه به فاز استنتاج بود. از آنجا که یافتن بهترین خروجی نیازمند اجرای ده‌ها گام گرادیان کاهشی است، آن‌ها این گام‌های متوالی را دقیقاً مانند لایه‌های یک شبکه‌ی عصبی در نظر گرفتند که اصطلاحاً باز\enfootnote{unroll} شده است. به عبارت ملموس‌تر، الگوریتم بهینه‌سازی به‌جای اینکه یک فرایند محاسباتیِ مجزا تلقی شود، مانند یک شبکه‌ی عصبی بازگشتی بسیار عمیق در نظر گرفته می‌شود که هر لایه‌ی آن یک گام گرادیانی است. با این مدل‌سازی، خطای خروجی نهایی در گام آخر محاسبه شده و گرادیان این خطا، مستقیماً از طریق تمام گام‌های بهینه‌سازی به عقب (پس‌انتشار) برمی‌گردد تا وزن‌های شبکه‌ی انرژی به‌روزرسانی شوند. این ظرافت ریاضی به مدل اجازه می‌دهد تا چشم‌انداز انرژی خود را به‌گونه‌ای نرم و هموار صیقل دهد که فرایند جستجو با گرادیان در زمان تست، بدون سرگردانی و با سرعت و دقت بسیار بالاتری به سمت پاسخ صحیح همگرا شود.

با این وجود، استفاده از گرادیان برای جستجو در فضای خروجی، پیش‌نیازهایی دارد که برای زبان طبیعی دردسرساز است. در بسیاری از وظایف ساختاریافته، خروجی ماهیتاً مجموعه‌ای از انتخاب‌های گسسته است (مثلاً در دسته‌بندی چندبرچسبی، خروجی یک بردار دودویی $y \in {\{0, 1\}}^L$ است که حضور یا عدم‌حضور هر برچسب را نشان می‌دهد). از آنجا که مشتق‌گیری از مقادیر صفر و یک ناممکن است، گیگلی و همکاران (۲۰۱۷)~\cite{gygli2017dvn} در توسعه‌ی «شبکه‌های ارزش عمیق»\enfootnote{deep value networks} (\lr{DVN}) از تکنیک آزادسازی پیوسته استفاده کردند. مدل \lr{DVN} برخلاف \lr{SPEN} که یک انرژی انتزاعی تولید می‌کرد، مستقیماً آموزش می‌بیند تا یک معیار ارزیابی واقعی (مانند امتیاز \lr{F1}) را برای خروجی پیش‌بینی کند. در فاز استنتاج، برای دور زدن مشکل مشتق‌ناپذیری، آن‌ها فضای گسسته‌ی خروجی را به یک فضای پیوسته ($y \in {[0, 1]}^L$) بسط دادند. شهود پشت آزادسازی پیوسته این است که متغیرهای دودویی را به‌عنوان احتمال حضور یا درجات عضویت پیوسته در بازه‌ی $[0, 1]$ تفسیر می‌کند؛ بدین‌ترتیب، فضای تصمیم‌گیری از یک شبکه‌ی گسسته‌ی دندانه‌دار و مشتق‌ناپذیر به یک چشم‌انداز پیوسته و هموار بدل می‌شود که اجازه می‌دهد بردار خروجی با گرادیان صعودی در مسیر افزایش کیفیت شناور شده و در گام نهایی با اعمال یک آستانه‌ی تصمیم به صفر و یک قطعی گرد شود.

علی‌رغم موفقیت تکنیک‌هایی نظیر آزادسازی پیوسته در تقریب فضای خروجی، جستجوی مبتنی بر گرادیان همچنان یک روش تقریبی است که خطر گرفتار شدن در بهینه‌های محلی را به‌همراه دارد. گرابر و همکاران (۲۰۱۸)~\cite{graber2018deep} تلاش کردند تا به‌جای تقریب، یک راهکار برای استنتاج ریاضی دقیق در شبکه‌های عمیق ارائه دهند. آن‌ها متوجه شدند که استنتاج دقیق تنها در مدل‌های ساده و خطی (که ساختار گرافیِ ثابت دارند) امکان‌پذیر است. بنابراین، یک معماری ترکیبی پیشنهاد دادند؛ یک بخش از مدل، وابستگی‌های ساده و محلی را با ساختارهای قابل‌حل مدل می‌کند و بخش دیگر، یک شبکه‌ی عصبی عمیق است که وابستگی‌های پیچیده و غیرخطی را در نظر می‌گیرد. برای اینکه این دو بخش با هم به توافق برسند، از ترفند آزادسازی لاگرانژی استفاده کردند که بهینه‌سازی را به‌صورت تناوبی میان بخش ساده و بخش پیچیده تقسیم می‌کند. درک شهودی آزادسازی لاگرانژی در این بافتار، تقسیم کار هوشمندانه است: مدل مسئله را به دو بخش ساختار محلی با حل دقیق و وابستگی‌های عمیق سراسری تجزیه می‌کند و متغیرهای دوگان (ضرایب لاگرانژ) به‌عنوان نیروهای جریمه‌کننده، این دو بخش را وادار می‌سازند تا از طریق گفت‌وگوی تکرارشونده بر سر یک پیکربندی واحد به اجماع برسند.

اهمیت بررسی این پژوهش در روشن ساختن یک محدودیت بنیادین است. راه‌حل‌هایی نظیر تجزیه‌ی لاگرانژی یا جستجوی گرادیانی (در \lr{SPEN} و \lr{DVN}) وابستگی مطلقی به این شرط دارند که طول و ابعاد فضای خروجی از پیش مشخص باشد. در متون باز و خودهمبسته، خروجی یک گراف ثابت نیست؛ بلکه دنباله‌ای با طول متغیر و انشعاباتی نامحدود به وسعت واژگان یک زبان است. اعمال گام‌های گرادیانی بر روی کلمات یک جمله‌ی در حال تولید، یا تلاش برای آزادسازی ریاضی یک درخت جستجوی بی‌نهایت، از نظر محاسباتی کاملاً بدتعریف و ناممکن است. این بن‌بست تئوریک اثبات می‌کند که برای استفاده از شبکه‌های انرژی در دامنه‌ی زبان طبیعی، باید پارادایم جستجو و استنتاج در فضای خروجی را به‌طور کامل کنار گذاشت و انرژی را در قالب یک ناظر بیرونی به‌کار گرفت.

با این حال، با وجود این ناکامی ساختاری در رویارویی با متون باز، راهکارهای ارائه‌شده برای تغییر معماری و تسریع استنتاج در این مدل‌ها، حاوی ایده‌ها و الگوهای ارزشمندی هستند که می‌توانند در طراحی رویکردهای جایگزین و تطبیق‌یافته در آینده راهگشا باشند.


\paragraph{استنتاج مستهلک‌شده و پویایی آموزش مشترک.}

پس از اثبات کندی و دشواری بهینه‌سازی تکرارشونده در زمان استنتاج، تو و گیمپل (۲۰۱۸)~\cite{tu2018learning} رویکرد استنتاج مستهلک‌شده\enfootnote{amortized inference} را برای شبکه‌های انرژی معرفی کردند. متدولوژی آن‌ها جایگزین کردن فرایند پرهزینه‌ی جستجو با یک شبکه‌ی استنتاج پیش‌خور است. از نظر ریاضی، اگر $E(x, y; \theta)$ تابع انرژی باشد، یک شبکه‌ی ثانویه با پارامترهای $\phi$ آموزش می‌بیند تا خروجی بهینه را مستقیماً تقریب بزند ($\hat{y} = I(x; \phi) \approx y^* = \arg\min_y E(x, y; \theta)$). از منظر شهودی، استهلاک در اینجا یعنی بار سنگین محاسباتِ بهینه‌سازی تکرارشونده در زمان آزمون، یک‌بار برای همیشه در وزن‌های شبکه‌ی استنتاج ذخیره و مستهلک می‌شود. این شبکه همانند یک متخصص باتجربه، با یک نگاه پیش‌خور به ورودی $x$، پیکربندی بهینه را مستقیماً در زمان ثابت $\mathcal{O}(1)$ حدس می‌زند، بدون آنکه نیازی به غلتیدن گام‌به‌گام در شیب‌های انرژی باشد. برای آموزش مشترک این معماری، از یک تابع زیان ساختاریافته‌ی حاشیه‌دار استفاده می‌شود تا تضمین گردد انرژی هر فرضیه‌ی نادرست یا خروجی خطادار ($\hat{y}$)، حداقل به اندازه‌ی یک حاشیه‌ی امنِ متناسب با خطای ساختاریافته ($\Delta(\hat{y}, y^*)$)، بالاتر از انرژیِ حقیقت پایه باشد ($E(x, \hat{y}; \theta) - E(x, y^*; \theta) \ge \Delta(\hat{y}, y^*)$). این تقطیر دانش، زمان استنتاج را از یک جستجوی طولانی به یک عبور پیش‌رو با زمان ثابت $\mathcal{O}(1)$ کاهش می‌دهد.

با این وجود، آموزش هم‌زمان و تناوبیِ این دو شبکه با ناپایداری‌های شدیدی همراه است. تو و همکاران (۲۰۲۰)~\cite{tu2020improving} با آسیب‌شناسی این فرایند نشان دادند که شبکه‌ی انرژی در طول آموزش غالباً حفره‌های کاذب انرژی برای خروجی‌های نادرست ایجاد می‌کند. در این شرایط، شبکه‌ی استنتاج به‌جای یادگیری ساختار صحیح، فریب خورده و پیش‌بینی‌های خود را به سمت این نواحی سوق می‌دهد. آن‌ها برای مهار این ناپایداری، یک تابع زیان ترکیبی طراحی کردند که علاوه بر شبکه‌ی استنتاج استاندارد، یک شبکه‌ی استنتاج افزوده-با-هزینه\enfootnote{cost-augmented} را نیز آموزش می‌دهد؛ این شبکه‌ی دوم به‌گونه‌ای طراحی شده است تا خروجی‌هایی تولید کند که دارای بیشترین خطا نسبت به خروجی صحیح هستند اما مدل انرژی به اشتباه به آن‌ها امتیاز خوبی می‌دهد. شهود این سازوکار همانند آموزش یک حریف آزمایشیِ سخت‌گیر است؛ شبکه‌ی افزوده-با-هزینه گسل‌ها و دره‌های کاذبی را که مدل انرژی به غلط امن پنداشته است شناسایی و برجسته می‌کند تا مدل انرژی با پر کردن و تنبیه این خروجی‌ها، ساختار چشم‌انداز بهینه‌سازی را تثبیت نماید.

از سوی دیگر، آموزش تمامی این شبکه‌های انرژی تا پیش از این، نیازمند دسترسی به مجموعه‌داده‌هایی با خروجی‌های ساختاریافته‌ی کاملاً صحیح (حقیقت پایه) بود. برای رفع این وابستگی مطلق، روشناس و همکاران (۲۰۱۹)~\cite{rooshenas2019searchguided} در رویکردی مبتنی بر نظارت ضعیف، نشان دادند که چگونه می‌توان از توابع پاداش خارجی (مانند یک سیستم امتیازدهنده‌ی ساده) به‌جای خروجی‌های طلایی استفاده کرد. در این متدولوژی، الگوریتم از جستجوی تصادفی بریده‌شده\enfootnote{truncated randomized search} برای تولید خروجی‌های متنوع استفاده کرده، آن‌ها را با تابع پاداش می‌سنجد و خروجی‌هایی را که بالاترین پاداش را می‌گیرند به‌عنوان شبه‌حقیقت پایه برای آموزش شبکه‌ی انرژی در نظر می‌گیرد.

در نهایت، از آنجا که شبکه‌ی استنتاج پیش‌خور تنها یک تقریب‌زننده‌ی سریع است و لزوماً خروجیِ دارای کمترین انرژی مطلق را پیدا نمی‌کند، پان و همکاران (۲۰۲۰)~\cite{pan2020adversarial} یک مکانیزم متخاصم با الهام از شبکه‌های مولد متخاصم معرفی کردند. در این ساختار، خروجی خام شبکه‌ی استنتاج صرفاً به‌عنوان یک حدس اولیه در نظر گرفته می‌شود و سپس با اعمال چند گام محدود از الگوریتم گرادیان کاهشی، این خروجی مستقیماً در چشم‌انداز انرژی مدل بهینه‌سازی و دقیق‌تر می‌گردد. اگرچه این متدولوژی‌ها با ارائه‌ی ایده‌هایی نظیر آموزش شبکه‌های دوگانه، بهره‌گیری از پاداش‌های خارجی و اصلاح خروجی با گرادیان، معماری مدل‌های انرژی را برای وظایف ساختاریافته به‌شکل چشمگیری بهبود بخشیدند، اما هیچ‌یک نتوانستند بر مانع بنیادین تولید زبان طبیعی غلبه کنند و این چالش در این پارادایم همچنان پابرجاست.


\paragraph{سایر کاربردها و تفاسیر.}

کاربرد مفاهیم شبکه‌های انرژی و استنتاج مستهلک‌شده تنها به فاز پیش‌بینی محدود نماند و به‌عنوان ابزاری برای تفسیرپذیری مدل‌های پیچیده نیز مورد توجه قرار گرفت. همان‌طور که در بخش~\ref{sec:3-1-3-3} بررسی شد، سیدصالحی و همکاران (۲۰۲۴)~\cite{seyyedsalehi2024sointer} چارچوب \lr{SOInter} را برای تفسیر مدل‌های دارای خروجی ساختاریافته معرفی کردند. این متدولوژی با الهام از معماری‌های دوگانه‌ی انرژی، یک شبکه‌ی انرژی سراسری را آموزش می‌دهد تا وابستگی‌های متقابل در خروجی مدل جعبه‌سیاه را درک کند. سپس، از این شبکه‌ی انرژی به‌عنوان یک ناظر برای هدایت یک شبکه‌ی مفسر (که عملکردی مشابه با شبکه‌ی استنتاج دارد) بهره می‌برد تا مهم‌ترین ویژگی‌های ورودی که کمترین انرژی را در ساختار خروجی ایجاد می‌کنند، انتخاب شوند. 

با وجود قدرت نظری این رویکرد، اعمال مستقیم آن برای ممیزی با دو محدودیت روبه‌روست. نخست، ارزیابی در سطح ویژگی‌های خرد (نظیر کلمات) و اعمال ماسک‌های گسسته بر روی یک پرامپت متنی، انسجام نحوی را تخریب کرده و به بروز خطای شوک توزیعی و تولید داده‌های \lr{OOD} منجر می‌گردد. اگرچه در متدولوژی \lr{SOInter} برای مهار نسبی این آسیب، شبکه‌ی انرژی در زمان آموزش مفسر با خروجی این ورودی‌های مخدوش به‌طور هم‌زمان تنظیم ظریف می‌شود تا با توزیع جدید سازگاری پیدا کند، اما محدودیت اصلی و دوم همچنان پابرجاست؛ این روش نیز مانند سایر معماری‌های مذکور در این بخش، ماهیتاً برای فضاهای خروجی با ابعاد ثابت و ساختار واضح طراحی شده است و پردازش توالی‌های نامحدود در آن بدتعریف است.

در نقطه‌ی مقابل رویکردهای مبتنی بر انرژی، رویکردی شکل گرفت تا بررسی شود آیا معماری‌های خودهمبسته‌ی استاندارد ذاتاً ظرفیت مدل‌سازی ساختارهای پیچیده‌ی درون‌متنی را دارند یا خیر. لیو و همکاران (۲۰۲۲)~\cite{liu2022autoregressive} با معرفی پیش‌بینی ساختاریافته‌ی خودهمبسته\enfootnote{autoregressive structured prediction} (\lr{ASP}) قابلیت‌مدل‌های زبانی پیش‌آموزش‌دیده را در وظایف پیش‌بینی ساختاریافته (نظیر استخراج رابطه یا ارجاع متقابل) ارزیابی کردند. بررسی‌های آن‌ها نشان داد که اگر این ساختارهای درهم‌تنیده صرفاً به یک رشته‌ی متنی مسطح تبدیل شوند، شبکه‌ی خودهمبسته در ردیابی وابستگی‌های چندجانبه دچار افت عملکرد می‌شود. متدولوژی آن‌ها برای فعال‌سازی این ظرفیت ساختاری در مدل، تغییر فرمت خروجی بود. شبکه به‌جای تولید متن مسطح، آموزش می‌بیند تا یک توالی از کنش‌های ساختارساز تولید کند. در چارچوب \lr{ASP}، مدل زبانی کنش‌هایی نظیر تعیین کرانه‌های یک قطعه‌ی متنی (با باز و بسته کردن نمادین براکت‌ها)، ارجاع به قطعات پردازش‌شده‌ی پیشین و تخصیص برچسب را صادر می‌کند تا ساختار گراف وابستگی‌ها به‌صورت گام‌به‌گام و دقیق بنا شود.

موفقیت این رویکرد حامل یک بینش تئوریک مهم است. این متدولوژی اثبات می‌کند که مکانیزم پردازشی در معماری مدل‌های زبانی، ظرفیت ذاتی و قدرتمندی برای استخراج ساختار و ردیابی وابستگی‌های عمیق از درون متن دارد. با وجود اینکه متوقف کردن تولید زبان طبیعی و تبدیل خروجی به یک توالی از نشانه‌های کنترلی و ساختاری برای ارزیابی متون مدل‌های جعبه‌سیاه کاربردی نیست، اما این توانایی پنهان در کشف ساختار بسیار الهام‌بخش است. این ویژگی نشان می‌دهد که معماری پایه‌ی این شبکه‌ها محدود به تولید حریصانه‌ی توکن‌ها نیست و توانایی ارزیابی قیود ساختاری را داراست؛ ظرفیتی که بستر را برای اعمال توابع ارزیاب کل‌نگر در متن آزاد فراهم می‌سازد.


\subsection{مدل‌های انرژی در پردازش زبان طبیعی}\label{sec:3-2-3}

پس از اثبات پشتوانه نظری شبکه‌های انرژی (بخش~\ref{sec:2-4}) و مشاهده‌ی کارایی عملی آن‌ها در مسائل پیش‌بینی ساختاریافته (بخش~\ref{sec:3-2-2})، این بخش به چگونگی تطبیق و ورود این معماری‌ها به حوزه‌ی پردازش زبان طبیعی می‌پردازد. چالش اصلی در این گذار، تفاوت ماهوی در ساختار تولید است؛ مدل‌های زبانی به‌طور سنتی متن را توکن‌به‌توکن تولید می‌کنند، در حالی که مدل‌های انرژی ماهیتاً ارزیابانِ کلان‌نگری هستند که کل ساختار خروجی را یکپارچه می‌سنجند. برای درک چگونگی پیوند انرژی با متن، ابتدا باید تلاش‌های اولیه‌ی محققان برای تغییر پارادایم تولید از حالت متوالی به حالت موازی بررسی شود. این تلاش‌های پیشگامانه اگرچه بستر لازم برای پردازش یکپارچه را فراهم کردند، اما با بحران ازهم‌گسیختگی معنایی رویه‌رو بودند که در نهایت، ضرورت وجود یک ارزیاب سراسری مبتنی بر انرژی را توجیه نمود.


\paragraph{بحران انسجام در تولید موازی و ضرورت ارزیابی سراسری.}

بستر تاریخی پیش از ورود مدل‌های انرژی به زبان طبیعی، با تلاش برای رفع محدودیت‌های تولید خودهمبسته شکل گرفت. مدل‌های زبانی استاندارد، متن را به‌صورت چپ‌به‌راست تولید می‌کنند ($P(Y \mid X) = \prod_{t=1}^T P(y_t \mid y_{<t}, X)$). اگرچه این روش انسجام گرامری را تضمین می‌کند، اما به‌دلیل ماهیت ترتیبی، فرایندی بسیار کند است. گو و همکاران (۲۰۱۸)~\cite{gu2018nonautoregressive} برای تسریع این روند، معماری ترجمه‌ی ماشینی غیرخودهمبسته\enfootnote{non-autoregressive translation} (\lr{NAT}) را معرفی کردند. در این ساختار، ترنسفورمر تمام توکن‌های خروجی را به‌طور موازی و مستقل از یکدیگر پیش‌بینی می‌کند ($P(Y \mid X) = \prod_{t=1}^T P(y_t \mid X)$). از منظر شهودی، در رابطه‌ی خودهمبسته، هر کلمه مشروط به تاریخچه‌ی تمام کلمات قبلی ($y_{<t}$) صادر می‌شود؛ این زنجیره‌ی متوالی مانع از تداخل نحوی شده اما فرایند تولید را به صف تک‌ستونی و کند مبدل می‌سازد. در مقابل، رابطه‌ی غیرخودهمبسته این اتصال شرطی را کاملاً قطع کرده و فرض می‌کند توکن‌ها مستقل از یکدیگر و صرفاً با اتکا به ورودی $X$ تولید می‌شوند. این استقلال شرطی اگرچه تولید موازی در زمان ثابت را ممکن می‌سازد، اما ارتباط توکن‌ها با یکدیگر را از بین برده و مدل را در مواجهه با چندین الگوی بیانی معتبر، دچار سردرگمی و التقاط می‌سازد. برای مدیریت مشکل نامشخص بودن طول توالی در زمان تولید یکباره، آن‌ها از متغیری پنهان به‌نام باروری استفاده کردند؛ این ماژول پیش از تولید نهایی، تخمین می‌زند که هر کلمه‌ی ورودی باید به‌عنوان الگو به چند توکن در خروجی بسط یا کپی شود تا یک اسکلت اولیه با طول مشخص برای توالیِ هدف شکل بگیرد.

اگرچه این معماری سرعت استنتاج را به‌شدت افزایش داد، اما یک آسیب‌پذیری واقعی و بنیادین به‌نام بحران استقلال شرطی را آشکار ساخت. از آنجا که توکن‌ها در زمان تولید از حضور یکدیگر بی‌خبرند، مدل توانایی حفظ ساختار سراسری را از دست می‌دهد. این امر منجر به بروز خطای چندنمایی می‌گردد؛ در زبان طبیعی یک جمله می‌تواند چندین بیان کاملاً متفاوت اما معتبر داشته باشد. از آنجا که مدل \lr{NAT} نمی‌تواند با توکن‌های قبلی هماهنگ شود تا صرفاً یکی از این مسیرهای معتبر را تا انتها دنبال کند، کلماتی از الگوهای بیانیِ مختلف را با هم ترکیب می‌کند که در عمل به خروجی‌های نامفهوم، تکرار متوالی کلمات یا حذف بخش‌های مهمی از متن منجر می‌شود.

در پاسخ به این فروپاشی انسجام، غزوینی‌نژاد و همکاران (۲۰۱۹)~\cite{ghazvininejad2019maskpredict} پارادایم پالایش تکرارشونده\enfootnote{iterative refinement} را پیشنهاد دادند. آن‌ها برای پیاده‌سازی این رویکرد، معماری مدل‌های زبانی ماسک‌شده‌ی شرطی\enfootnote{conditional masked language models} (\lr{CMLM}) را طراحی کرده و الگوریتم رمزگشاییِ \lr{Mask-Predict} را بر روی آن اجرا نمودند. در این الگوریتم، مدل ابتدا یک خروجی خام به‌طور هم‌زمان تولید می‌کند. سپس در چندین گام، توکن‌هایی که کمترین احتمال (پایین‌ترین سطح اطمینان) را دارند ماسک کرده و آن‌ها را مجدداً با توجه به بافتار توکن‌های مطمئن‌تر پیش‌بینی می‌کند.

در تکامل همین معماری، ساوینوف و همکاران (۲۰۲۲)~\cite{savinov2022stepunrolled} چارچوب \lr{SUNDAE} را توسعه دادند. مسئله‌ی محوری آن‌ها پدیده‌ی خطای انباشته ناشی از تفاوت فاز آموزش و استنتاج بود؛ در \lr{CMLM}ِ استاندارد، مدل روی ماسک‌های کاملاً تصادفی در متون بی‌نقص انسانی آموزش می‌دید، اما در زمان تست مجبور بود اشتباهات ساختاری خودش را ترمیم کند. روش \lr{SUNDAE} این مشکل را با باز کردن گام‌های نویزگیری در زمان آموزش حل کرد؛ به‌گونه‌ای که مدل یاد گرفت خروجی ناقص خودش را دریافت کرده و آن را در گام‌های متوالی تصحیح کند.

با وجود اینکه رویکردهای مبتنی بر پالایش، کیفیت متون غیرخودهمبسته را تا حد مناسبی بهبود بخشیدند، اما همچنان برای ارزیابی درستی متن، منحصراً به احتمالات محلی هر توکن اتکا می‌کردند. عمق این بن‌بست ساختاری (یعنی تلاش برای تولید موازیِ بی‌نقص بدون داشتن یک ارزیاب کلان) در پژوهش گو و کونگ (۲۰۲۱)~\cite{gu2021fully} به‌صورت تجربی اثبات گردید. آن‌ها با انجام تحلیل‌های گسترده نشان دادند که رسیدن به کیفیت مدل‌های خودهمبسته در یک مدلِ کاملاً موازی (بدون استفاده از پالایش تکرارشونده در زمان تست)، نیازمند اعمال ترکیبی سنگین از ترفندهای مهندسی است. این ترفندها شامل استفاده از رمزگذارهای بسیار عمیق و به‌ویژه تقطیر دانش توالی از یک مدل معلم بود تا پیچیدگی و چندنمایی داده‌های آموزشی به‌صورت مصنوعی کاهش یافته و مدل بتواند تولید موازی را تحمل کند.

این هزینه‌های سنگین مهندسی و نیاز به دستکاری توزیع داده‌ها، یک خلأ تئوریک را در پردازش زبان طبیعی برجسته ساخت. پژوهش‌ها نشان دادند که ارزیابی احتمال توکن‌ها به‌صورت مجزا، برای داوریِ کیفیت یک توالی متنی کافی نیست. زبان طبیعی برای حفظ ساختار درهم‌تنیده‌ی خود، فراتر از تولید موازی، نیازمند یک مکانیزم ارزیاب سراسری است که بتواند کیفیت کل توالی را در قالب یک امتیاز داوری کند.


\paragraph{استقرار شبکه‌های انرژی به‌عنوان ارزیابان ساختار متنی.}

تلاش برای خروج از بن‌بست ارزیابی‌های توکن‌محور، منجر به ظهور پارادایم مدل‌سازی بر پایه‌ی انرژی گردید که با نگاهی کل‌نگر، توالی را به‌عنوان یک واحد یکپارچه مدل‌سازی می‌کنند. با این حال، همان‌طور که پیش‌تر اشاره شد، بزرگترین مانع در اعمال مستقیم \lr{EBM}ها بر روی زبان طبیعی، ماهیت گسسته و ترکیبیاتی متن است. در مدل‌سازی انرژی، محاسبه‌ی دقیق احتمالات نیازمند تخمین تابع افراز ($Z_\theta = \sum_x \exp(-\mathcal{E}_\theta(x))$) است. دلیل شهودی این حل‌ناپذیری آن است که برای تبدیل انرژی به یک توزیع احتمال استاندارد، باید مجموع نماییِ انرژی تمامی جملات ممکن در مخرج کسر ($Z_\theta$) محاسبه شود؛ در فضایی با اندازه‌ی واژگان $|\mathcal{V}|$ و طول جمله $T$، این جمع روی $|\mathcal{V}|^T$ حالت ترکیبیاتیِ نجومی (مثلاً $50000^{40}$) گسترده شده است و برخلاف متغیرهای پیوسته، گرادیان انرژی نیز برای هدایت نمونه‌برداری گسسته در میان کلمات کارایی ندارد.

بختیین و همکاران (۲۰۱۹)~\cite{bakhtin2019realfake} با ارائه‌ی یک چارچوب متمایزگر برای تشخیص متن تولیدشده توسط ماشین از متن انسانی، راهکاری عملی برای دور زدن این بن‌بست ارائه دادند. آن‌ها به‌جای درگیری با محاسبه‌ی مستقیم تابع افراز، از تکنیک آموزش تقابلی بهره بردند. در این معماری، تابع انرژی صرفاً یک لایه‌ی خطی ساده نیست، بلکه پایه‌ی آن یک شبکه‌ی عصبی عمیق نظیر ترنسفورمر دوسویه است. این مدل، به‌عنوان متمایزگر عمل کرده، کل توالی را به‌صورت یکپارچه دریافت می‌کند و مستقیماً یک اسکالر را به‌عنوان انرژی ($\mathcal{E}(x)$) تولید می‌نماید. برای آموزش این شبکه، متون انسانی به‌عنوان نمونه‌های مثبت ($x^+$) و متون تولیدشده توسط یک مدل زبانی خودهمبسته‌ی پیش‌آموزش‌دیده به‌عنوان نمونه‌های منفی ($x^-$) در نظر گرفته می‌شوند. تابع هدف، کمینه‌سازی زیان آنتروپی متقاطع دودویی (\lr{BCE}) بر مقادیر انرژی است:
\begin{equation}\label{eq:chap3_nlp_ebm_bce}
    \mathcal{L}_{\text{\lr{BCE}}} = -\log(\sigma(-\mathcal{E}(x^+))) - \log(1 - \sigma(-\mathcal{E}(\hat{x}^-))) = -\log(\sigma(-\mathcal{E}(x^+))) - \log(\sigma(\mathcal{E}(\hat{x}^-)))
\end{equation}
که در آن $\sigma(z) = \frac{1}{1 + \exp(-z)}$ تابع سیگموئید استاندارد است و $\hat{x}^-$ نمایانگرِ سخت‌ترین نمونه‌ی منفی است؛ یعنی توالی منتخبی که در میان مجموعه‌ای از $k$ نمونه‌ی مختلف تولیدشده توسط مولد پایه، به اشتباه کمترین انرژی را از دید متمایزگر دریافت کرده است ($\hat{x}^- = \arg\min_{i=1}^k \mathcal{E}(x^-_i)$). از منظر شهودی، عبارت $\sigma(-\mathcal{E}(x))$ بازتاب‌دهنده‌ی احتمال تعلق متن $x$ به دسته‌ی متون واقعیِ انسانی است؛ زیرا انرژی کمتر (منفی‌تر)، این احتمال را به یک نزدیک می‌سازد. بر این مبنا:
نخست، برای نمونه‌های اصیل انسانی ($x^+$)، کمینه‌سازی جمله‌ی اول مستلزم آن است که مدل به آن‌ها انرژی پایینی اختصاص دهد ($\mathcal{E}(x^+) \to -\infty$) تا $\sigma(-\mathcal{E}(x^+)) \to 1$ شده و جریمه‌ی $-\log(1) = 0$ حاصل گردد.
دوم، برای سخت‌ترین نمونه‌های منفی ماشینی ($\hat{x}^-$)، جمله‌ی دوم مدل را ملزم می‌کند که انرژی آن‌ها را به‌شدت بالا ببرد ($\mathcal{E}(\hat{x}^-) \to +\infty$) تا احتمال واقعی بودن آن‌ها به صفر میل کرده و متقابلاً $\sigma(\mathcal{E}(\hat{x}^-)) \to 1$ شود.
انتخاب سخت‌ترین نمونه‌ی منفی تضمین می‌کند که متمایزگر فعالانه مناطقی از فضای متن را هدف بگیرد که مدل مولد به غلط در آن‌ها خطاهای ساختاری واقع‌نما تولید کرده است. این ساختار، مدل را قادر می‌سازد تا بدون نیاز به نمونه‌برداری‌های پرهزینه و پیچیده‌ی درونی، تفاوت‌های ساختاری و کلان میان یک متن ارگانیک و یک متن ماشینی را در قالب یک عدد یکپارچه یاد بگیرد.

موفقیت این ترنسفورمرهای دوسویه در نقش متمایزگر، پایه‌گذار بسط نظری آن‌ها برای بهبود مستقیم فرایند تولید متن گردید. دنگ و همکاران (۲۰۲۰) و بختیین و همکاران (۲۰۲۱)~\cite{deng2020residual, bakhtin2021residual} چارچوب «مدل‌های انرژی پسماند»\enfootnote{residual energy-based models} (\lr{REBM}) را توسعه دادند. هدف این ساختار، جبران نارسایی‌های مدل‌های خودهمبسته (نظیر سوگیری مواجهه\enfootnote{exposure bias} و تکرارهای ساختاری) از طریق اعمال یک ارزیاب کلان بود. آن‌ها توزیع احتمالِ توأم ($P_\theta(x)$) را به‌صورت ضرب توزیع یک مولد خودهمبسته‌ی پایه‌ای ($P_{\text{\lr{LM}}}(x)$) و یک عبارت انرژی سراسری یادگیرنده ($\mathcal{E}_\theta(x)$) فرمول‌بندی کردند:
\begin{equation}\label{eq:chap3_residual_ebm_p}
    P_\theta(x) = \frac{P_{\text{\lr{LM}}}(x) \exp(-\mathcal{E}_\theta(x))}{Z_\theta}
\end{equation}
شهود نهفته در عبارت مدل‌های پسماند آن است که یادگیری کل ساختار پیچیده، گرامر، و واژگان زبان از نقطه‌ی صفر صرفاً به کمک یک مدل انرژی، فوق‌العاده سنگین و ناپایدار است. در فرمول فوق، مدل زبانی پایه ($P_{\text{\lr{LM}}}$) بخش عمده‌ی گرامر و ساختار محلی را مدل می‌کند، در حالی که فاکتور انرژی ($\exp(-\mathcal{E}_\theta(x))$) تنها در فضای پسماند این مدل به‌عنوان یک فیلتر اصلاحی سراسری عمل کرده و توالی‌هایی را که با وجود صحت محلی کلمات فاقد انسجام کلان هستند جریمه می‌کند. برای آموزش این عبارت انرژی، نویسندگان از تکنیک تخمین تقابلی نویز شرطی\enfootnote{conditional noise contrastive estimation} (\lr{CNCE}) بهره بردند. تابع هدف این روش، بسطی بر تابع \lr{BCE} چارچوب متمایزگر پیشین در معادله‌ی~\ref{eq:chap3_nlp_ebm_bce} است:
\begin{equation}\label{eq:chap3_residual_ebm_cnce}
    \max_\theta \mathbb{E}_{x_+ \sim P_{\text{\lr{data}}}} \left[ \log \frac{1}{1 + \exp(\mathcal{E}_\theta(x_+))} \right] + \mathbb{E}_{x_- \sim P_{\text{\lr{LM}}}} \left[ \log \frac{1}{1 + \exp(-\mathcal{E}_\theta(x_-))} \right]
\end{equation}
درک شهودی این تابع هدفِ بیشینه‌سازی بر پایه‌ی برآورد تقابلی نویز است که در آن مدل زبانی پایه نقش توزیع نویز را ایفا می‌کند:
عبارت نخست با کاهش دادن مقدار انرژی برای نمونه‌های واقعی انسانی ($\mathcal{E}_\theta(x_+) \to -\infty$)، توان مثبت نمایی $\exp(\mathcal{E}_\theta(x_+))$ را به صفر میل داده، مخرج کسر را به یک نزدیک کرده و مقدار لگاریتم را به حداکثر ممکن (صفر) می‌رساند.
به‌طور متقابل، عبارت دوم با تخصیص انرژی بالا به نمونه‌های نویز یا همان متون تولیدشده توسط مولد پایه ($\mathcal{E}_\theta(x_-) \to +\infty$)، توان منفی نمایی $\exp(-\mathcal{E}_\theta(x_-))$ را ناپدید کرده و مقدار این عبارت را نیز به حداکثر می‌رساند. این تقارن، مدل انرژی را وادار می‌سازد که مرزهای کیفی میان متن اصیل و ماشینی را به دقت ترسیم کند.

نوآوری اصلی چارچوب \lr{REBM} در فاز استنتاج است که در آن از نمونه‌برداری اهمیتی خودنرمال‌ساز\enfootnote{self-normalizing importance sampling} استفاده می‌شود. در این فاز، هزاران توالی کاندیدا توسط مولد پایه تولید شده و سپس با وزن‌های اهمیتیِ متناسب با $\exp(-\mathcal{E}_\theta(x^i))$ توسط مدل انرژی بازنمونه‌برداری و رتبه‌بندی مجدد می‌گردند تا گزینه‌های دارای انرژی پایین برگزیده شوند. اگرچه این رویکرد از نظر تجربی کیفیت و انسجام متون را به‌شدت ارتقا داد، اما اجرای متوالی دو مدل سنگین و ارزیابی تک‌به‌تکِ انبوهی از کاندیداها، بار محاسباتی فاجعه‌باری ایجاد کرد که استقرار این سیستم‌ها را در محیط‌های بلادرنگ عملاً محدود می‌ساخت.

برای غلبه بر این گلوگاه محاسباتی، تو و همکاران (۲۰۲۰)~\cite{tu2020engine} چارچوب \lr{ENGINE}\enfootnote{energy-based inference networks} را با هدف ترکیب دقت ارزیاب‌های انرژی با سرعت مدل‌های غیرخودهمبسته (\lr{NAT}) معرفی کردند. این پژوهش، مسئله را در قالب پارادایم استنتاج مستهلک‌شده (که پیش‌تر در بخش~\ref{sec:3-2-2} معرفی شد) بازتعریف کرد؛ به این معنا که به‌جای اجرای جستجوی ترکیبیاتی و پرهزینه برای یافتن توالیِ دارای کمترین انرژی در زمان تست، یک شبکه‌ی ثانویه برای انجام مستقیم و سریع این پیش‌بینی آموزش می‌بیند.

در چارچوب \lr{ENGINE}، یک مدل زبانی خودهمبسته (\lr{AR}) پیش‌آموزش‌دیده، نقش تابع انرژی ($\mathcal{E}_\Theta$) را ایفا می‌کند؛ به این معنا که انرژی یک توالی، معادل با منفیِ لگاریتم احتمال آن در نگاه این مدل معلم تعریف می‌شود ($\mathcal{E} = -\log P_\Theta$)؛ لذا هرچه جمله از دید معلم محتمل‌تر و طبیعی‌تر باشد، انرژی آن کمتر است. هدف این است که یک مدل \lr{NAT} به‌عنوان شبکه‌ی استنتاج ($A_\Psi$) آموزش ببیند تا مستقیماً پارامترهای خود را به‌گونه‌ای تنظیم کند که توالی‌های تولیدی آن، کمترین انرژی ممکن را نزد مدل معلم به‌دست آورند:
\begin{equation}\label{eq:chap3_engine_inference}
    \hat{\Psi} = \arg\min_\Psi \sum_{\langle x,y \rangle \in \mathcal{D}} \mathcal{E}_\Theta(x, A_\Psi(x))
\end{equation}
بزرگترین چالش فنی در این مسیر، سد مشتق‌ناپذیری در فضای گسسته‌ی متن است که مانع از جریان یافتنِ گرادیان از تابع انرژی (معلم) به سمت شبکه‌ی استنتاج (\lr{NAT}) می‌شود. نویسندگان برای حل این معضل، انرژی کل را به مجموعه‌ای از انرژی‌های محلی ($e_t$) برای هر جایگاه توکن تجزیه کردند:
\begin{equation}\label{eq:chap3_engine_energy}
    \mathcal{E}_\Theta(x,y) = \sum_{t=1}^{|y|} e_t(x,y)
\end{equation}
که در آن هر مؤلفه‌ی انرژی محلی به‌صورت زیر تعریف می‌شود:
\begin{equation}\label{eq:chap3_engine_local_energy}
    e_t(x, y) = - O_2{(z_t)}^\top \log P_\Theta(\cdot \mid O_1(z_0), O_1(z_1), \dots, O_1(z_{t-1}), x)
\end{equation}
در این فرمول‌بندی، $z_t$ بردار لوجیت‌های خروجیِ شبکه‌ی استنتاج در جایگاه $t$ است و دو عملگر واسط $O_1$ و $O_2$ یک درگاه کاملاً مشتق‌پذیر میان دو شبکه ایجاد می‌کنند:
عملگر $O_1$ خروجی‌های خام (لوجیت‌های) شبکه‌ی استنتاج را بدون گسسته‌سازی خشن، به فضایی پیوسته (مانند توزیع هموار \lr{Softmax} یا تقریب \lr{Straight-Through}) نگاشت می‌کند تا به‌عنوان بافتار تاریخچه‌ی ورودی توسط رمزگشای مدل معلم قابل‌پردازش باشد.
عملگر $O_2$ نیز وظیفه دارد میان توزیع احتمالی تولیدشده توسط مدل \lr{NAT} در گام $t$ و پیش‌بینی مدل معلم در آن گام، یک معیار هزینه یا انرژی محلی به فرم آنتروپی متقاطع محاسبه کند.
این ساختار یکپارچه اجازه می‌دهد گرادیان خطای مدل معلم مستقیماً به پارامترهای $\Psi$ شبکه‌ی استنتاج پس‌انتشار یابد؛ بدین‌ترتیب، \lr{ENGINE} توانست بدون نیاز به داده‌های تقطیرشده، سرعت بالای مدل‌های غیرخودهمبسته را با دانش عمیق مدل‌های خودهمبسته تلفیق کرده و به نتایجی در سطح هنر دست یابد.

با این حال، مقایسه‌ی این رویکرد با مدل‌های \lr{REBM} نشان‌دهنده یک مصالحه‌ی میان دقت ساختاری و سرعت عملیاتی است. اگرچه \lr{ENGINE} با استهلاک بار محاسباتی در یک شبکه‌ی استنتاج، تولید متن را به‌شدت تسریع می‌کند، اما پایداری آن به کیفیت تقریب‌های گرادیان در عملگرهای واسط وابسته است. علاوه بر این، تکیه‌ی مطلق تابع انرژی به یک مدل خودهمبسته ازپیش‌تعیین‌شده، هم یک مزیت و هم یک محدودیت بنیادین محسوب می‌شود؛ از یک سو فرایند آموزش را بسیار کارآمد می‌کند، اما از سوی دیگر، توانایی شبکه‌ی استنتاج را به سقف توانمندی‌های همان مدل معلم محدود ساخته و مانع از دستیابی به دیدگاه‌های ارزیاب کلان‌نگری می‌شود که فراتر از احتمالات گام‌به‌گامِ مدل‌های خودهمبسته حرکت می‌کنند. از طرف دیگر، این وابستگی ساختاری باعث می‌شود \lr{ENGINE} چارچوبی با انعطاف‌پذیری کمتر باشد که در مقایسه با چارچوب کلیِ \lr{REBM}، انطباق آن با شرایط و وظایف گوناگون دشوارتر است. لذا، هر یک از این پارادایم‌ها بسته به نیازمندی‌ها، جایگاه تحقیقاتی و کاربردی خاص خود را دارند.


\paragraph{پیشروی در فضاهای پیوسته و همسوسازی در مقیاس‌های کلان.}

دامنه‌ی نفوذ شبکه‌های انرژی از بهبود مدل‌های خودهمبسته فراتر رفته و به‌عنوان راهکاری برای داوری کلان، کنترل خروجی در فضاهای پیوسته، و همسوسازی مدل‌ها با ترجیحات انسانی نیز به‌کار گرفته شده است. 

یکی از مهم‌ترین گام‌ها در اثبات برتری مدل‌های انرژی برای ارزیابی متون، توسط باتاچاریا و همکاران (۲۰۲۱)~\cite{bhattacharyya2021reranking} برداشته شد. مسئله‌ی محوری این پژوهش، تضاد بنیادین میان آموزش مبتنی بر برآورد بیشینه‌ی درست‌نمایی (\lr{MLE}) در مولدهای خودهمبسته و معیارهای ارزیابی سراسری کیفیت ترجمه بود. مولدهای استاندارد توالی را به‌صورت محلی بهینه‌سازی می‌کنند و در جستجوی پرتو غالباً از انتخاب بهترین توالی یکپارچه باز می‌مانند. متدولوژیِ \lr{EBR}\enfootnote{energy-based reranking} برای حل این معضل، از یک مدل انرژیِ ترنسفورمری به‌عنوان ارزیاب پساعملی بهره می‌برد. در این چارچوب، مولد پایه لیستی از کاندیداهای برتر تولید کرده و سپس مدل انرژی (که با یک تابع زیان حاشیه‌ای آموزش دیده است) آن‌ها را بر اساس کیفیت سراسری‌شان امتیازدهی و رتبه‌بندی مجدد می‌کند. این پژوهش تثبیت می‌کند که مدل‌های انرژی ذاتاً ابزارهای بهتری برای داوری ساختار کلان و انسجام توالی هستند.

در حوزه‌ی تولید متن مقید، چین و همکاران (۲۰۲۲)~\cite{qin2022cold} با معرفی چارچوب \lr{COLD Decoding} نشان دادند که چگونه می‌توان محدودیت‌های معنایی پیچیده (مانند کنترل لحن یا گنجاندن کلمات کلیدی) را بدون تنظیم ظریف بر مدل‌های زبانی اعمال کرد. متدولوژی آن‌ها این محدودیت‌ها را در قالب مجموعه‌ای از توابع انرژی فرمول‌بندی می‌کند. از آنجا که فضای متن گسسته است، آن‌ها از تکنیک آزادسازی پیوسته روی لوجیت‌های خروجی استفاده می‌کنند تا توزیع کلمات به متغیرهای پیوسته بسط یابد. در زمان استنتاج، با استفاده از دینامیک لانژون، این بازنمایی‌های پیوسته در چندین گام متوالی با کمک گرادیان انرژی به‌روزرسانی می‌شوند تا محدودیت‌ها ارضا گردند. از منظر شهودی، دینامیک لانژون در این فرایند همانند حرکت یک ذره در یک میدان پتانسیل همراه با لرزش حرارتی عمل می‌کند: گرادیان منفی انرژی ($-\nabla_y E$) ذره را در جهت ارضای محدودیت‌ها به سمت قعر دره‌های انرژی سوق می‌دهد، در حالی که نویز تصادفیِ کنترل‌شده مانع از گیر افتادن در چاله‌های بهینه‌ی محلی می‌شود؛ در انتها، این موقعیت‌های پیوسته به نزدیک‌ترین توکن‌های گسسته در واژگان نگاشت می‌گردند. این رویکرد قدرت توابع انرژی را در هدایت مدل‌های زبانی نشان می‌دهد، اما به‌دلیل نیاز به محاسبه‌ی گرادیان در طول مراحل نمونه‌برداری، بار سنگینی در زمان تولید تحمیل می‌کند.

در سطح معماری‌های مولد غیرخودهمبسته، زو و همکاران (۲۰۲۵)~\cite{xu2025energybased} چارچوب «مدل‌های زبانی انتشاری مبتنی بر انرژی»\enfootnote{energy-based diffusion language models} (\lr{EDLM}) را برای رفع نقص مدل‌های انتشاری گسسته توسعه دادند. این مدل‌ها در زمان استنتاج، با کاهش تعداد گام‌های نمونه‌برداری، دچار افت شدید کیفیت می‌شوند؛ خطایی که ناشی از ناتوانی مدل در درک وابستگی‌های سراسری توکن‌ها در یک گام واحد است. متدولوژی \lr{EDLM} برای ترمیم این خطا، یک مدل انرژی پسماند را مستقیماً وارد فرایند انتقال معکوس می‌کند تا در هر گام، تقریب احتمالات تصحیح شود و سپس از نمونه‌برداری اهمیتیِ موازی برای استخراج خروجی استفاده می‌گردد. این معماری ثابت می‌کند که توابع انرژی حتی ظرفیت کالیبره کردن و ترمیم پویاییِ ساختاری پیشرفته‌ترین مدل‌های موازی را نیز دارند.

در نهایت، قابلیت‌های ارزیاب کلان در پژوهش لوچاب و ژانگ (۲۰۲۵)~\cite{lochab2025energybased} به حوزه‌ی همسوسازی با ترجیحات انسانی گسترش یافت. آن‌ها چارچوب «مدل پاداش مبتنی بر انرژی»\enfootnote{energy-based reward model} (\lr{EBRM}) را برای مهار پدیده‌ی تقلب در پاداش معرفی کردند. مدل‌های پاداش استاندارد، ترجیحات انسانی را به‌صورت یک اسکالر قطعی تقریب می‌زنند و توانایی پوشش عدم‌قطعیت‌های موجود در داده‌های انسانی را ندارند. روش \lr{EBRM} این ترجیحات را در قالب یک توزیع انرژی مدل‌سازی می‌کند ($P_\theta(y \mid x) \propto \exp(-E_\theta(x, y))$). به بیان شهودی، به‌جای فشرده‌سازی بازخورد چندوجهی و پرنوسان انسان به یک نمره‌ی اسکالر نقطه‌ای که مستعد فریب خوردن است، مدل پاداش یک چشم‌انداز احتمالی تشکیل می‌دهد که عدم‌قطعیت‌های داده‌های انسانی را منعکس می‌سازد. این مدل با استفاده از آموزش تقابلی آگاه به نویز برچسب\enfootnote{label-noise-aware contrastive training} آموزش می‌بیند و به‌عنوان یک پالایش‌گر پساعملی عمل می‌کند.

این متدولوژی‌ها در یک نگاه جامع اثبات می‌کنند که مدل‌های انرژی از تواناترین ابزارها برای داوری یکپارچه و کلان ویژگی‌های متنی (از مطلوبیت و ایمنی تا زبان‌مایگی و انسجام ترجمه) هستند. لذا انرژیِ محاسباتی می‌تواند به‌عنوان یک امتیازدهنده‌ی سبک، فشرده و کلی عمل کرده و ضعف مدل‌های زبانی سنتی را در درک ساختار و ارزیابی کیفیت سراسری به‌خوبی پوشش دهد.



% =========================================================
% بدنه سوم
% =========================================================



\section{جمع‌بندی و شناسایی شکاف تحقیقاتی}\label{sec:3-3}

مرور جامع پارادایم‌های تفسیرپذیری در این فصل، پرده از یک بن‌بست عملیاتی و نظری در ارزیابی سیستم‌های هوش مصنوعی مولد و مدرن برمی‌دارد. همان‌طور که در بخش~\ref{sec:3-1-1} واکاوی شد، رویکردهای تفسیرپذیری درونی به‌دلیل تحمیل محدودیت‌های ساختاری، ایجاد افت شدید در عملکرد و ناتوانی در همگام‌سازی با ماهیت توزیع‌شده‌ی زبان، برای معماری‌های پیچیده‌ی زبانی عملاً ناکارآمدند. از سوی دیگر، متدولوژی‌های قدرتمند پساعملی جعبه‌سفید (بخش~\ref{sec:3-1-2}) – اعم از روش‌های مبتنی بر سیگنال‌های پایه‌ی گرادیان، کاوشگرهای خطی در فضای پنهان، و رهیافت‌های تفسیرپذیری مکانیکی – وابستگی مطلقی به دسترسی کامل به معماری داخلی، ماتریس پارامترها و گراف محاسباتی دارند. اگرچه این متدولوژی‌ها در ارزیابی‌های تجربی مدل‌های منبع‌باز به‌عنوان خطوط مبنای قوی و معتبر عمل می‌کنند، اما اعمال مستقیم آن‌ها برای حسابرسیِ مدل‌های مرزی و تجاری که در انحصار رابط‌های برنامه‌نویسی (\lr{API}) قرار دارند و دسترسی به وزن‌ها یا گرادیان‌ها در آن‌ها مسدود است، ناممکن خواهد بود.

در مواجهه با این انسداد، پناه بردن به روش‌های پساعملی جعبه‌سیاه رایج (بخش~\ref{sec:3-1-3}) نیز با نارسایی‌های عمیقی همراه است. رویکردهای مبتنی بر خودتوضیحی و استدلال درون‌بافتاری، لزوماً بازتاب‌دهنده‌ی سازوکار واقعی اثرگذاری ورودی بر خروجی نبوده و مدل را در معرض خطر توهمات تبیینی و توجیهات پساعملی قرار می‌دهند. همچنین، رویکردهای متقدم اختلالی و مستقل از مدل (نظیر \lr{LIME}، \lr{SHAP} و نسخه‌های متنی آن‌ها) که بر ارزیابی متوالی توکن‌ها استوارند، در بستر زبان طبیعی در چهار تله‌ی مهلک گرفتار می‌شوند: نخست، تخریب گرامر و ساختار نحوی که به تولید داده‌های \lr{OOD} و فروپاشی معنایی منجر می‌شود؛ دوم، کوری معنایی، به این صورت که با تمرکز بر نوسانات آماری توکن‌ها در مواجهه با ماهیت انعطاف‌پذیر زبان (و پدیده‌ی پارافریز)، تغییرات ظاهری گرامر را به اشتباه به‌عنوان فقدان استدلال تفسیر می‌کنند؛ سوم، پیش‌نیازهای پنهان استعلام، چرا که این روش‌ها برای ارزیابی نوسانات، نیازمند استخراج توزیع احتمالات دقیق (لوجیت‌های لایه‌ی آخر) برای تک‌تک توکن‌های خروجی هستند که در سیستم‌های تجاری جعبه‌سیاه، مسدود و غیرقابل‌دسترس است؛ و چهارم، پیچیدگی ترکیبیاتی و تحمیل هزینه‌ی نمایی ($\mathcal{O}(2^N)$) برای استعلام از \lr{API}ها که در متون آزاد رشد غیرقابل‌تصوری دارد. در نهایت، همین نارسایی‌های توزیعی باعث شده تا معیارهای ارزیابی موجود در این دامنه نیز (بخش~\ref{sec:3-1-3-4}) دچار تله‌های منطقی و قانون گودهارت شوند.

برای خروج از این بن‌بست چندجانبه، واکاوی ادبیات ابزارهای نوین در بخش~\ref{sec:3-2} نشان می‌دهد که کلید این مسئله می‌تواند در گرو تغییر بنیادین در واحد تحلیل و دگرگونی در معماری ارزیاب یافت شود. ارتقای واحد معنایی از سطح نشانه‌های خُرد نحوی (توکن) به جمله (بخش~\ref{sec:3-2-1})، راهکاری اصولی و برآمده از علوم شناختی برای حفظ یکپارچگی گزاره‌ای است. با اتکا به جملات به‌عنوان مفاهیم پایه‌ی معنایی، می‌توان بدون تخریب ساختار گرامری، میزان اثرگذاری پیش‌آیندها و گزاره‌های کلان ورودی را بررسی کرد و از گرداب داده‌های خارج از توزیع و کوری معنایی رهایی یافت.

از سوی دیگر، همان‌گونه که در بررسی معماری‌های نوین (بخش‌های~\ref{sec:3-2-2} و~\ref{sec:3-2-3}) نشان داده شد، \lr{EBM}ها ظرفیت نظری بالایی برای ارزیابی کل‌نگر و ساختاریافته‌ی متون ارائه می‌دهند. این مدل‌ها نشان داده‌اند که می‌توانند با عبور از ارزیابی‌های تقلیل‌گرایانه‌ی توکن‌به‌توکن، انطباق‌پذیری سراسری یک توالی را داوری کنند. با این وجود، استقرار مستقیم این متدولوژی‌ها در پارادایم تفسیرپذیری جعبه‌سیاه با موانع ساختاری عمیقی روبه‌روست. رویکردهایی نظیر چارچوب \lr{COLD} برای هدایت تولید، نیازمند دسترسی به فضای پیوسته‌ی لوجیت‌ها و محاسبه‌ی متوالی گرادیان انرژی در زمان اجرا هستند. از طرفی، مدل‌هایی نظیر \lr{REBM} بار محاسباتی فاجعه‌باری را در فاز استنتاج برای نمونه‌برداری و رتبه‌بندی مجدد کاندیداها تحمیل می‌کنند. این وابستگی‌های مطلق به دسترسی جعبه‌سفید و توان محاسباتی بالا، با رسالت حسابرسی سریع و بیرونی \lr{API}های تجاری در تضاد کامل است.

برای غلبه بر این گلوگاه محاسباتی در زمان اجرا، ادبیات پیش‌بینی ساختاریافته به رویکرد استنتاج مستهلک‌شده (نظیر چارچوب \lr{ENGINE}) و آموزش مفسرهای جانشین (نظیر \lr{SOInter}) متمایل شده است. این پارادایم نشان می‌دهد که می‌توان با آموزش یک شبکه‌ی ثانویه، بار محاسباتی و جستجوی نمایی را از زمان استنتاج به فاز آموزش منتقل کرد. با این حال، همان‌طور که پیش‌تر نقد شد، این معماری‌ها غالباً برای فضاهای خروجی با ابعاد ثابت بهینه‌سازی شده‌اند و تعمیم آن‌ها به متون آزاد و نامحدود، آن هم به‌گونه‌ای که نیازی به استخراج گرادیان از شبکه‌ی اصلی نداشته باشد و منجر به تولید داده‌های \lr{OOD} نگردد، همچنان یک مسئله‌ی باز باقی مانده است.

مجموع این بررسی‌ها در این فصل نشان می‌دهد که در اکوسیستم فعلی پردازش زبان طبیعی، یک شکاف تحقیقاتی مشخص در ممیزی پساعملی مدل‌های زبانی بزرگ به چشم می‌خورد؛ این خلأ متوجه توسعه‌ی رویکردهایی است که بتوانند به‌طور هم‌زمان چندین الزام کلیدی را ارضا کنند: نخست، فعالیت در سطح گزاره‌ها و جملات به‌عنوان مفاهیم معنایی پایدار به‌منظور خنثی‌سازی خطای شوک توزیعی و افزونگی واژگانی؛ دوم، پایبندی به رژیم جعبه‌سیاه و استنتاج صرفاً بر مبنای متن خام خروجی بدون نیاز به گرادیان‌ها یا لوجیت‌های مدل هدف؛ سوم، بهره‌گیری از یک ارزیابی کل‌نگر و سراسری بر پایه‌ی مدل‌های انرژی برای سنجش انطباق‌پذیری متنی؛ و چهارم، انتقال بار محاسباتی به فاز آموزش به کمک مفسرهای جانشین به‌منظور دستیابی به زمان اجرای مستقل از تعداد جملات در فاز استنتاج.

بررسی ادبیات نشان می‌دهد که اگرچه هر یک از این مؤلفه‌ها به‌صورت مجزا در شاخه‌های متفاوتی از یادگیری ماشین و پردازش زبان طبیعی توسعه یافته‌اند، اما چگونگی تلفیق ساختاریافته‌ی آن‌ها در قالب یک چارچوب ممیزی منسجم و ارزیابی آن در کنار خط‌پایه‌های مقایسه‌ای معتبر (اعم از روش‌های جعبه‌سیاه و روش‌های جعبه‌سفید و مبتنی بر گرادیان با حفظ مزیت اطلاعات درونی آن‌ها)، همچنان یکی از افق‌های تحقیقاتی باز و مهم در این حوزه به شمار می‌رود.